% F083｜三种证明方法，按具体结构配图。
\clearpage
\begin{proofbox}
\textbf{\textcolor{CProof}{证明一：直接代入 $-x$，看指数如何互换}}
% Formula-based samples; see F083 checks/build_tikz_diagrams.py.
\begin{center}\begin{tikzpicture}[font=\small,>=stealth,line join=round]
\draw[gray!40,rounded corners=3pt] (0,0) rectangle (14.2,8.4);
\node at (7.1,7.75) {奇1：$y=a^x-a^{-x}$};
\begin{scope}[shift={(4.400000,3.900000)},x=1.700000cm,y=0.450000cm]
\draw[->,gray] (-2,0)--(2,0) node[right] {$x$};
\draw[->,gray] (0,-5)--(0,5) node[above] {$y$};
\draw[gray] (-1,-.035/0.450000)--(-1,.035/0.450000);
\node[anchor=north] at (-1,-.07/0.450000) {$-1$};
\draw[gray] (1,-.035/0.450000)--(1,.035/0.450000);
\node[anchor=north] at (1,-.07/0.450000) {$1$};
\begin{scope}
\clip (-2,-5) rectangle (2,5);
\draw[blue,thick,solid] plot coordinates {(-2.0000000,-3.7500000) (-1.9933333,-3.7304008) (-1.9866667,-3.7108813) (-1.9800000,-3.6914409) (-1.9733333,-3.6720795) (-1.9666667,-3.6527964) (-1.9600000,-3.6335913) (-1.9533333,-3.6144639) (-1.9466667,-3.5954136) (-1.9400000,-3.5764400) (-1.9333333,-3.5575429) (-1.9266667,-3.5387217) (-1.9200000,-3.5199761) (-1.9133333,-3.5013056) (-1.9066667,-3.4827099) (-1.9000000,-3.4641886) (-1.8933333,-3.4457412) (-1.8866667,-3.4273675) (-1.8800000,-3.4090669) (-1.8733333,-3.3908391) (-1.8666667,-3.3726837) (-1.8600000,-3.3546003) (-1.8533333,-3.3365886) (-1.8466667,-3.3186481) (-1.8400000,-3.3007785) (-1.8333333,-3.2829794) (-1.8266667,-3.2652503) (-1.8200000,-3.2475910) (-1.8133333,-3.2300011) (-1.8066667,-3.2124801) (-1.8000000,-3.1950277) (-1.7933333,-3.1776435) (-1.7866667,-3.1603272) (-1.7800000,-3.1430783) (-1.7733333,-3.1258966) (-1.7666667,-3.1087817) (-1.7600000,-3.0917331) (-1.7533333,-3.0747505) (-1.7466667,-3.0578336) (-1.7400000,-3.0409820) (-1.7333333,-3.0241953) (-1.7266667,-3.0074732) (-1.7200000,-2.9908153) (-1.7133333,-2.9742213) (-1.7066667,-2.9576908) (-1.7000000,-2.9412235) (-1.6933333,-2.9248189) (-1.6866667,-2.9084769) (-1.6800000,-2.8921969) (-1.6733333,-2.8759787) (-1.6666667,-2.8598218) (-1.6600000,-2.8437261) (-1.6533333,-2.8276911) (-1.6466667,-2.8117164) (-1.6400000,-2.7958018) (-1.6333333,-2.7799469) (-1.6266667,-2.7641514) (-1.6200000,-2.7484149) (-1.6133333,-2.7327371) (-1.6066667,-2.7171176) (-1.6000000,-2.7015562) (-1.5933333,-2.6860524) (-1.5866667,-2.6706060) (-1.5800000,-2.6552166) (-1.5733333,-2.6398839) (-1.5666667,-2.6246076) (-1.5600000,-2.6093874) (-1.5533333,-2.5942228) (-1.5466667,-2.5791137) (-1.5400000,-2.5640596) (-1.5333333,-2.5490603) (-1.5266667,-2.5341154) (-1.5200000,-2.5192246) (-1.5133333,-2.5043876) (-1.5066667,-2.4896041) (-1.5000000,-2.4748737) (-1.4933333,-2.4601962) (-1.4866667,-2.4455713) (-1.4800000,-2.4309985) (-1.4733333,-2.4164777) (-1.4666667,-2.4020085) (-1.4600000,-2.3875905) (-1.4533333,-2.3732235) (-1.4466667,-2.3589073) (-1.4400000,-2.3446414) (-1.4333333,-2.3304255) (-1.4266667,-2.3162594) (-1.4200000,-2.3021428) (-1.4133333,-2.2880753) (-1.4066667,-2.2740567) (-1.4000000,-2.2600867) (-1.3933333,-2.2461649) (-1.3866667,-2.2322911) (-1.3800000,-2.2184649) (-1.3733333,-2.2046861) (-1.3666667,-2.1909544) (-1.3600000,-2.1772695) (-1.3533333,-2.1636311) (-1.3466667,-2.1500388) (-1.3400000,-2.1364925) (-1.3333333,-2.1229918) (-1.3266667,-2.1095365) (-1.3200000,-2.0961262) (-1.3133333,-2.0827606) (-1.3066667,-2.0694395) (-1.3000000,-2.0561626) (-1.2933333,-2.0429296) (-1.2866667,-2.0297403) (-1.2800000,-2.0165943) (-1.2733333,-2.0034913) (-1.2666667,-1.9904311) (-1.2600000,-1.9774134) (-1.2533333,-1.9644380) (-1.2466667,-1.9515045) (-1.2400000,-1.9386127) (-1.2333333,-1.9257622) (-1.2266667,-1.9129529) (-1.2200000,-1.9001845) (-1.2133333,-1.8874566) (-1.2066667,-1.8747690) (-1.2000000,-1.8621214) (-1.1933333,-1.8495136) (-1.1866667,-1.8369453) (-1.1800000,-1.8244163) (-1.1733333,-1.8119262) (-1.1666667,-1.7994747) (-1.1600000,-1.7870617) (-1.1533333,-1.7746869) (-1.1466667,-1.7623500) (-1.1400000,-1.7500507) (-1.1333333,-1.7377887) (-1.1266667,-1.7255639) (-1.1200000,-1.7133759) (-1.1133333,-1.7012245) (-1.1066667,-1.6891094) (-1.1000000,-1.6770304) (-1.0933333,-1.6649872) (-1.0866667,-1.6529796) (-1.0800000,-1.6410073) (-1.0733333,-1.6290700) (-1.0666667,-1.6171674) (-1.0600000,-1.6052995) (-1.0533333,-1.5934658) (-1.0466667,-1.5816661) (-1.0400000,-1.5699002) (-1.0333333,-1.5581678) (-1.0266667,-1.5464687) (-1.0200000,-1.5348026) (-1.0133333,-1.5231693) (-1.0066667,-1.5115685) (-1.0000000,-1.5000000) (-0.9933333,-1.4884635) (-0.9866667,-1.4769588) (-0.9800000,-1.4654857) (-0.9733333,-1.4540438) (-0.9666667,-1.4426330) (-0.9600000,-1.4312530) (-0.9533333,-1.4199035) (-0.9466667,-1.4085844) (-0.9400000,-1.3972954) (-0.9333333,-1.3860361) (-0.9266667,-1.3748065) (-0.9200000,-1.3636063) (-0.9133333,-1.3524351) (-0.9066667,-1.3412929) (-0.9000000,-1.3301793) (-0.8933333,-1.3190940) (-0.8866667,-1.3080370) (-0.8800000,-1.2970079) (-0.8733333,-1.2860064) (-0.8666667,-1.2750325) (-0.8600000,-1.2640858) (-0.8533333,-1.2531660) (-0.8466667,-1.2422730) (-0.8400000,-1.2314066) (-0.8333333,-1.2205664) (-0.8266667,-1.2097523) (-0.8200000,-1.1989640) (-0.8133333,-1.1882014) (-0.8066667,-1.1774641) (-0.8000000,-1.1667519) (-0.7933333,-1.1560647) (-0.7866667,-1.1454022) (-0.7800000,-1.1347641) (-0.7733333,-1.1241502) (-0.7666667,-1.1135604) (-0.7600000,-1.1029943) (-0.7533333,-1.0924518) (-0.7466667,-1.0819326) (-0.7400000,-1.0714365) (-0.7333333,-1.0609633) (-0.7266667,-1.0505127) (-0.7200000,-1.0400846) (-0.7133333,-1.0296787) (-0.7066667,-1.0192947) (-0.7000000,-1.0089326) (-0.6933333,-0.9985920) (-0.6866667,-0.9882727) (-0.6800000,-0.9779745) (-0.6733333,-0.9676972) (-0.6666667,-0.9574405) (-0.6600000,-0.9472043) (-0.6533333,-0.9369884) (-0.6466667,-0.9267924) (-0.6400000,-0.9166162) (-0.6333333,-0.9064596) (-0.6266667,-0.8963224) (-0.6200000,-0.8862043) (-0.6133333,-0.8761051) (-0.6066667,-0.8660246) (-0.6000000,-0.8559626) (-0.5933333,-0.8459189) (-0.5866667,-0.8358933) (-0.5800000,-0.8258855) (-0.5733333,-0.8158953) (-0.5666667,-0.8059226) (-0.5600000,-0.7959671) (-0.5533333,-0.7860285) (-0.5466667,-0.7761068) (-0.5400000,-0.7662016) (-0.5333333,-0.7563128) (-0.5266667,-0.7464401) (-0.5200000,-0.7365834) (-0.5133333,-0.7267424) (-0.5066667,-0.7169169) (-0.5000000,-0.7071068) (-0.4933333,-0.6973117) (-0.4866667,-0.6875315) (-0.4800000,-0.6777660) (-0.4733333,-0.6680150) (-0.4666667,-0.6582783) (-0.4600000,-0.6485556) (-0.4533333,-0.6388467) (-0.4466667,-0.6291515) (-0.4400000,-0.6194697) (-0.4333333,-0.6098012) (-0.4266667,-0.6001456) (-0.4200000,-0.5905029) (-0.4133333,-0.5808728) (-0.4066667,-0.5712551) (-0.4000000,-0.5616496) (-0.3933333,-0.5520561) (-0.3866667,-0.5424744) (-0.3800000,-0.5329043) (-0.3733333,-0.5233455) (-0.3666667,-0.5137979) (-0.3600000,-0.5042613) (-0.3533333,-0.4947355) (-0.3466667,-0.4852202) (-0.3400000,-0.4757153) (-0.3333333,-0.4662205) (-0.3266667,-0.4567357) (-0.3200000,-0.4472607) (-0.3133333,-0.4377952) (-0.3066667,-0.4283390) (-0.3000000,-0.4188920) (-0.2933333,-0.4094540) (-0.2866667,-0.4000246) (-0.2800000,-0.3906039) (-0.2733333,-0.3811914) (-0.2666667,-0.3717871) (-0.2600000,-0.3623908) (-0.2533333,-0.3530022) (-0.2466667,-0.3436211) (-0.2400000,-0.3342473) (-0.2333333,-0.3248807) (-0.2266667,-0.3155211) (-0.2200000,-0.3061682) (-0.2133333,-0.2968218) (-0.2066667,-0.2874817) (-0.2000000,-0.2781478) (-0.1933333,-0.2688198) (-0.1866667,-0.2594976) (-0.1800000,-0.2501809) (-0.1733333,-0.2408695) (-0.1666667,-0.2315633) (-0.1600000,-0.2222621) (-0.1533333,-0.2129656) (-0.1466667,-0.2036736) (-0.1400000,-0.1943860) (-0.1333333,-0.1851025) (-0.1266667,-0.1758230) (-0.1200000,-0.1665472) (-0.1133333,-0.1572750) (-0.1066667,-0.1480062) (-0.1000000,-0.1387405) (-0.0933333,-0.1294777) (-0.0866667,-0.1202178) (-0.0800000,-0.1109604) (-0.0733333,-0.1017054) (-0.0666667,-0.0924525) (-0.0600000,-0.0832016) (-0.0533333,-0.0739525) (-0.0466667,-0.0647050) (-0.0400000,-0.0554589) (-0.0333333,-0.0462139) (-0.0266667,-0.0369700) (-0.0200000,-0.0277268) (-0.0133333,-0.0184842) (-0.0066667,-0.0092420) (0.0000000,0.0000000) (0.0066667,0.0092420) (0.0133333,0.0184842) (0.0200000,0.0277268) (0.0266667,0.0369700) (0.0333333,0.0462139) (0.0400000,0.0554589) (0.0466667,0.0647050) (0.0533333,0.0739525) (0.0600000,0.0832016) (0.0666667,0.0924525) (0.0733333,0.1017054) (0.0800000,0.1109604) (0.0866667,0.1202178) (0.0933333,0.1294777) (0.1000000,0.1387405) (0.1066667,0.1480062) (0.1133333,0.1572750) (0.1200000,0.1665472) (0.1266667,0.1758230) (0.1333333,0.1851025) (0.1400000,0.1943860) (0.1466667,0.2036736) (0.1533333,0.2129656) (0.1600000,0.2222621) (0.1666667,0.2315633) (0.1733333,0.2408695) (0.1800000,0.2501809) (0.1866667,0.2594976) (0.1933333,0.2688198) (0.2000000,0.2781478) (0.2066667,0.2874817) (0.2133333,0.2968218) (0.2200000,0.3061682) (0.2266667,0.3155211) (0.2333333,0.3248807) (0.2400000,0.3342473) (0.2466667,0.3436211) (0.2533333,0.3530022) (0.2600000,0.3623908) (0.2666667,0.3717871) (0.2733333,0.3811914) (0.2800000,0.3906039) (0.2866667,0.4000246) (0.2933333,0.4094540) (0.3000000,0.4188920) (0.3066667,0.4283390) (0.3133333,0.4377952) (0.3200000,0.4472607) (0.3266667,0.4567357) (0.3333333,0.4662205) (0.3400000,0.4757153) (0.3466667,0.4852202) (0.3533333,0.4947355) (0.3600000,0.5042613) (0.3666667,0.5137979) (0.3733333,0.5233455) (0.3800000,0.5329043) (0.3866667,0.5424744) (0.3933333,0.5520561) (0.4000000,0.5616496) (0.4066667,0.5712551) (0.4133333,0.5808728) (0.4200000,0.5905029) (0.4266667,0.6001456) (0.4333333,0.6098012) (0.4400000,0.6194697) (0.4466667,0.6291515) (0.4533333,0.6388467) (0.4600000,0.6485556) (0.4666667,0.6582783) (0.4733333,0.6680150) (0.4800000,0.6777660) (0.4866667,0.6875315) (0.4933333,0.6973117) (0.5000000,0.7071068) (0.5066667,0.7169169) (0.5133333,0.7267424) (0.5200000,0.7365834) (0.5266667,0.7464401) (0.5333333,0.7563128) (0.5400000,0.7662016) (0.5466667,0.7761068) (0.5533333,0.7860285) (0.5600000,0.7959671) (0.5666667,0.8059226) (0.5733333,0.8158953) (0.5800000,0.8258855) (0.5866667,0.8358933) (0.5933333,0.8459189) (0.6000000,0.8559626) (0.6066667,0.8660246) (0.6133333,0.8761051) (0.6200000,0.8862043) (0.6266667,0.8963224) (0.6333333,0.9064596) (0.6400000,0.9166162) (0.6466667,0.9267924) (0.6533333,0.9369884) (0.6600000,0.9472043) (0.6666667,0.9574405) (0.6733333,0.9676972) (0.6800000,0.9779745) (0.6866667,0.9882727) (0.6933333,0.9985920) (0.7000000,1.0089326) (0.7066667,1.0192947) (0.7133333,1.0296787) (0.7200000,1.0400846) (0.7266667,1.0505127) (0.7333333,1.0609633) (0.7400000,1.0714365) (0.7466667,1.0819326) (0.7533333,1.0924518) (0.7600000,1.1029943) (0.7666667,1.1135604) (0.7733333,1.1241502) (0.7800000,1.1347641) (0.7866667,1.1454022) (0.7933333,1.1560647) (0.8000000,1.1667519) (0.8066667,1.1774641) (0.8133333,1.1882014) (0.8200000,1.1989640) (0.8266667,1.2097523) (0.8333333,1.2205664) (0.8400000,1.2314066) (0.8466667,1.2422730) (0.8533333,1.2531660) (0.8600000,1.2640858) (0.8666667,1.2750325) (0.8733333,1.2860064) (0.8800000,1.2970079) (0.8866667,1.3080370) (0.8933333,1.3190940) (0.9000000,1.3301793) (0.9066667,1.3412929) (0.9133333,1.3524351) (0.9200000,1.3636063) (0.9266667,1.3748065) (0.9333333,1.3860361) (0.9400000,1.3972954) (0.9466667,1.4085844) (0.9533333,1.4199035) (0.9600000,1.4312530) (0.9666667,1.4426330) (0.9733333,1.4540438) (0.9800000,1.4654857) (0.9866667,1.4769588) (0.9933333,1.4884635) (1.0000000,1.5000000) (1.0066667,1.5115685) (1.0133333,1.5231693) (1.0200000,1.5348026) (1.0266667,1.5464687) (1.0333333,1.5581678) (1.0400000,1.5699002) (1.0466667,1.5816661) (1.0533333,1.5934658) (1.0600000,1.6052995) (1.0666667,1.6171674) (1.0733333,1.6290700) (1.0800000,1.6410073) (1.0866667,1.6529796) (1.0933333,1.6649872) (1.1000000,1.6770304) (1.1066667,1.6891094) (1.1133333,1.7012245) (1.1200000,1.7133759) (1.1266667,1.7255639) (1.1333333,1.7377887) (1.1400000,1.7500507) (1.1466667,1.7623500) (1.1533333,1.7746869) (1.1600000,1.7870617) (1.1666667,1.7994747) (1.1733333,1.8119262) (1.1800000,1.8244163) (1.1866667,1.8369453) (1.1933333,1.8495136) (1.2000000,1.8621214) (1.2066667,1.8747690) (1.2133333,1.8874566) (1.2200000,1.9001845) (1.2266667,1.9129529) (1.2333333,1.9257622) (1.2400000,1.9386127) (1.2466667,1.9515045) (1.2533333,1.9644380) (1.2600000,1.9774134) (1.2666667,1.9904311) (1.2733333,2.0034913) (1.2800000,2.0165943) (1.2866667,2.0297403) (1.2933333,2.0429296) (1.3000000,2.0561626) (1.3066667,2.0694395) (1.3133333,2.0827606) (1.3200000,2.0961262) (1.3266667,2.1095365) (1.3333333,2.1229918) (1.3400000,2.1364925) (1.3466667,2.1500388) (1.3533333,2.1636311) (1.3600000,2.1772695) (1.3666667,2.1909544) (1.3733333,2.2046861) (1.3800000,2.2184649) (1.3866667,2.2322911) (1.3933333,2.2461649) (1.4000000,2.2600867) (1.4066667,2.2740567) (1.4133333,2.2880753) (1.4200000,2.3021428) (1.4266667,2.3162594) (1.4333333,2.3304255) (1.4400000,2.3446414) (1.4466667,2.3589073) (1.4533333,2.3732235) (1.4600000,2.3875905) (1.4666667,2.4020085) (1.4733333,2.4164777) (1.4800000,2.4309985) (1.4866667,2.4455713) (1.4933333,2.4601962) (1.5000000,2.4748737) (1.5066667,2.4896041) (1.5133333,2.5043876) (1.5200000,2.5192246) (1.5266667,2.5341154) (1.5333333,2.5490603) (1.5400000,2.5640596) (1.5466667,2.5791137) (1.5533333,2.5942228) (1.5600000,2.6093874) (1.5666667,2.6246076) (1.5733333,2.6398839) (1.5800000,2.6552166) (1.5866667,2.6706060) (1.5933333,2.6860524) (1.6000000,2.7015562) (1.6066667,2.7171176) (1.6133333,2.7327371) (1.6200000,2.7484149) (1.6266667,2.7641514) (1.6333333,2.7799469) (1.6400000,2.7958018) (1.6466667,2.8117164) (1.6533333,2.8276911) (1.6600000,2.8437261) (1.6666667,2.8598218) (1.6733333,2.8759787) (1.6800000,2.8921969) (1.6866667,2.9084769) (1.6933333,2.9248189) (1.7000000,2.9412235) (1.7066667,2.9576908) (1.7133333,2.9742213) (1.7200000,2.9908153) (1.7266667,3.0074732) (1.7333333,3.0241953) (1.7400000,3.0409820) (1.7466667,3.0578336) (1.7533333,3.0747505) (1.7600000,3.0917331) (1.7666667,3.1087817) (1.7733333,3.1258966) (1.7800000,3.1430783) (1.7866667,3.1603272) (1.7933333,3.1776435) (1.8000000,3.1950277) (1.8066667,3.2124801) (1.8133333,3.2300011) (1.8200000,3.2475910) (1.8266667,3.2652503) (1.8333333,3.2829794) (1.8400000,3.3007785) (1.8466667,3.3186481) (1.8533333,3.3365886) (1.8600000,3.3546003) (1.8666667,3.3726837) (1.8733333,3.3908391) (1.8800000,3.4090669) (1.8866667,3.4273675) (1.8933333,3.4457412) (1.9000000,3.4641886) (1.9066667,3.4827099) (1.9133333,3.5013056) (1.9200000,3.5199761) (1.9266667,3.5387217) (1.9333333,3.5575429) (1.9400000,3.5764400) (1.9466667,3.5954136) (1.9533333,3.6144639) (1.9600000,3.6335913) (1.9666667,3.6527964) (1.9733333,3.6720795) (1.9800000,3.6914409) (1.9866667,3.7108813) (1.9933333,3.7304008) (2.0000000,3.7500000)};
\draw[orange!85!black,thick,dashed] plot coordinates {(-2.0000000,3.7500000) (-1.9933333,3.7304008) (-1.9866667,3.7108813) (-1.9800000,3.6914409) (-1.9733333,3.6720795) (-1.9666667,3.6527964) (-1.9600000,3.6335913) (-1.9533333,3.6144639) (-1.9466667,3.5954136) (-1.9400000,3.5764400) (-1.9333333,3.5575429) (-1.9266667,3.5387217) (-1.9200000,3.5199761) (-1.9133333,3.5013056) (-1.9066667,3.4827099) (-1.9000000,3.4641886) (-1.8933333,3.4457412) (-1.8866667,3.4273675) (-1.8800000,3.4090669) (-1.8733333,3.3908391) (-1.8666667,3.3726837) (-1.8600000,3.3546003) (-1.8533333,3.3365886) (-1.8466667,3.3186481) (-1.8400000,3.3007785) (-1.8333333,3.2829794) (-1.8266667,3.2652503) (-1.8200000,3.2475910) (-1.8133333,3.2300011) (-1.8066667,3.2124801) (-1.8000000,3.1950277) (-1.7933333,3.1776435) (-1.7866667,3.1603272) (-1.7800000,3.1430783) (-1.7733333,3.1258966) (-1.7666667,3.1087817) (-1.7600000,3.0917331) (-1.7533333,3.0747505) (-1.7466667,3.0578336) (-1.7400000,3.0409820) (-1.7333333,3.0241953) (-1.7266667,3.0074732) (-1.7200000,2.9908153) (-1.7133333,2.9742213) (-1.7066667,2.9576908) (-1.7000000,2.9412235) (-1.6933333,2.9248189) (-1.6866667,2.9084769) (-1.6800000,2.8921969) (-1.6733333,2.8759787) (-1.6666667,2.8598218) (-1.6600000,2.8437261) (-1.6533333,2.8276911) (-1.6466667,2.8117164) (-1.6400000,2.7958018) (-1.6333333,2.7799469) (-1.6266667,2.7641514) (-1.6200000,2.7484149) (-1.6133333,2.7327371) (-1.6066667,2.7171176) (-1.6000000,2.7015562) (-1.5933333,2.6860524) (-1.5866667,2.6706060) (-1.5800000,2.6552166) (-1.5733333,2.6398839) (-1.5666667,2.6246076) (-1.5600000,2.6093874) (-1.5533333,2.5942228) (-1.5466667,2.5791137) (-1.5400000,2.5640596) (-1.5333333,2.5490603) (-1.5266667,2.5341154) (-1.5200000,2.5192246) (-1.5133333,2.5043876) (-1.5066667,2.4896041) (-1.5000000,2.4748737) (-1.4933333,2.4601962) (-1.4866667,2.4455713) (-1.4800000,2.4309985) (-1.4733333,2.4164777) (-1.4666667,2.4020085) (-1.4600000,2.3875905) (-1.4533333,2.3732235) (-1.4466667,2.3589073) (-1.4400000,2.3446414) (-1.4333333,2.3304255) (-1.4266667,2.3162594) (-1.4200000,2.3021428) (-1.4133333,2.2880753) (-1.4066667,2.2740567) (-1.4000000,2.2600867) (-1.3933333,2.2461649) (-1.3866667,2.2322911) (-1.3800000,2.2184649) (-1.3733333,2.2046861) (-1.3666667,2.1909544) (-1.3600000,2.1772695) (-1.3533333,2.1636311) (-1.3466667,2.1500388) (-1.3400000,2.1364925) (-1.3333333,2.1229918) (-1.3266667,2.1095365) (-1.3200000,2.0961262) (-1.3133333,2.0827606) (-1.3066667,2.0694395) (-1.3000000,2.0561626) (-1.2933333,2.0429296) (-1.2866667,2.0297403) (-1.2800000,2.0165943) (-1.2733333,2.0034913) (-1.2666667,1.9904311) (-1.2600000,1.9774134) (-1.2533333,1.9644380) (-1.2466667,1.9515045) (-1.2400000,1.9386127) (-1.2333333,1.9257622) (-1.2266667,1.9129529) (-1.2200000,1.9001845) (-1.2133333,1.8874566) (-1.2066667,1.8747690) (-1.2000000,1.8621214) (-1.1933333,1.8495136) (-1.1866667,1.8369453) (-1.1800000,1.8244163) (-1.1733333,1.8119262) (-1.1666667,1.7994747) (-1.1600000,1.7870617) (-1.1533333,1.7746869) (-1.1466667,1.7623500) (-1.1400000,1.7500507) (-1.1333333,1.7377887) (-1.1266667,1.7255639) (-1.1200000,1.7133759) (-1.1133333,1.7012245) (-1.1066667,1.6891094) (-1.1000000,1.6770304) (-1.0933333,1.6649872) (-1.0866667,1.6529796) (-1.0800000,1.6410073) (-1.0733333,1.6290700) (-1.0666667,1.6171674) (-1.0600000,1.6052995) (-1.0533333,1.5934658) (-1.0466667,1.5816661) (-1.0400000,1.5699002) (-1.0333333,1.5581678) (-1.0266667,1.5464687) (-1.0200000,1.5348026) (-1.0133333,1.5231693) (-1.0066667,1.5115685) (-1.0000000,1.5000000) (-0.9933333,1.4884635) (-0.9866667,1.4769588) (-0.9800000,1.4654857) (-0.9733333,1.4540438) (-0.9666667,1.4426330) (-0.9600000,1.4312530) (-0.9533333,1.4199035) (-0.9466667,1.4085844) (-0.9400000,1.3972954) (-0.9333333,1.3860361) (-0.9266667,1.3748065) (-0.9200000,1.3636063) (-0.9133333,1.3524351) (-0.9066667,1.3412929) (-0.9000000,1.3301793) (-0.8933333,1.3190940) (-0.8866667,1.3080370) (-0.8800000,1.2970079) (-0.8733333,1.2860064) (-0.8666667,1.2750325) (-0.8600000,1.2640858) (-0.8533333,1.2531660) (-0.8466667,1.2422730) (-0.8400000,1.2314066) (-0.8333333,1.2205664) (-0.8266667,1.2097523) (-0.8200000,1.1989640) (-0.8133333,1.1882014) (-0.8066667,1.1774641) (-0.8000000,1.1667519) (-0.7933333,1.1560647) (-0.7866667,1.1454022) (-0.7800000,1.1347641) (-0.7733333,1.1241502) (-0.7666667,1.1135604) (-0.7600000,1.1029943) (-0.7533333,1.0924518) (-0.7466667,1.0819326) (-0.7400000,1.0714365) (-0.7333333,1.0609633) (-0.7266667,1.0505127) (-0.7200000,1.0400846) (-0.7133333,1.0296787) (-0.7066667,1.0192947) (-0.7000000,1.0089326) (-0.6933333,0.9985920) (-0.6866667,0.9882727) (-0.6800000,0.9779745) (-0.6733333,0.9676972) (-0.6666667,0.9574405) (-0.6600000,0.9472043) (-0.6533333,0.9369884) (-0.6466667,0.9267924) (-0.6400000,0.9166162) (-0.6333333,0.9064596) (-0.6266667,0.8963224) (-0.6200000,0.8862043) (-0.6133333,0.8761051) (-0.6066667,0.8660246) (-0.6000000,0.8559626) (-0.5933333,0.8459189) (-0.5866667,0.8358933) (-0.5800000,0.8258855) (-0.5733333,0.8158953) (-0.5666667,0.8059226) (-0.5600000,0.7959671) (-0.5533333,0.7860285) (-0.5466667,0.7761068) (-0.5400000,0.7662016) (-0.5333333,0.7563128) (-0.5266667,0.7464401) (-0.5200000,0.7365834) (-0.5133333,0.7267424) (-0.5066667,0.7169169) (-0.5000000,0.7071068) (-0.4933333,0.6973117) (-0.4866667,0.6875315) (-0.4800000,0.6777660) (-0.4733333,0.6680150) (-0.4666667,0.6582783) (-0.4600000,0.6485556) (-0.4533333,0.6388467) (-0.4466667,0.6291515) (-0.4400000,0.6194697) (-0.4333333,0.6098012) (-0.4266667,0.6001456) (-0.4200000,0.5905029) (-0.4133333,0.5808728) (-0.4066667,0.5712551) (-0.4000000,0.5616496) (-0.3933333,0.5520561) (-0.3866667,0.5424744) (-0.3800000,0.5329043) (-0.3733333,0.5233455) (-0.3666667,0.5137979) (-0.3600000,0.5042613) (-0.3533333,0.4947355) (-0.3466667,0.4852202) (-0.3400000,0.4757153) (-0.3333333,0.4662205) (-0.3266667,0.4567357) (-0.3200000,0.4472607) (-0.3133333,0.4377952) (-0.3066667,0.4283390) (-0.3000000,0.4188920) (-0.2933333,0.4094540) (-0.2866667,0.4000246) (-0.2800000,0.3906039) (-0.2733333,0.3811914) (-0.2666667,0.3717871) (-0.2600000,0.3623908) (-0.2533333,0.3530022) (-0.2466667,0.3436211) (-0.2400000,0.3342473) (-0.2333333,0.3248807) (-0.2266667,0.3155211) (-0.2200000,0.3061682) (-0.2133333,0.2968218) (-0.2066667,0.2874817) (-0.2000000,0.2781478) (-0.1933333,0.2688198) (-0.1866667,0.2594976) (-0.1800000,0.2501809) (-0.1733333,0.2408695) (-0.1666667,0.2315633) (-0.1600000,0.2222621) (-0.1533333,0.2129656) (-0.1466667,0.2036736) (-0.1400000,0.1943860) (-0.1333333,0.1851025) (-0.1266667,0.1758230) (-0.1200000,0.1665472) (-0.1133333,0.1572750) (-0.1066667,0.1480062) (-0.1000000,0.1387405) (-0.0933333,0.1294777) (-0.0866667,0.1202178) (-0.0800000,0.1109604) (-0.0733333,0.1017054) (-0.0666667,0.0924525) (-0.0600000,0.0832016) (-0.0533333,0.0739525) (-0.0466667,0.0647050) (-0.0400000,0.0554589) (-0.0333333,0.0462139) (-0.0266667,0.0369700) (-0.0200000,0.0277268) (-0.0133333,0.0184842) (-0.0066667,0.0092420) (0.0000000,0.0000000) (0.0066667,-0.0092420) (0.0133333,-0.0184842) (0.0200000,-0.0277268) (0.0266667,-0.0369700) (0.0333333,-0.0462139) (0.0400000,-0.0554589) (0.0466667,-0.0647050) (0.0533333,-0.0739525) (0.0600000,-0.0832016) (0.0666667,-0.0924525) (0.0733333,-0.1017054) (0.0800000,-0.1109604) (0.0866667,-0.1202178) (0.0933333,-0.1294777) (0.1000000,-0.1387405) (0.1066667,-0.1480062) (0.1133333,-0.1572750) (0.1200000,-0.1665472) (0.1266667,-0.1758230) (0.1333333,-0.1851025) (0.1400000,-0.1943860) (0.1466667,-0.2036736) (0.1533333,-0.2129656) (0.1600000,-0.2222621) (0.1666667,-0.2315633) (0.1733333,-0.2408695) (0.1800000,-0.2501809) (0.1866667,-0.2594976) (0.1933333,-0.2688198) (0.2000000,-0.2781478) (0.2066667,-0.2874817) (0.2133333,-0.2968218) (0.2200000,-0.3061682) (0.2266667,-0.3155211) (0.2333333,-0.3248807) (0.2400000,-0.3342473) (0.2466667,-0.3436211) (0.2533333,-0.3530022) (0.2600000,-0.3623908) (0.2666667,-0.3717871) (0.2733333,-0.3811914) (0.2800000,-0.3906039) (0.2866667,-0.4000246) (0.2933333,-0.4094540) (0.3000000,-0.4188920) (0.3066667,-0.4283390) (0.3133333,-0.4377952) (0.3200000,-0.4472607) (0.3266667,-0.4567357) (0.3333333,-0.4662205) (0.3400000,-0.4757153) (0.3466667,-0.4852202) (0.3533333,-0.4947355) (0.3600000,-0.5042613) (0.3666667,-0.5137979) (0.3733333,-0.5233455) (0.3800000,-0.5329043) (0.3866667,-0.5424744) (0.3933333,-0.5520561) (0.4000000,-0.5616496) (0.4066667,-0.5712551) (0.4133333,-0.5808728) (0.4200000,-0.5905029) (0.4266667,-0.6001456) (0.4333333,-0.6098012) (0.4400000,-0.6194697) (0.4466667,-0.6291515) (0.4533333,-0.6388467) (0.4600000,-0.6485556) (0.4666667,-0.6582783) (0.4733333,-0.6680150) (0.4800000,-0.6777660) (0.4866667,-0.6875315) (0.4933333,-0.6973117) (0.5000000,-0.7071068) (0.5066667,-0.7169169) (0.5133333,-0.7267424) (0.5200000,-0.7365834) (0.5266667,-0.7464401) (0.5333333,-0.7563128) (0.5400000,-0.7662016) (0.5466667,-0.7761068) (0.5533333,-0.7860285) (0.5600000,-0.7959671) (0.5666667,-0.8059226) (0.5733333,-0.8158953) (0.5800000,-0.8258855) (0.5866667,-0.8358933) (0.5933333,-0.8459189) (0.6000000,-0.8559626) (0.6066667,-0.8660246) (0.6133333,-0.8761051) (0.6200000,-0.8862043) (0.6266667,-0.8963224) (0.6333333,-0.9064596) (0.6400000,-0.9166162) (0.6466667,-0.9267924) (0.6533333,-0.9369884) (0.6600000,-0.9472043) (0.6666667,-0.9574405) (0.6733333,-0.9676972) (0.6800000,-0.9779745) (0.6866667,-0.9882727) (0.6933333,-0.9985920) (0.7000000,-1.0089326) (0.7066667,-1.0192947) (0.7133333,-1.0296787) (0.7200000,-1.0400846) (0.7266667,-1.0505127) (0.7333333,-1.0609633) (0.7400000,-1.0714365) (0.7466667,-1.0819326) (0.7533333,-1.0924518) (0.7600000,-1.1029943) (0.7666667,-1.1135604) (0.7733333,-1.1241502) (0.7800000,-1.1347641) (0.7866667,-1.1454022) (0.7933333,-1.1560647) (0.8000000,-1.1667519) (0.8066667,-1.1774641) (0.8133333,-1.1882014) (0.8200000,-1.1989640) (0.8266667,-1.2097523) (0.8333333,-1.2205664) (0.8400000,-1.2314066) (0.8466667,-1.2422730) (0.8533333,-1.2531660) (0.8600000,-1.2640858) (0.8666667,-1.2750325) (0.8733333,-1.2860064) (0.8800000,-1.2970079) (0.8866667,-1.3080370) (0.8933333,-1.3190940) (0.9000000,-1.3301793) (0.9066667,-1.3412929) (0.9133333,-1.3524351) (0.9200000,-1.3636063) (0.9266667,-1.3748065) (0.9333333,-1.3860361) (0.9400000,-1.3972954) (0.9466667,-1.4085844) (0.9533333,-1.4199035) (0.9600000,-1.4312530) (0.9666667,-1.4426330) (0.9733333,-1.4540438) (0.9800000,-1.4654857) (0.9866667,-1.4769588) (0.9933333,-1.4884635) (1.0000000,-1.5000000) (1.0066667,-1.5115685) (1.0133333,-1.5231693) (1.0200000,-1.5348026) (1.0266667,-1.5464687) (1.0333333,-1.5581678) (1.0400000,-1.5699002) (1.0466667,-1.5816661) (1.0533333,-1.5934658) (1.0600000,-1.6052995) (1.0666667,-1.6171674) (1.0733333,-1.6290700) (1.0800000,-1.6410073) (1.0866667,-1.6529796) (1.0933333,-1.6649872) (1.1000000,-1.6770304) (1.1066667,-1.6891094) (1.1133333,-1.7012245) (1.1200000,-1.7133759) (1.1266667,-1.7255639) (1.1333333,-1.7377887) (1.1400000,-1.7500507) (1.1466667,-1.7623500) (1.1533333,-1.7746869) (1.1600000,-1.7870617) (1.1666667,-1.7994747) (1.1733333,-1.8119262) (1.1800000,-1.8244163) (1.1866667,-1.8369453) (1.1933333,-1.8495136) (1.2000000,-1.8621214) (1.2066667,-1.8747690) (1.2133333,-1.8874566) (1.2200000,-1.9001845) (1.2266667,-1.9129529) (1.2333333,-1.9257622) (1.2400000,-1.9386127) (1.2466667,-1.9515045) (1.2533333,-1.9644380) (1.2600000,-1.9774134) (1.2666667,-1.9904311) (1.2733333,-2.0034913) (1.2800000,-2.0165943) (1.2866667,-2.0297403) (1.2933333,-2.0429296) (1.3000000,-2.0561626) (1.3066667,-2.0694395) (1.3133333,-2.0827606) (1.3200000,-2.0961262) (1.3266667,-2.1095365) (1.3333333,-2.1229918) (1.3400000,-2.1364925) (1.3466667,-2.1500388) (1.3533333,-2.1636311) (1.3600000,-2.1772695) (1.3666667,-2.1909544) (1.3733333,-2.2046861) (1.3800000,-2.2184649) (1.3866667,-2.2322911) (1.3933333,-2.2461649) (1.4000000,-2.2600867) (1.4066667,-2.2740567) (1.4133333,-2.2880753) (1.4200000,-2.3021428) (1.4266667,-2.3162594) (1.4333333,-2.3304255) (1.4400000,-2.3446414) (1.4466667,-2.3589073) (1.4533333,-2.3732235) (1.4600000,-2.3875905) (1.4666667,-2.4020085) (1.4733333,-2.4164777) (1.4800000,-2.4309985) (1.4866667,-2.4455713) (1.4933333,-2.4601962) (1.5000000,-2.4748737) (1.5066667,-2.4896041) (1.5133333,-2.5043876) (1.5200000,-2.5192246) (1.5266667,-2.5341154) (1.5333333,-2.5490603) (1.5400000,-2.5640596) (1.5466667,-2.5791137) (1.5533333,-2.5942228) (1.5600000,-2.6093874) (1.5666667,-2.6246076) (1.5733333,-2.6398839) (1.5800000,-2.6552166) (1.5866667,-2.6706060) (1.5933333,-2.6860524) (1.6000000,-2.7015562) (1.6066667,-2.7171176) (1.6133333,-2.7327371) (1.6200000,-2.7484149) (1.6266667,-2.7641514) (1.6333333,-2.7799469) (1.6400000,-2.7958018) (1.6466667,-2.8117164) (1.6533333,-2.8276911) (1.6600000,-2.8437261) (1.6666667,-2.8598218) (1.6733333,-2.8759787) (1.6800000,-2.8921969) (1.6866667,-2.9084769) (1.6933333,-2.9248189) (1.7000000,-2.9412235) (1.7066667,-2.9576908) (1.7133333,-2.9742213) (1.7200000,-2.9908153) (1.7266667,-3.0074732) (1.7333333,-3.0241953) (1.7400000,-3.0409820) (1.7466667,-3.0578336) (1.7533333,-3.0747505) (1.7600000,-3.0917331) (1.7666667,-3.1087817) (1.7733333,-3.1258966) (1.7800000,-3.1430783) (1.7866667,-3.1603272) (1.7933333,-3.1776435) (1.8000000,-3.1950277) (1.8066667,-3.2124801) (1.8133333,-3.2300011) (1.8200000,-3.2475910) (1.8266667,-3.2652503) (1.8333333,-3.2829794) (1.8400000,-3.3007785) (1.8466667,-3.3186481) (1.8533333,-3.3365886) (1.8600000,-3.3546003) (1.8666667,-3.3726837) (1.8733333,-3.3908391) (1.8800000,-3.4090669) (1.8866667,-3.4273675) (1.8933333,-3.4457412) (1.9000000,-3.4641886) (1.9066667,-3.4827099) (1.9133333,-3.5013056) (1.9200000,-3.5199761) (1.9266667,-3.5387217) (1.9333333,-3.5575429) (1.9400000,-3.5764400) (1.9466667,-3.5954136) (1.9533333,-3.6144639) (1.9600000,-3.6335913) (1.9666667,-3.6527964) (1.9733333,-3.6720795) (1.9800000,-3.6914409) (1.9866667,-3.7108813) (1.9933333,-3.7304008) (2.0000000,-3.7500000)};
\draw[violet!75!black,densely dotted] (-1,-1.5000000)--(1,1.5000000);
\fill[violet!75!black] (-1,-1.5000000) circle(2.3pt);
\fill[violet!75!black] (1,1.5000000) circle(2.3pt);
\end{scope}
\end{scope}
\node[anchor=west,align=left,text width=4.35cm] at (8.9,4.4) {蓝实线：$a=2$\\[8pt]橙虚线：$a=\frac12$\\[14pt]两种底数下，\\点对均关于原点对称。};
\node[text width=12.8cm,align=center] at (7.1,.65) {$D=\mathbb R$；$f(0)=0$。指数互换后，差变号。};
\end{tikzpicture}\end{center}


先检查定义域，再对任意允许的 $x$ 计算。图中的点对提示“差反号”；一般证明为
\[
\begin{aligned}
a^{-x}-a^x&=-(a^x-a^{-x}),\\
a^{-x}+a^x&=a^x+a^{-x}.
\end{aligned}
\]
所以奇1为奇函数、偶1为偶函数，均在 $\mathbb R$ 上有定义。
同样，$|-x|=|x|$、$-(-x)^2=-x^2$，得到
\[
a^{|-x|}=a^{|x|},\quad a^{-(-x)^2}=a^{-x^2},\quad
\log_a|-x|=\log_a|x|\ (x\ne0).
\]
故偶2、偶3、偶4均为偶函数。对数那一式先排除零，不能省略。
\end{proofbox}
\clearpage
\begin{proofbox}
\textbf{\textcolor{CProof}{证明二：倒数代换——先处理指数分式}}
% Formula-based samples; see F083 checks/build_tikz_diagrams.py.
\begin{center}\begin{tikzpicture}[font=\small,>=stealth,line join=round]
\draw[gray!40,rounded corners=3pt] (0,0) rectangle (14.2,8.4);
\node at (7.1,7.75) {奇3：$y=\dfrac{a^x+1}{a^x-1}$};
\begin{scope}[shift={(4.400000,3.900000)},x=0.850000cm,y=0.450000cm]
\draw[gray,dotted] (-4,-1)--(4,-1);
\node[anchor=east] at (-4-.08/0.850000,-1) {$-1$};
\draw[gray,dotted] (-4,1)--(4,1);
\node[anchor=east] at (-4-.08/0.850000,1) {$1$};
\draw[->,gray] (-4,0)--(4,0) node[right] {$x$};
\draw[->,gray] (0,-5)--(0,5) node[above] {$y$};
\draw[gray] (-1,-.035/0.450000)--(-1,.035/0.450000);
\node[anchor=north] at (-1,-.07/0.450000) {$-1$};
\draw[gray] (1,-.035/0.450000)--(1,.035/0.450000);
\node[anchor=north] at (1,-.07/0.450000) {$1$};
\begin{scope}
\clip (-4,-5) rectangle (4,5);
\draw[blue,thick,solid] plot coordinates {(-4.0000000,-1.1333333) (-3.9933333,-1.1339923) (-3.9866667,-1.1346547) (-3.9800000,-1.1353205) (-3.9733333,-1.1359899) (-3.9666667,-1.1366628) (-3.9600000,-1.1373393) (-3.9533333,-1.1380193) (-3.9466667,-1.1387029) (-3.9400000,-1.1393901) (-3.9333333,-1.1400809) (-3.9266667,-1.1407754) (-3.9200000,-1.1414735) (-3.9133333,-1.1421753) (-3.9066667,-1.1428809) (-3.9000000,-1.1435902) (-3.8933333,-1.1443032) (-3.8866667,-1.1450201) (-3.8800000,-1.1457407) (-3.8733333,-1.1464652) (-3.8666667,-1.1471935) (-3.8600000,-1.1479256) (-3.8533333,-1.1486617) (-3.8466667,-1.1494017) (-3.8400000,-1.1501456) (-3.8333333,-1.1508935) (-3.8266667,-1.1516454) (-3.8200000,-1.1524013) (-3.8133333,-1.1531612) (-3.8066667,-1.1539252) (-3.8000000,-1.1546933) (-3.7933333,-1.1554655) (-3.7866667,-1.1562418) (-3.7800000,-1.1570223) (-3.7733333,-1.1578069) (-3.7666667,-1.1585958) (-3.7600000,-1.1593889) (-3.7533333,-1.1601863) (-3.7466667,-1.1609879) (-3.7400000,-1.1617939) (-3.7333333,-1.1626042) (-3.7266667,-1.1634189) (-3.7200000,-1.1642379) (-3.7133333,-1.1650614) (-3.7066667,-1.1658893) (-3.7000000,-1.1667217) (-3.6933333,-1.1675586) (-3.6866667,-1.1684000) (-3.6800000,-1.1692460) (-3.6733333,-1.1700966) (-3.6666667,-1.1709517) (-3.6600000,-1.1718115) (-3.6533333,-1.1726760) (-3.6466667,-1.1735452) (-3.6400000,-1.1744191) (-3.6333333,-1.1752978) (-3.6266667,-1.1761812) (-3.6200000,-1.1770695) (-3.6133333,-1.1779626) (-3.6066667,-1.1788606) (-3.6000000,-1.1797634) (-3.5933333,-1.1806713) (-3.5866667,-1.1815841) (-3.5800000,-1.1825018) (-3.5733333,-1.1834246) (-3.5666667,-1.1843525) (-3.5600000,-1.1852855) (-3.5533333,-1.1862236) (-3.5466667,-1.1871668) (-3.5400000,-1.1881153) (-3.5333333,-1.1890689) (-3.5266667,-1.1900278) (-3.5200000,-1.1909920) (-3.5133333,-1.1919615) (-3.5066667,-1.1929364) (-3.5000000,-1.1939167) (-3.4933333,-1.1949024) (-3.4866667,-1.1958935) (-3.4800000,-1.1968901) (-3.4733333,-1.1978923) (-3.4666667,-1.1989000) (-3.4600000,-1.1999133) (-3.4533333,-1.2009323) (-3.4466667,-1.2019569) (-3.4400000,-1.2029873) (-3.4333333,-1.2040233) (-3.4266667,-1.2050652) (-3.4200000,-1.2061129) (-3.4133333,-1.2071664) (-3.4066667,-1.2082258) (-3.4000000,-1.2092912) (-3.3933333,-1.2103625) (-3.3866667,-1.2114399) (-3.3800000,-1.2125233) (-3.3733333,-1.2136127) (-3.3666667,-1.2147083) (-3.3600000,-1.2158101) (-3.3533333,-1.2169181) (-3.3466667,-1.2180323) (-3.3400000,-1.2191528) (-3.3333333,-1.2202796) (-3.3266667,-1.2214129) (-3.3200000,-1.2225525) (-3.3133333,-1.2236986) (-3.3066667,-1.2248511) (-3.3000000,-1.2260103) (-3.2933333,-1.2271760) (-3.2866667,-1.2283483) (-3.2800000,-1.2295273) (-3.2733333,-1.2307131) (-3.2666667,-1.2319056) (-3.2600000,-1.2331049) (-3.2533333,-1.2343110) (-3.2466667,-1.2355241) (-3.2400000,-1.2367441) (-3.2333333,-1.2379711) (-3.2266667,-1.2392051) (-3.2200000,-1.2404462) (-3.2133333,-1.2416945) (-3.2066667,-1.2429499) (-3.2000000,-1.2442126) (-3.1933333,-1.2454825) (-3.1866667,-1.2467598) (-3.1800000,-1.2480444) (-3.1733333,-1.2493365) (-3.1666667,-1.2506361) (-3.1600000,-1.2519432) (-3.1533333,-1.2532579) (-3.1466667,-1.2545802) (-3.1400000,-1.2559102) (-3.1333333,-1.2572480) (-3.1266667,-1.2585935) (-3.1200000,-1.2599469) (-3.1133333,-1.2613082) (-3.1066667,-1.2626775) (-3.1000000,-1.2640547) (-3.0933333,-1.2654401) (-3.0866667,-1.2668335) (-3.0800000,-1.2682352) (-3.0733333,-1.2696450) (-3.0666667,-1.2710632) (-3.0600000,-1.2724897) (-3.0533333,-1.2739247) (-3.0466667,-1.2753681) (-3.0400000,-1.2768200) (-3.0333333,-1.2782806) (-3.0266667,-1.2797498) (-3.0200000,-1.2812277) (-3.0133333,-1.2827143) (-3.0066667,-1.2842099) (-3.0000000,-1.2857143) (-2.9933333,-1.2872277) (-2.9866667,-1.2887501) (-2.9800000,-1.2902816) (-2.9733333,-1.2918223) (-2.9666667,-1.2933721) (-2.9600000,-1.2949313) (-2.9533333,-1.2964999) (-2.9466667,-1.2980778) (-2.9400000,-1.2996653) (-2.9333333,-1.3012623) (-2.9266667,-1.3028690) (-2.9200000,-1.3044853) (-2.9133333,-1.3061114) (-2.9066667,-1.3077474) (-2.9000000,-1.3093933) (-2.8933333,-1.3110492) (-2.8866667,-1.3127151) (-2.8800000,-1.3143912) (-2.8733333,-1.3160775) (-2.8666667,-1.3177741) (-2.8600000,-1.3194810) (-2.8533333,-1.3211984) (-2.8466667,-1.3229263) (-2.8400000,-1.3246648) (-2.8333333,-1.3264139) (-2.8266667,-1.3281738) (-2.8200000,-1.3299446) (-2.8133333,-1.3317263) (-2.8066667,-1.3335189) (-2.8000000,-1.3353227) (-2.7933333,-1.3371376) (-2.7866667,-1.3389637) (-2.7800000,-1.3408012) (-2.7733333,-1.3426501) (-2.7666667,-1.3445106) (-2.7600000,-1.3463826) (-2.7533333,-1.3482663) (-2.7466667,-1.3501617) (-2.7400000,-1.3520691) (-2.7333333,-1.3539884) (-2.7266667,-1.3559197) (-2.7200000,-1.3578632) (-2.7133333,-1.3598189) (-2.7066667,-1.3617869) (-2.7000000,-1.3637674) (-2.6933333,-1.3657604) (-2.6866667,-1.3677659) (-2.6800000,-1.3697843) (-2.6733333,-1.3718154) (-2.6666667,-1.3738594) (-2.6600000,-1.3759165) (-2.6533333,-1.3779867) (-2.6466667,-1.3800701) (-2.6400000,-1.3821668) (-2.6333333,-1.3842770) (-2.6266667,-1.3864007) (-2.6200000,-1.3885381) (-2.6133333,-1.3906892) (-2.6066667,-1.3928542) (-2.6000000,-1.3950331) (-2.5933333,-1.3972262) (-2.5866667,-1.3994335) (-2.5800000,-1.4016551) (-2.5733333,-1.4038911) (-2.5666667,-1.4061417) (-2.5600000,-1.4084069) (-2.5533333,-1.4106869) (-2.5466667,-1.4129819) (-2.5400000,-1.4152919) (-2.5333333,-1.4176170) (-2.5266667,-1.4199574) (-2.5200000,-1.4223133) (-2.5133333,-1.4246846) (-2.5066667,-1.4270716) (-2.5000000,-1.4294745) (-2.4933333,-1.4318932) (-2.4866667,-1.4343280) (-2.4800000,-1.4367790) (-2.4733333,-1.4392463) (-2.4666667,-1.4417301) (-2.4600000,-1.4442305) (-2.4533333,-1.4467476) (-2.4466667,-1.4492816) (-2.4400000,-1.4518327) (-2.4333333,-1.4544009) (-2.4266667,-1.4569864) (-2.4200000,-1.4595894) (-2.4133333,-1.4622100) (-2.4066667,-1.4648484) (-2.4000000,-1.4675047) (-2.3933333,-1.4701791) (-2.3866667,-1.4728717) (-2.3800000,-1.4755827) (-2.3733333,-1.4783122) (-2.3666667,-1.4810605) (-2.3600000,-1.4838276) (-2.3533333,-1.4866138) (-2.3466667,-1.4894192) (-2.3400000,-1.4922439) (-2.3333333,-1.4950882) (-2.3266667,-1.4979522) (-2.3200000,-1.5008361) (-2.3133333,-1.5037401) (-2.3066667,-1.5066643) (-2.3000000,-1.5096090) (-2.2933333,-1.5125742) (-2.2866667,-1.5155603) (-2.2800000,-1.5185674) (-2.2733333,-1.5215956) (-2.2666667,-1.5246452) (-2.2600000,-1.5277163) (-2.2533333,-1.5308093) (-2.2466667,-1.5339241) (-2.2400000,-1.5370612) (-2.2333333,-1.5402206) (-2.2266667,-1.5434025) (-2.2200000,-1.5466072) (-2.2133333,-1.5498350) (-2.2066667,-1.5530859) (-2.2000000,-1.5563602) (-2.1933333,-1.5596582) (-2.1866667,-1.5629800) (-2.1800000,-1.5663259) (-2.1733333,-1.5696961) (-2.1666667,-1.5730908) (-2.1600000,-1.5765103) (-2.1533333,-1.5799548) (-2.1466667,-1.5834246) (-2.1400000,-1.5869198) (-2.1333333,-1.5904408) (-2.1266667,-1.5939877) (-2.1200000,-1.5975609) (-2.1133333,-1.6011605) (-2.1066667,-1.6047868) (-2.1000000,-1.6084402) (-2.0933333,-1.6121208) (-2.0866667,-1.6158289) (-2.0800000,-1.6195648) (-2.0733333,-1.6233288) (-2.0666667,-1.6271211) (-2.0600000,-1.6309421) (-2.0533333,-1.6347919) (-2.0466667,-1.6386710) (-2.0400000,-1.6425795) (-2.0333333,-1.6465178) (-2.0266667,-1.6504862) (-2.0200000,-1.6544850) (-2.0133333,-1.6585144) (-2.0066667,-1.6625749) (-2.0000000,-1.6666667) (-1.9933333,-1.6707901) (-1.9866667,-1.6749454) (-1.9800000,-1.6791331) (-1.9733333,-1.6833534) (-1.9666667,-1.6876066) (-1.9600000,-1.6918931) (-1.9533333,-1.6962133) (-1.9466667,-1.7005675) (-1.9400000,-1.7049560) (-1.9333333,-1.7093792) (-1.9266667,-1.7138375) (-1.9200000,-1.7183313) (-1.9133333,-1.7228609) (-1.9066667,-1.7274267) (-1.9000000,-1.7320291) (-1.8933333,-1.7366684) (-1.8866667,-1.7413452) (-1.8800000,-1.7460597) (-1.8733333,-1.7508124) (-1.8666667,-1.7556038) (-1.8600000,-1.7604341) (-1.8533333,-1.7653039) (-1.8466667,-1.7702136) (-1.8400000,-1.7751637) (-1.8333333,-1.7801545) (-1.8266667,-1.7851865) (-1.8200000,-1.7902602) (-1.8133333,-1.7953761) (-1.8066667,-1.8005345) (-1.8000000,-1.8057361) (-1.7933333,-1.8109813) (-1.7866667,-1.8162705) (-1.7800000,-1.8216043) (-1.7733333,-1.8269832) (-1.7666667,-1.8324078) (-1.7600000,-1.8378784) (-1.7533333,-1.8433957) (-1.7466667,-1.8489602) (-1.7400000,-1.8545724) (-1.7333333,-1.8602330) (-1.7266667,-1.8659424) (-1.7200000,-1.8717012) (-1.7133333,-1.8775101) (-1.7066667,-1.8833696) (-1.7000000,-1.8892803) (-1.6933333,-1.8952429) (-1.6866667,-1.9012579) (-1.6800000,-1.9073259) (-1.6733333,-1.9134477) (-1.6666667,-1.9196239) (-1.6600000,-1.9258551) (-1.6533333,-1.9321420) (-1.6466667,-1.9384853) (-1.6400000,-1.9448856) (-1.6333333,-1.9513438) (-1.6266667,-1.9578604) (-1.6200000,-1.9644362) (-1.6133333,-1.9710721) (-1.6066667,-1.9777686) (-1.6000000,-1.9845266) (-1.5933333,-1.9913469) (-1.5866667,-1.9982302) (-1.5800000,-2.0051774) (-1.5733333,-2.0121893) (-1.5666667,-2.0192666) (-1.5600000,-2.0264103) (-1.5533333,-2.0336213) (-1.5466667,-2.0409003) (-1.5400000,-2.0482482) (-1.5333333,-2.0556661) (-1.5266667,-2.0631547) (-1.5200000,-2.0707151) (-1.5133333,-2.0783482) (-1.5066667,-2.0860549) (-1.5000000,-2.0938363) (-1.4933333,-2.1016934) (-1.4866667,-2.1096271) (-1.4800000,-2.1176385) (-1.4733333,-2.1257288) (-1.4666667,-2.1338988) (-1.4600000,-2.1421499) (-1.4533333,-2.1504830) (-1.4466667,-2.1588994) (-1.4400000,-2.1674001) (-1.4333333,-2.1759864) (-1.4266667,-2.1846595) (-1.4200000,-2.1934206) (-1.4133333,-2.2022709) (-1.4066667,-2.2112118) (-1.4000000,-2.2202445) (-1.3933333,-2.2293703) (-1.3866667,-2.2385906) (-1.3800000,-2.2479069) (-1.3733333,-2.2573204) (-1.3666667,-2.2668325) (-1.3600000,-2.2764449) (-1.3533333,-2.2861589) (-1.3466667,-2.2959761) (-1.3400000,-2.3058980) (-1.3333333,-2.3159262) (-1.3266667,-2.3260623) (-1.3200000,-2.3363079) (-1.3133333,-2.3466647) (-1.3066667,-2.3571344) (-1.3000000,-2.3677189) (-1.2933333,-2.3784197) (-1.2866667,-2.3892388) (-1.2800000,-2.4001780) (-1.2733333,-2.4112393) (-1.2666667,-2.4224244) (-1.2600000,-2.4337355) (-1.2533333,-2.4451745) (-1.2466667,-2.4567435) (-1.2400000,-2.4684446) (-1.2333333,-2.4802799) (-1.2266667,-2.4922517) (-1.2200000,-2.5043621) (-1.2133333,-2.5166135) (-1.2066667,-2.5290081) (-1.2000000,-2.5415485) (-1.1933333,-2.5542371) (-1.1866667,-2.5670763) (-1.1800000,-2.5800687) (-1.1733333,-2.5932169) (-1.1666667,-2.6065236) (-1.1600000,-2.6199916) (-1.1533333,-2.6336237) (-1.1466667,-2.6474226) (-1.1400000,-2.6613914) (-1.1333333,-2.6755331) (-1.1266667,-2.6898507) (-1.1200000,-2.7043473) (-1.1133333,-2.7190263) (-1.1066667,-2.7338909) (-1.1000000,-2.7489444) (-1.0933333,-2.7641904) (-1.0866667,-2.7796324) (-1.0800000,-2.7952740) (-1.0733333,-2.8111190) (-1.0666667,-2.8271711) (-1.0600000,-2.8434343) (-1.0533333,-2.8599126) (-1.0466667,-2.8766102) (-1.0400000,-2.8935312) (-1.0333333,-2.9106799) (-1.0266667,-2.9280609) (-1.0200000,-2.9456787) (-1.0133333,-2.9635379) (-1.0066667,-2.9816433) (-1.0000000,-3.0000000) (-0.9933333,-3.0186129) (-0.9866667,-3.0374873) (-0.9800000,-3.0566284) (-0.9733333,-3.0760419) (-0.9666667,-3.0957332) (-0.9600000,-3.1157083) (-0.9533333,-3.1359731) (-0.9466667,-3.1565336) (-0.9400000,-3.1773963) (-0.9333333,-3.1985676) (-0.9266667,-3.2200541) (-0.9200000,-3.2418627) (-0.9133333,-3.2640005) (-0.9066667,-3.2864747) (-0.9000000,-3.3092929) (-0.8933333,-3.3324626) (-0.8866667,-3.3559920) (-0.8800000,-3.3798891) (-0.8733333,-3.4041623) (-0.8666667,-3.4288205) (-0.8600000,-3.4538724) (-0.8533333,-3.4793274) (-0.8466667,-3.5051949) (-0.8400000,-3.5314849) (-0.8333333,-3.5582074) (-0.8266667,-3.5853730) (-0.8200000,-3.6129923) (-0.8133333,-3.6410767) (-0.8066667,-3.6696377) (-0.8000000,-3.6986870) (-0.7933333,-3.7282372) (-0.7866667,-3.7583009) (-0.7800000,-3.7888912) (-0.7733333,-3.8200219) (-0.7666667,-3.8517070) (-0.7600000,-3.8839611) (-0.7533333,-3.9167994) (-0.7466667,-3.9502374) (-0.7400000,-3.9842914) (-0.7333333,-4.0189782) (-0.7266667,-4.0543153) (-0.7200000,-4.0903206) (-0.7133333,-4.1270130) (-0.7066667,-4.1644119) (-0.7000000,-4.2025375) (-0.6933333,-4.2414107) (-0.6866667,-4.2810533) (-0.6800000,-4.3214880) (-0.6733333,-4.3627384) (-0.6666667,-4.4048288) (-0.6600000,-4.4477847) (-0.6533333,-4.4916327) (-0.6466667,-4.5364003) (-0.6400000,-4.5821163) (-0.6333333,-4.6288107) (-0.6266667,-4.6765146) (-0.6200000,-4.7252607) (-0.6133333,-4.7750829) (-0.6066667,-4.8260167) (-0.6000000,-4.8780992) (-0.5933333,-4.9313691) (-0.5866667,-4.9858669)};
\draw[blue,thick,solid] plot coordinates {(0.5866667,4.9858669) (0.5933333,4.9313691) (0.6000000,4.8780992) (0.6066667,4.8260167) (0.6133333,4.7750829) (0.6200000,4.7252607) (0.6266667,4.6765146) (0.6333333,4.6288107) (0.6400000,4.5821163) (0.6466667,4.5364003) (0.6533333,4.4916327) (0.6600000,4.4477847) (0.6666667,4.4048288) (0.6733333,4.3627384) (0.6800000,4.3214880) (0.6866667,4.2810533) (0.6933333,4.2414107) (0.7000000,4.2025375) (0.7066667,4.1644119) (0.7133333,4.1270130) (0.7200000,4.0903206) (0.7266667,4.0543153) (0.7333333,4.0189782) (0.7400000,3.9842914) (0.7466667,3.9502374) (0.7533333,3.9167994) (0.7600000,3.8839611) (0.7666667,3.8517070) (0.7733333,3.8200219) (0.7800000,3.7888912) (0.7866667,3.7583009) (0.7933333,3.7282372) (0.8000000,3.6986870) (0.8066667,3.6696377) (0.8133333,3.6410767) (0.8200000,3.6129923) (0.8266667,3.5853730) (0.8333333,3.5582074) (0.8400000,3.5314849) (0.8466667,3.5051949) (0.8533333,3.4793274) (0.8600000,3.4538724) (0.8666667,3.4288205) (0.8733333,3.4041623) (0.8800000,3.3798891) (0.8866667,3.3559920) (0.8933333,3.3324626) (0.9000000,3.3092929) (0.9066667,3.2864747) (0.9133333,3.2640005) (0.9200000,3.2418627) (0.9266667,3.2200541) (0.9333333,3.1985676) (0.9400000,3.1773963) (0.9466667,3.1565336) (0.9533333,3.1359731) (0.9600000,3.1157083) (0.9666667,3.0957332) (0.9733333,3.0760419) (0.9800000,3.0566284) (0.9866667,3.0374873) (0.9933333,3.0186129) (1.0000000,3.0000000) (1.0066667,2.9816433) (1.0133333,2.9635379) (1.0200000,2.9456787) (1.0266667,2.9280609) (1.0333333,2.9106799) (1.0400000,2.8935312) (1.0466667,2.8766102) (1.0533333,2.8599126) (1.0600000,2.8434343) (1.0666667,2.8271711) (1.0733333,2.8111190) (1.0800000,2.7952740) (1.0866667,2.7796324) (1.0933333,2.7641904) (1.1000000,2.7489444) (1.1066667,2.7338909) (1.1133333,2.7190263) (1.1200000,2.7043473) (1.1266667,2.6898507) (1.1333333,2.6755331) (1.1400000,2.6613914) (1.1466667,2.6474226) (1.1533333,2.6336237) (1.1600000,2.6199916) (1.1666667,2.6065236) (1.1733333,2.5932169) (1.1800000,2.5800687) (1.1866667,2.5670763) (1.1933333,2.5542371) (1.2000000,2.5415485) (1.2066667,2.5290081) (1.2133333,2.5166135) (1.2200000,2.5043621) (1.2266667,2.4922517) (1.2333333,2.4802799) (1.2400000,2.4684446) (1.2466667,2.4567435) (1.2533333,2.4451745) (1.2600000,2.4337355) (1.2666667,2.4224244) (1.2733333,2.4112393) (1.2800000,2.4001780) (1.2866667,2.3892388) (1.2933333,2.3784197) (1.3000000,2.3677189) (1.3066667,2.3571344) (1.3133333,2.3466647) (1.3200000,2.3363079) (1.3266667,2.3260623) (1.3333333,2.3159262) (1.3400000,2.3058980) (1.3466667,2.2959761) (1.3533333,2.2861589) (1.3600000,2.2764449) (1.3666667,2.2668325) (1.3733333,2.2573204) (1.3800000,2.2479069) (1.3866667,2.2385906) (1.3933333,2.2293703) (1.4000000,2.2202445) (1.4066667,2.2112118) (1.4133333,2.2022709) (1.4200000,2.1934206) (1.4266667,2.1846595) (1.4333333,2.1759864) (1.4400000,2.1674001) (1.4466667,2.1588994) (1.4533333,2.1504830) (1.4600000,2.1421499) (1.4666667,2.1338988) (1.4733333,2.1257288) (1.4800000,2.1176385) (1.4866667,2.1096271) (1.4933333,2.1016934) (1.5000000,2.0938363) (1.5066667,2.0860549) (1.5133333,2.0783482) (1.5200000,2.0707151) (1.5266667,2.0631547) (1.5333333,2.0556661) (1.5400000,2.0482482) (1.5466667,2.0409003) (1.5533333,2.0336213) (1.5600000,2.0264103) (1.5666667,2.0192666) (1.5733333,2.0121893) (1.5800000,2.0051774) (1.5866667,1.9982302) (1.5933333,1.9913469) (1.6000000,1.9845266) (1.6066667,1.9777686) (1.6133333,1.9710721) (1.6200000,1.9644362) (1.6266667,1.9578604) (1.6333333,1.9513438) (1.6400000,1.9448856) (1.6466667,1.9384853) (1.6533333,1.9321420) (1.6600000,1.9258551) (1.6666667,1.9196239) (1.6733333,1.9134477) (1.6800000,1.9073259) (1.6866667,1.9012579) (1.6933333,1.8952429) (1.7000000,1.8892803) (1.7066667,1.8833696) (1.7133333,1.8775101) (1.7200000,1.8717012) (1.7266667,1.8659424) (1.7333333,1.8602330) (1.7400000,1.8545724) (1.7466667,1.8489602) (1.7533333,1.8433957) (1.7600000,1.8378784) (1.7666667,1.8324078) (1.7733333,1.8269832) (1.7800000,1.8216043) (1.7866667,1.8162705) (1.7933333,1.8109813) (1.8000000,1.8057361) (1.8066667,1.8005345) (1.8133333,1.7953761) (1.8200000,1.7902602) (1.8266667,1.7851865) (1.8333333,1.7801545) (1.8400000,1.7751637) (1.8466667,1.7702136) (1.8533333,1.7653039) (1.8600000,1.7604341) (1.8666667,1.7556038) (1.8733333,1.7508124) (1.8800000,1.7460597) (1.8866667,1.7413452) (1.8933333,1.7366684) (1.9000000,1.7320291) (1.9066667,1.7274267) (1.9133333,1.7228609) (1.9200000,1.7183313) (1.9266667,1.7138375) (1.9333333,1.7093792) (1.9400000,1.7049560) (1.9466667,1.7005675) (1.9533333,1.6962133) (1.9600000,1.6918931) (1.9666667,1.6876066) (1.9733333,1.6833534) (1.9800000,1.6791331) (1.9866667,1.6749454) (1.9933333,1.6707901) (2.0000000,1.6666667) (2.0066667,1.6625749) (2.0133333,1.6585144) (2.0200000,1.6544850) (2.0266667,1.6504862) (2.0333333,1.6465178) (2.0400000,1.6425795) (2.0466667,1.6386710) (2.0533333,1.6347919) (2.0600000,1.6309421) (2.0666667,1.6271211) (2.0733333,1.6233288) (2.0800000,1.6195648) (2.0866667,1.6158289) (2.0933333,1.6121208) (2.1000000,1.6084402) (2.1066667,1.6047868) (2.1133333,1.6011605) (2.1200000,1.5975609) (2.1266667,1.5939877) (2.1333333,1.5904408) (2.1400000,1.5869198) (2.1466667,1.5834246) (2.1533333,1.5799548) (2.1600000,1.5765103) (2.1666667,1.5730908) (2.1733333,1.5696961) (2.1800000,1.5663259) (2.1866667,1.5629800) (2.1933333,1.5596582) (2.2000000,1.5563602) (2.2066667,1.5530859) (2.2133333,1.5498350) (2.2200000,1.5466072) (2.2266667,1.5434025) (2.2333333,1.5402206) (2.2400000,1.5370612) (2.2466667,1.5339241) (2.2533333,1.5308093) (2.2600000,1.5277163) (2.2666667,1.5246452) (2.2733333,1.5215956) (2.2800000,1.5185674) (2.2866667,1.5155603) (2.2933333,1.5125742) (2.3000000,1.5096090) (2.3066667,1.5066643) (2.3133333,1.5037401) (2.3200000,1.5008361) (2.3266667,1.4979522) (2.3333333,1.4950882) (2.3400000,1.4922439) (2.3466667,1.4894192) (2.3533333,1.4866138) (2.3600000,1.4838276) (2.3666667,1.4810605) (2.3733333,1.4783122) (2.3800000,1.4755827) (2.3866667,1.4728717) (2.3933333,1.4701791) (2.4000000,1.4675047) (2.4066667,1.4648484) (2.4133333,1.4622100) (2.4200000,1.4595894) (2.4266667,1.4569864) (2.4333333,1.4544009) (2.4400000,1.4518327) (2.4466667,1.4492816) (2.4533333,1.4467476) (2.4600000,1.4442305) (2.4666667,1.4417301) (2.4733333,1.4392463) (2.4800000,1.4367790) (2.4866667,1.4343280) (2.4933333,1.4318932) (2.5000000,1.4294745) (2.5066667,1.4270716) (2.5133333,1.4246846) (2.5200000,1.4223133) (2.5266667,1.4199574) (2.5333333,1.4176170) (2.5400000,1.4152919) (2.5466667,1.4129819) (2.5533333,1.4106869) (2.5600000,1.4084069) (2.5666667,1.4061417) (2.5733333,1.4038911) (2.5800000,1.4016551) (2.5866667,1.3994335) (2.5933333,1.3972262) (2.6000000,1.3950331) (2.6066667,1.3928542) (2.6133333,1.3906892) (2.6200000,1.3885381) (2.6266667,1.3864007) (2.6333333,1.3842770) (2.6400000,1.3821668) (2.6466667,1.3800701) (2.6533333,1.3779867) (2.6600000,1.3759165) (2.6666667,1.3738594) (2.6733333,1.3718154) (2.6800000,1.3697843) (2.6866667,1.3677659) (2.6933333,1.3657604) (2.7000000,1.3637674) (2.7066667,1.3617869) (2.7133333,1.3598189) (2.7200000,1.3578632) (2.7266667,1.3559197) (2.7333333,1.3539884) (2.7400000,1.3520691) (2.7466667,1.3501617) (2.7533333,1.3482663) (2.7600000,1.3463826) (2.7666667,1.3445106) (2.7733333,1.3426501) (2.7800000,1.3408012) (2.7866667,1.3389637) (2.7933333,1.3371376) (2.8000000,1.3353227) (2.8066667,1.3335189) (2.8133333,1.3317263) (2.8200000,1.3299446) (2.8266667,1.3281738) (2.8333333,1.3264139) (2.8400000,1.3246648) (2.8466667,1.3229263) (2.8533333,1.3211984) (2.8600000,1.3194810) (2.8666667,1.3177741) (2.8733333,1.3160775) (2.8800000,1.3143912) (2.8866667,1.3127151) (2.8933333,1.3110492) (2.9000000,1.3093933) (2.9066667,1.3077474) (2.9133333,1.3061114) (2.9200000,1.3044853) (2.9266667,1.3028690) (2.9333333,1.3012623) (2.9400000,1.2996653) (2.9466667,1.2980778) (2.9533333,1.2964999) (2.9600000,1.2949313) (2.9666667,1.2933721) (2.9733333,1.2918223) (2.9800000,1.2902816) (2.9866667,1.2887501) (2.9933333,1.2872277) (3.0000000,1.2857143) (3.0066667,1.2842099) (3.0133333,1.2827143) (3.0200000,1.2812277) (3.0266667,1.2797498) (3.0333333,1.2782806) (3.0400000,1.2768200) (3.0466667,1.2753681) (3.0533333,1.2739247) (3.0600000,1.2724897) (3.0666667,1.2710632) (3.0733333,1.2696450) (3.0800000,1.2682352) (3.0866667,1.2668335) (3.0933333,1.2654401) (3.1000000,1.2640547) (3.1066667,1.2626775) (3.1133333,1.2613082) (3.1200000,1.2599469) (3.1266667,1.2585935) (3.1333333,1.2572480) (3.1400000,1.2559102) (3.1466667,1.2545802) (3.1533333,1.2532579) (3.1600000,1.2519432) (3.1666667,1.2506361) (3.1733333,1.2493365) (3.1800000,1.2480444) (3.1866667,1.2467598) (3.1933333,1.2454825) (3.2000000,1.2442126) (3.2066667,1.2429499) (3.2133333,1.2416945) (3.2200000,1.2404462) (3.2266667,1.2392051) (3.2333333,1.2379711) (3.2400000,1.2367441) (3.2466667,1.2355241) (3.2533333,1.2343110) (3.2600000,1.2331049) (3.2666667,1.2319056) (3.2733333,1.2307131) (3.2800000,1.2295273) (3.2866667,1.2283483) (3.2933333,1.2271760) (3.3000000,1.2260103) (3.3066667,1.2248511) (3.3133333,1.2236986) (3.3200000,1.2225525) (3.3266667,1.2214129) (3.3333333,1.2202796) (3.3400000,1.2191528) (3.3466667,1.2180323) (3.3533333,1.2169181) (3.3600000,1.2158101) (3.3666667,1.2147083) (3.3733333,1.2136127) (3.3800000,1.2125233) (3.3866667,1.2114399) (3.3933333,1.2103625) (3.4000000,1.2092912) (3.4066667,1.2082258) (3.4133333,1.2071664) (3.4200000,1.2061129) (3.4266667,1.2050652) (3.4333333,1.2040233) (3.4400000,1.2029873) (3.4466667,1.2019569) (3.4533333,1.2009323) (3.4600000,1.1999133) (3.4666667,1.1989000) (3.4733333,1.1978923) (3.4800000,1.1968901) (3.4866667,1.1958935) (3.4933333,1.1949024) (3.5000000,1.1939167) (3.5066667,1.1929364) (3.5133333,1.1919615) (3.5200000,1.1909920) (3.5266667,1.1900278) (3.5333333,1.1890689) (3.5400000,1.1881153) (3.5466667,1.1871668) (3.5533333,1.1862236) (3.5600000,1.1852855) (3.5666667,1.1843525) (3.5733333,1.1834246) (3.5800000,1.1825018) (3.5866667,1.1815841) (3.5933333,1.1806713) (3.6000000,1.1797634) (3.6066667,1.1788606) (3.6133333,1.1779626) (3.6200000,1.1770695) (3.6266667,1.1761812) (3.6333333,1.1752978) (3.6400000,1.1744191) (3.6466667,1.1735452) (3.6533333,1.1726760) (3.6600000,1.1718115) (3.6666667,1.1709517) (3.6733333,1.1700966) (3.6800000,1.1692460) (3.6866667,1.1684000) (3.6933333,1.1675586) (3.7000000,1.1667217) (3.7066667,1.1658893) (3.7133333,1.1650614) (3.7200000,1.1642379) (3.7266667,1.1634189) (3.7333333,1.1626042) (3.7400000,1.1617939) (3.7466667,1.1609879) (3.7533333,1.1601863) (3.7600000,1.1593889) (3.7666667,1.1585958) (3.7733333,1.1578069) (3.7800000,1.1570223) (3.7866667,1.1562418) (3.7933333,1.1554655) (3.8000000,1.1546933) (3.8066667,1.1539252) (3.8133333,1.1531612) (3.8200000,1.1524013) (3.8266667,1.1516454) (3.8333333,1.1508935) (3.8400000,1.1501456) (3.8466667,1.1494017) (3.8533333,1.1486617) (3.8600000,1.1479256) (3.8666667,1.1471935) (3.8733333,1.1464652) (3.8800000,1.1457407) (3.8866667,1.1450201) (3.8933333,1.1443032) (3.9000000,1.1435902) (3.9066667,1.1428809) (3.9133333,1.1421753) (3.9200000,1.1414735) (3.9266667,1.1407754) (3.9333333,1.1400809) (3.9400000,1.1393901) (3.9466667,1.1387029) (3.9533333,1.1380193) (3.9600000,1.1373393) (3.9666667,1.1366628) (3.9733333,1.1359899) (3.9800000,1.1353205) (3.9866667,1.1346547) (3.9933333,1.1339923) (4.0000000,1.1333333)};
\draw[orange!85!black,thick,dashed] plot coordinates {(-4.0000000,1.1333333) (-3.9933333,1.1339923) (-3.9866667,1.1346547) (-3.9800000,1.1353205) (-3.9733333,1.1359899) (-3.9666667,1.1366628) (-3.9600000,1.1373393) (-3.9533333,1.1380193) (-3.9466667,1.1387029) (-3.9400000,1.1393901) (-3.9333333,1.1400809) (-3.9266667,1.1407754) (-3.9200000,1.1414735) (-3.9133333,1.1421753) (-3.9066667,1.1428809) (-3.9000000,1.1435902) (-3.8933333,1.1443032) (-3.8866667,1.1450201) (-3.8800000,1.1457407) (-3.8733333,1.1464652) (-3.8666667,1.1471935) (-3.8600000,1.1479256) (-3.8533333,1.1486617) (-3.8466667,1.1494017) (-3.8400000,1.1501456) (-3.8333333,1.1508935) (-3.8266667,1.1516454) (-3.8200000,1.1524013) (-3.8133333,1.1531612) (-3.8066667,1.1539252) (-3.8000000,1.1546933) (-3.7933333,1.1554655) (-3.7866667,1.1562418) (-3.7800000,1.1570223) (-3.7733333,1.1578069) (-3.7666667,1.1585958) (-3.7600000,1.1593889) (-3.7533333,1.1601863) (-3.7466667,1.1609879) (-3.7400000,1.1617939) (-3.7333333,1.1626042) (-3.7266667,1.1634189) (-3.7200000,1.1642379) (-3.7133333,1.1650614) (-3.7066667,1.1658893) (-3.7000000,1.1667217) (-3.6933333,1.1675586) (-3.6866667,1.1684000) (-3.6800000,1.1692460) (-3.6733333,1.1700966) (-3.6666667,1.1709517) (-3.6600000,1.1718115) (-3.6533333,1.1726760) (-3.6466667,1.1735452) (-3.6400000,1.1744191) (-3.6333333,1.1752978) (-3.6266667,1.1761812) (-3.6200000,1.1770695) (-3.6133333,1.1779626) (-3.6066667,1.1788606) (-3.6000000,1.1797634) (-3.5933333,1.1806713) (-3.5866667,1.1815841) (-3.5800000,1.1825018) (-3.5733333,1.1834246) (-3.5666667,1.1843525) (-3.5600000,1.1852855) (-3.5533333,1.1862236) (-3.5466667,1.1871668) (-3.5400000,1.1881153) (-3.5333333,1.1890689) (-3.5266667,1.1900278) (-3.5200000,1.1909920) (-3.5133333,1.1919615) (-3.5066667,1.1929364) (-3.5000000,1.1939167) (-3.4933333,1.1949024) (-3.4866667,1.1958935) (-3.4800000,1.1968901) (-3.4733333,1.1978923) (-3.4666667,1.1989000) (-3.4600000,1.1999133) (-3.4533333,1.2009323) (-3.4466667,1.2019569) (-3.4400000,1.2029873) (-3.4333333,1.2040233) (-3.4266667,1.2050652) (-3.4200000,1.2061129) (-3.4133333,1.2071664) (-3.4066667,1.2082258) (-3.4000000,1.2092912) (-3.3933333,1.2103625) (-3.3866667,1.2114399) (-3.3800000,1.2125233) (-3.3733333,1.2136127) (-3.3666667,1.2147083) (-3.3600000,1.2158101) (-3.3533333,1.2169181) (-3.3466667,1.2180323) (-3.3400000,1.2191528) (-3.3333333,1.2202796) (-3.3266667,1.2214129) (-3.3200000,1.2225525) (-3.3133333,1.2236986) (-3.3066667,1.2248511) (-3.3000000,1.2260103) (-3.2933333,1.2271760) (-3.2866667,1.2283483) (-3.2800000,1.2295273) (-3.2733333,1.2307131) (-3.2666667,1.2319056) (-3.2600000,1.2331049) (-3.2533333,1.2343110) (-3.2466667,1.2355241) (-3.2400000,1.2367441) (-3.2333333,1.2379711) (-3.2266667,1.2392051) (-3.2200000,1.2404462) (-3.2133333,1.2416945) (-3.2066667,1.2429499) (-3.2000000,1.2442126) (-3.1933333,1.2454825) (-3.1866667,1.2467598) (-3.1800000,1.2480444) (-3.1733333,1.2493365) (-3.1666667,1.2506361) (-3.1600000,1.2519432) (-3.1533333,1.2532579) (-3.1466667,1.2545802) (-3.1400000,1.2559102) (-3.1333333,1.2572480) (-3.1266667,1.2585935) (-3.1200000,1.2599469) (-3.1133333,1.2613082) (-3.1066667,1.2626775) (-3.1000000,1.2640547) (-3.0933333,1.2654401) (-3.0866667,1.2668335) (-3.0800000,1.2682352) (-3.0733333,1.2696450) (-3.0666667,1.2710632) (-3.0600000,1.2724897) (-3.0533333,1.2739247) (-3.0466667,1.2753681) (-3.0400000,1.2768200) (-3.0333333,1.2782806) (-3.0266667,1.2797498) (-3.0200000,1.2812277) (-3.0133333,1.2827143) (-3.0066667,1.2842099) (-3.0000000,1.2857143) (-2.9933333,1.2872277) (-2.9866667,1.2887501) (-2.9800000,1.2902816) (-2.9733333,1.2918223) (-2.9666667,1.2933721) (-2.9600000,1.2949313) (-2.9533333,1.2964999) (-2.9466667,1.2980778) (-2.9400000,1.2996653) (-2.9333333,1.3012623) (-2.9266667,1.3028690) (-2.9200000,1.3044853) (-2.9133333,1.3061114) (-2.9066667,1.3077474) (-2.9000000,1.3093933) (-2.8933333,1.3110492) (-2.8866667,1.3127151) (-2.8800000,1.3143912) (-2.8733333,1.3160775) (-2.8666667,1.3177741) (-2.8600000,1.3194810) (-2.8533333,1.3211984) (-2.8466667,1.3229263) (-2.8400000,1.3246648) (-2.8333333,1.3264139) (-2.8266667,1.3281738) (-2.8200000,1.3299446) (-2.8133333,1.3317263) (-2.8066667,1.3335189) (-2.8000000,1.3353227) (-2.7933333,1.3371376) (-2.7866667,1.3389637) (-2.7800000,1.3408012) (-2.7733333,1.3426501) (-2.7666667,1.3445106) (-2.7600000,1.3463826) (-2.7533333,1.3482663) (-2.7466667,1.3501617) (-2.7400000,1.3520691) (-2.7333333,1.3539884) (-2.7266667,1.3559197) (-2.7200000,1.3578632) (-2.7133333,1.3598189) (-2.7066667,1.3617869) (-2.7000000,1.3637674) (-2.6933333,1.3657604) (-2.6866667,1.3677659) (-2.6800000,1.3697843) (-2.6733333,1.3718154) (-2.6666667,1.3738594) (-2.6600000,1.3759165) (-2.6533333,1.3779867) (-2.6466667,1.3800701) (-2.6400000,1.3821668) (-2.6333333,1.3842770) (-2.6266667,1.3864007) (-2.6200000,1.3885381) (-2.6133333,1.3906892) (-2.6066667,1.3928542) (-2.6000000,1.3950331) (-2.5933333,1.3972262) (-2.5866667,1.3994335) (-2.5800000,1.4016551) (-2.5733333,1.4038911) (-2.5666667,1.4061417) (-2.5600000,1.4084069) (-2.5533333,1.4106869) (-2.5466667,1.4129819) (-2.5400000,1.4152919) (-2.5333333,1.4176170) (-2.5266667,1.4199574) (-2.5200000,1.4223133) (-2.5133333,1.4246846) (-2.5066667,1.4270716) (-2.5000000,1.4294745) (-2.4933333,1.4318932) (-2.4866667,1.4343280) (-2.4800000,1.4367790) (-2.4733333,1.4392463) (-2.4666667,1.4417301) (-2.4600000,1.4442305) (-2.4533333,1.4467476) (-2.4466667,1.4492816) (-2.4400000,1.4518327) (-2.4333333,1.4544009) (-2.4266667,1.4569864) (-2.4200000,1.4595894) (-2.4133333,1.4622100) (-2.4066667,1.4648484) (-2.4000000,1.4675047) (-2.3933333,1.4701791) (-2.3866667,1.4728717) (-2.3800000,1.4755827) (-2.3733333,1.4783122) (-2.3666667,1.4810605) (-2.3600000,1.4838276) (-2.3533333,1.4866138) (-2.3466667,1.4894192) (-2.3400000,1.4922439) (-2.3333333,1.4950882) (-2.3266667,1.4979522) (-2.3200000,1.5008361) (-2.3133333,1.5037401) (-2.3066667,1.5066643) (-2.3000000,1.5096090) (-2.2933333,1.5125742) (-2.2866667,1.5155603) (-2.2800000,1.5185674) (-2.2733333,1.5215956) (-2.2666667,1.5246452) (-2.2600000,1.5277163) (-2.2533333,1.5308093) (-2.2466667,1.5339241) (-2.2400000,1.5370612) (-2.2333333,1.5402206) (-2.2266667,1.5434025) (-2.2200000,1.5466072) (-2.2133333,1.5498350) (-2.2066667,1.5530859) (-2.2000000,1.5563602) (-2.1933333,1.5596582) (-2.1866667,1.5629800) (-2.1800000,1.5663259) (-2.1733333,1.5696961) (-2.1666667,1.5730908) (-2.1600000,1.5765103) (-2.1533333,1.5799548) (-2.1466667,1.5834246) (-2.1400000,1.5869198) (-2.1333333,1.5904408) (-2.1266667,1.5939877) (-2.1200000,1.5975609) (-2.1133333,1.6011605) (-2.1066667,1.6047868) (-2.1000000,1.6084402) (-2.0933333,1.6121208) (-2.0866667,1.6158289) (-2.0800000,1.6195648) (-2.0733333,1.6233288) (-2.0666667,1.6271211) (-2.0600000,1.6309421) (-2.0533333,1.6347919) (-2.0466667,1.6386710) (-2.0400000,1.6425795) (-2.0333333,1.6465178) (-2.0266667,1.6504862) (-2.0200000,1.6544850) (-2.0133333,1.6585144) (-2.0066667,1.6625749) (-2.0000000,1.6666667) (-1.9933333,1.6707901) (-1.9866667,1.6749454) (-1.9800000,1.6791331) (-1.9733333,1.6833534) (-1.9666667,1.6876066) (-1.9600000,1.6918931) (-1.9533333,1.6962133) (-1.9466667,1.7005675) (-1.9400000,1.7049560) (-1.9333333,1.7093792) (-1.9266667,1.7138375) (-1.9200000,1.7183313) (-1.9133333,1.7228609) (-1.9066667,1.7274267) (-1.9000000,1.7320291) (-1.8933333,1.7366684) (-1.8866667,1.7413452) (-1.8800000,1.7460597) (-1.8733333,1.7508124) (-1.8666667,1.7556038) (-1.8600000,1.7604341) (-1.8533333,1.7653039) (-1.8466667,1.7702136) (-1.8400000,1.7751637) (-1.8333333,1.7801545) (-1.8266667,1.7851865) (-1.8200000,1.7902602) (-1.8133333,1.7953761) (-1.8066667,1.8005345) (-1.8000000,1.8057361) (-1.7933333,1.8109813) (-1.7866667,1.8162705) (-1.7800000,1.8216043) (-1.7733333,1.8269832) (-1.7666667,1.8324078) (-1.7600000,1.8378784) (-1.7533333,1.8433957) (-1.7466667,1.8489602) (-1.7400000,1.8545724) (-1.7333333,1.8602330) (-1.7266667,1.8659424) (-1.7200000,1.8717012) (-1.7133333,1.8775101) (-1.7066667,1.8833696) (-1.7000000,1.8892803) (-1.6933333,1.8952429) (-1.6866667,1.9012579) (-1.6800000,1.9073259) (-1.6733333,1.9134477) (-1.6666667,1.9196239) (-1.6600000,1.9258551) (-1.6533333,1.9321420) (-1.6466667,1.9384853) (-1.6400000,1.9448856) (-1.6333333,1.9513438) (-1.6266667,1.9578604) (-1.6200000,1.9644362) (-1.6133333,1.9710721) (-1.6066667,1.9777686) (-1.6000000,1.9845266) (-1.5933333,1.9913469) (-1.5866667,1.9982302) (-1.5800000,2.0051774) (-1.5733333,2.0121893) (-1.5666667,2.0192666) (-1.5600000,2.0264103) (-1.5533333,2.0336213) (-1.5466667,2.0409003) (-1.5400000,2.0482482) (-1.5333333,2.0556661) (-1.5266667,2.0631547) (-1.5200000,2.0707151) (-1.5133333,2.0783482) (-1.5066667,2.0860549) (-1.5000000,2.0938363) (-1.4933333,2.1016934) (-1.4866667,2.1096271) (-1.4800000,2.1176385) (-1.4733333,2.1257288) (-1.4666667,2.1338988) (-1.4600000,2.1421499) (-1.4533333,2.1504830) (-1.4466667,2.1588994) (-1.4400000,2.1674001) (-1.4333333,2.1759864) (-1.4266667,2.1846595) (-1.4200000,2.1934206) (-1.4133333,2.2022709) (-1.4066667,2.2112118) (-1.4000000,2.2202445) (-1.3933333,2.2293703) (-1.3866667,2.2385906) (-1.3800000,2.2479069) (-1.3733333,2.2573204) (-1.3666667,2.2668325) (-1.3600000,2.2764449) (-1.3533333,2.2861589) (-1.3466667,2.2959761) (-1.3400000,2.3058980) (-1.3333333,2.3159262) (-1.3266667,2.3260623) (-1.3200000,2.3363079) (-1.3133333,2.3466647) (-1.3066667,2.3571344) (-1.3000000,2.3677189) (-1.2933333,2.3784197) (-1.2866667,2.3892388) (-1.2800000,2.4001780) (-1.2733333,2.4112393) (-1.2666667,2.4224244) (-1.2600000,2.4337355) (-1.2533333,2.4451745) (-1.2466667,2.4567435) (-1.2400000,2.4684446) (-1.2333333,2.4802799) (-1.2266667,2.4922517) (-1.2200000,2.5043621) (-1.2133333,2.5166135) (-1.2066667,2.5290081) (-1.2000000,2.5415485) (-1.1933333,2.5542371) (-1.1866667,2.5670763) (-1.1800000,2.5800687) (-1.1733333,2.5932169) (-1.1666667,2.6065236) (-1.1600000,2.6199916) (-1.1533333,2.6336237) (-1.1466667,2.6474226) (-1.1400000,2.6613914) (-1.1333333,2.6755331) (-1.1266667,2.6898507) (-1.1200000,2.7043473) (-1.1133333,2.7190263) (-1.1066667,2.7338909) (-1.1000000,2.7489444) (-1.0933333,2.7641904) (-1.0866667,2.7796324) (-1.0800000,2.7952740) (-1.0733333,2.8111190) (-1.0666667,2.8271711) (-1.0600000,2.8434343) (-1.0533333,2.8599126) (-1.0466667,2.8766102) (-1.0400000,2.8935312) (-1.0333333,2.9106799) (-1.0266667,2.9280609) (-1.0200000,2.9456787) (-1.0133333,2.9635379) (-1.0066667,2.9816433) (-1.0000000,3.0000000) (-0.9933333,3.0186129) (-0.9866667,3.0374873) (-0.9800000,3.0566284) (-0.9733333,3.0760419) (-0.9666667,3.0957332) (-0.9600000,3.1157083) (-0.9533333,3.1359731) (-0.9466667,3.1565336) (-0.9400000,3.1773963) (-0.9333333,3.1985676) (-0.9266667,3.2200541) (-0.9200000,3.2418627) (-0.9133333,3.2640005) (-0.9066667,3.2864747) (-0.9000000,3.3092929) (-0.8933333,3.3324626) (-0.8866667,3.3559920) (-0.8800000,3.3798891) (-0.8733333,3.4041623) (-0.8666667,3.4288205) (-0.8600000,3.4538724) (-0.8533333,3.4793274) (-0.8466667,3.5051949) (-0.8400000,3.5314849) (-0.8333333,3.5582074) (-0.8266667,3.5853730) (-0.8200000,3.6129923) (-0.8133333,3.6410767) (-0.8066667,3.6696377) (-0.8000000,3.6986870) (-0.7933333,3.7282372) (-0.7866667,3.7583009) (-0.7800000,3.7888912) (-0.7733333,3.8200219) (-0.7666667,3.8517070) (-0.7600000,3.8839611) (-0.7533333,3.9167994) (-0.7466667,3.9502374) (-0.7400000,3.9842914) (-0.7333333,4.0189782) (-0.7266667,4.0543153) (-0.7200000,4.0903206) (-0.7133333,4.1270130) (-0.7066667,4.1644119) (-0.7000000,4.2025375) (-0.6933333,4.2414107) (-0.6866667,4.2810533) (-0.6800000,4.3214880) (-0.6733333,4.3627384) (-0.6666667,4.4048288) (-0.6600000,4.4477847) (-0.6533333,4.4916327) (-0.6466667,4.5364003) (-0.6400000,4.5821163) (-0.6333333,4.6288107) (-0.6266667,4.6765146) (-0.6200000,4.7252607) (-0.6133333,4.7750829) (-0.6066667,4.8260167) (-0.6000000,4.8780992) (-0.5933333,4.9313691) (-0.5866667,4.9858669)};
\draw[orange!85!black,thick,dashed] plot coordinates {(0.5866667,-4.9858669) (0.5933333,-4.9313691) (0.6000000,-4.8780992) (0.6066667,-4.8260167) (0.6133333,-4.7750829) (0.6200000,-4.7252607) (0.6266667,-4.6765146) (0.6333333,-4.6288107) (0.6400000,-4.5821163) (0.6466667,-4.5364003) (0.6533333,-4.4916327) (0.6600000,-4.4477847) (0.6666667,-4.4048288) (0.6733333,-4.3627384) (0.6800000,-4.3214880) (0.6866667,-4.2810533) (0.6933333,-4.2414107) (0.7000000,-4.2025375) (0.7066667,-4.1644119) (0.7133333,-4.1270130) (0.7200000,-4.0903206) (0.7266667,-4.0543153) (0.7333333,-4.0189782) (0.7400000,-3.9842914) (0.7466667,-3.9502374) (0.7533333,-3.9167994) (0.7600000,-3.8839611) (0.7666667,-3.8517070) (0.7733333,-3.8200219) (0.7800000,-3.7888912) (0.7866667,-3.7583009) (0.7933333,-3.7282372) (0.8000000,-3.6986870) (0.8066667,-3.6696377) (0.8133333,-3.6410767) (0.8200000,-3.6129923) (0.8266667,-3.5853730) (0.8333333,-3.5582074) (0.8400000,-3.5314849) (0.8466667,-3.5051949) (0.8533333,-3.4793274) (0.8600000,-3.4538724) (0.8666667,-3.4288205) (0.8733333,-3.4041623) (0.8800000,-3.3798891) (0.8866667,-3.3559920) (0.8933333,-3.3324626) (0.9000000,-3.3092929) (0.9066667,-3.2864747) (0.9133333,-3.2640005) (0.9200000,-3.2418627) (0.9266667,-3.2200541) (0.9333333,-3.1985676) (0.9400000,-3.1773963) (0.9466667,-3.1565336) (0.9533333,-3.1359731) (0.9600000,-3.1157083) (0.9666667,-3.0957332) (0.9733333,-3.0760419) (0.9800000,-3.0566284) (0.9866667,-3.0374873) (0.9933333,-3.0186129) (1.0000000,-3.0000000) (1.0066667,-2.9816433) (1.0133333,-2.9635379) (1.0200000,-2.9456787) (1.0266667,-2.9280609) (1.0333333,-2.9106799) (1.0400000,-2.8935312) (1.0466667,-2.8766102) (1.0533333,-2.8599126) (1.0600000,-2.8434343) (1.0666667,-2.8271711) (1.0733333,-2.8111190) (1.0800000,-2.7952740) (1.0866667,-2.7796324) (1.0933333,-2.7641904) (1.1000000,-2.7489444) (1.1066667,-2.7338909) (1.1133333,-2.7190263) (1.1200000,-2.7043473) (1.1266667,-2.6898507) (1.1333333,-2.6755331) (1.1400000,-2.6613914) (1.1466667,-2.6474226) (1.1533333,-2.6336237) (1.1600000,-2.6199916) (1.1666667,-2.6065236) (1.1733333,-2.5932169) (1.1800000,-2.5800687) (1.1866667,-2.5670763) (1.1933333,-2.5542371) (1.2000000,-2.5415485) (1.2066667,-2.5290081) (1.2133333,-2.5166135) (1.2200000,-2.5043621) (1.2266667,-2.4922517) (1.2333333,-2.4802799) (1.2400000,-2.4684446) (1.2466667,-2.4567435) (1.2533333,-2.4451745) (1.2600000,-2.4337355) (1.2666667,-2.4224244) (1.2733333,-2.4112393) (1.2800000,-2.4001780) (1.2866667,-2.3892388) (1.2933333,-2.3784197) (1.3000000,-2.3677189) (1.3066667,-2.3571344) (1.3133333,-2.3466647) (1.3200000,-2.3363079) (1.3266667,-2.3260623) (1.3333333,-2.3159262) (1.3400000,-2.3058980) (1.3466667,-2.2959761) (1.3533333,-2.2861589) (1.3600000,-2.2764449) (1.3666667,-2.2668325) (1.3733333,-2.2573204) (1.3800000,-2.2479069) (1.3866667,-2.2385906) (1.3933333,-2.2293703) (1.4000000,-2.2202445) (1.4066667,-2.2112118) (1.4133333,-2.2022709) (1.4200000,-2.1934206) (1.4266667,-2.1846595) (1.4333333,-2.1759864) (1.4400000,-2.1674001) (1.4466667,-2.1588994) (1.4533333,-2.1504830) (1.4600000,-2.1421499) (1.4666667,-2.1338988) (1.4733333,-2.1257288) (1.4800000,-2.1176385) (1.4866667,-2.1096271) (1.4933333,-2.1016934) (1.5000000,-2.0938363) (1.5066667,-2.0860549) (1.5133333,-2.0783482) (1.5200000,-2.0707151) (1.5266667,-2.0631547) (1.5333333,-2.0556661) (1.5400000,-2.0482482) (1.5466667,-2.0409003) (1.5533333,-2.0336213) (1.5600000,-2.0264103) (1.5666667,-2.0192666) (1.5733333,-2.0121893) (1.5800000,-2.0051774) (1.5866667,-1.9982302) (1.5933333,-1.9913469) (1.6000000,-1.9845266) (1.6066667,-1.9777686) (1.6133333,-1.9710721) (1.6200000,-1.9644362) (1.6266667,-1.9578604) (1.6333333,-1.9513438) (1.6400000,-1.9448856) (1.6466667,-1.9384853) (1.6533333,-1.9321420) (1.6600000,-1.9258551) (1.6666667,-1.9196239) (1.6733333,-1.9134477) (1.6800000,-1.9073259) (1.6866667,-1.9012579) (1.6933333,-1.8952429) (1.7000000,-1.8892803) (1.7066667,-1.8833696) (1.7133333,-1.8775101) (1.7200000,-1.8717012) (1.7266667,-1.8659424) (1.7333333,-1.8602330) (1.7400000,-1.8545724) (1.7466667,-1.8489602) (1.7533333,-1.8433957) (1.7600000,-1.8378784) (1.7666667,-1.8324078) (1.7733333,-1.8269832) (1.7800000,-1.8216043) (1.7866667,-1.8162705) (1.7933333,-1.8109813) (1.8000000,-1.8057361) (1.8066667,-1.8005345) (1.8133333,-1.7953761) (1.8200000,-1.7902602) (1.8266667,-1.7851865) (1.8333333,-1.7801545) (1.8400000,-1.7751637) (1.8466667,-1.7702136) (1.8533333,-1.7653039) (1.8600000,-1.7604341) (1.8666667,-1.7556038) (1.8733333,-1.7508124) (1.8800000,-1.7460597) (1.8866667,-1.7413452) (1.8933333,-1.7366684) (1.9000000,-1.7320291) (1.9066667,-1.7274267) (1.9133333,-1.7228609) (1.9200000,-1.7183313) (1.9266667,-1.7138375) (1.9333333,-1.7093792) (1.9400000,-1.7049560) (1.9466667,-1.7005675) (1.9533333,-1.6962133) (1.9600000,-1.6918931) (1.9666667,-1.6876066) (1.9733333,-1.6833534) (1.9800000,-1.6791331) (1.9866667,-1.6749454) (1.9933333,-1.6707901) (2.0000000,-1.6666667) (2.0066667,-1.6625749) (2.0133333,-1.6585144) (2.0200000,-1.6544850) (2.0266667,-1.6504862) (2.0333333,-1.6465178) (2.0400000,-1.6425795) (2.0466667,-1.6386710) (2.0533333,-1.6347919) (2.0600000,-1.6309421) (2.0666667,-1.6271211) (2.0733333,-1.6233288) (2.0800000,-1.6195648) (2.0866667,-1.6158289) (2.0933333,-1.6121208) (2.1000000,-1.6084402) (2.1066667,-1.6047868) (2.1133333,-1.6011605) (2.1200000,-1.5975609) (2.1266667,-1.5939877) (2.1333333,-1.5904408) (2.1400000,-1.5869198) (2.1466667,-1.5834246) (2.1533333,-1.5799548) (2.1600000,-1.5765103) (2.1666667,-1.5730908) (2.1733333,-1.5696961) (2.1800000,-1.5663259) (2.1866667,-1.5629800) (2.1933333,-1.5596582) (2.2000000,-1.5563602) (2.2066667,-1.5530859) (2.2133333,-1.5498350) (2.2200000,-1.5466072) (2.2266667,-1.5434025) (2.2333333,-1.5402206) (2.2400000,-1.5370612) (2.2466667,-1.5339241) (2.2533333,-1.5308093) (2.2600000,-1.5277163) (2.2666667,-1.5246452) (2.2733333,-1.5215956) (2.2800000,-1.5185674) (2.2866667,-1.5155603) (2.2933333,-1.5125742) (2.3000000,-1.5096090) (2.3066667,-1.5066643) (2.3133333,-1.5037401) (2.3200000,-1.5008361) (2.3266667,-1.4979522) (2.3333333,-1.4950882) (2.3400000,-1.4922439) (2.3466667,-1.4894192) (2.3533333,-1.4866138) (2.3600000,-1.4838276) (2.3666667,-1.4810605) (2.3733333,-1.4783122) (2.3800000,-1.4755827) (2.3866667,-1.4728717) (2.3933333,-1.4701791) (2.4000000,-1.4675047) (2.4066667,-1.4648484) (2.4133333,-1.4622100) (2.4200000,-1.4595894) (2.4266667,-1.4569864) (2.4333333,-1.4544009) (2.4400000,-1.4518327) (2.4466667,-1.4492816) (2.4533333,-1.4467476) (2.4600000,-1.4442305) (2.4666667,-1.4417301) (2.4733333,-1.4392463) (2.4800000,-1.4367790) (2.4866667,-1.4343280) (2.4933333,-1.4318932) (2.5000000,-1.4294745) (2.5066667,-1.4270716) (2.5133333,-1.4246846) (2.5200000,-1.4223133) (2.5266667,-1.4199574) (2.5333333,-1.4176170) (2.5400000,-1.4152919) (2.5466667,-1.4129819) (2.5533333,-1.4106869) (2.5600000,-1.4084069) (2.5666667,-1.4061417) (2.5733333,-1.4038911) (2.5800000,-1.4016551) (2.5866667,-1.3994335) (2.5933333,-1.3972262) (2.6000000,-1.3950331) (2.6066667,-1.3928542) (2.6133333,-1.3906892) (2.6200000,-1.3885381) (2.6266667,-1.3864007) (2.6333333,-1.3842770) (2.6400000,-1.3821668) (2.6466667,-1.3800701) (2.6533333,-1.3779867) (2.6600000,-1.3759165) (2.6666667,-1.3738594) (2.6733333,-1.3718154) (2.6800000,-1.3697843) (2.6866667,-1.3677659) (2.6933333,-1.3657604) (2.7000000,-1.3637674) (2.7066667,-1.3617869) (2.7133333,-1.3598189) (2.7200000,-1.3578632) (2.7266667,-1.3559197) (2.7333333,-1.3539884) (2.7400000,-1.3520691) (2.7466667,-1.3501617) (2.7533333,-1.3482663) (2.7600000,-1.3463826) (2.7666667,-1.3445106) (2.7733333,-1.3426501) (2.7800000,-1.3408012) (2.7866667,-1.3389637) (2.7933333,-1.3371376) (2.8000000,-1.3353227) (2.8066667,-1.3335189) (2.8133333,-1.3317263) (2.8200000,-1.3299446) (2.8266667,-1.3281738) (2.8333333,-1.3264139) (2.8400000,-1.3246648) (2.8466667,-1.3229263) (2.8533333,-1.3211984) (2.8600000,-1.3194810) (2.8666667,-1.3177741) (2.8733333,-1.3160775) (2.8800000,-1.3143912) (2.8866667,-1.3127151) (2.8933333,-1.3110492) (2.9000000,-1.3093933) (2.9066667,-1.3077474) (2.9133333,-1.3061114) (2.9200000,-1.3044853) (2.9266667,-1.3028690) (2.9333333,-1.3012623) (2.9400000,-1.2996653) (2.9466667,-1.2980778) (2.9533333,-1.2964999) (2.9600000,-1.2949313) (2.9666667,-1.2933721) (2.9733333,-1.2918223) (2.9800000,-1.2902816) (2.9866667,-1.2887501) (2.9933333,-1.2872277) (3.0000000,-1.2857143) (3.0066667,-1.2842099) (3.0133333,-1.2827143) (3.0200000,-1.2812277) (3.0266667,-1.2797498) (3.0333333,-1.2782806) (3.0400000,-1.2768200) (3.0466667,-1.2753681) (3.0533333,-1.2739247) (3.0600000,-1.2724897) (3.0666667,-1.2710632) (3.0733333,-1.2696450) (3.0800000,-1.2682352) (3.0866667,-1.2668335) (3.0933333,-1.2654401) (3.1000000,-1.2640547) (3.1066667,-1.2626775) (3.1133333,-1.2613082) (3.1200000,-1.2599469) (3.1266667,-1.2585935) (3.1333333,-1.2572480) (3.1400000,-1.2559102) (3.1466667,-1.2545802) (3.1533333,-1.2532579) (3.1600000,-1.2519432) (3.1666667,-1.2506361) (3.1733333,-1.2493365) (3.1800000,-1.2480444) (3.1866667,-1.2467598) (3.1933333,-1.2454825) (3.2000000,-1.2442126) (3.2066667,-1.2429499) (3.2133333,-1.2416945) (3.2200000,-1.2404462) (3.2266667,-1.2392051) (3.2333333,-1.2379711) (3.2400000,-1.2367441) (3.2466667,-1.2355241) (3.2533333,-1.2343110) (3.2600000,-1.2331049) (3.2666667,-1.2319056) (3.2733333,-1.2307131) (3.2800000,-1.2295273) (3.2866667,-1.2283483) (3.2933333,-1.2271760) (3.3000000,-1.2260103) (3.3066667,-1.2248511) (3.3133333,-1.2236986) (3.3200000,-1.2225525) (3.3266667,-1.2214129) (3.3333333,-1.2202796) (3.3400000,-1.2191528) (3.3466667,-1.2180323) (3.3533333,-1.2169181) (3.3600000,-1.2158101) (3.3666667,-1.2147083) (3.3733333,-1.2136127) (3.3800000,-1.2125233) (3.3866667,-1.2114399) (3.3933333,-1.2103625) (3.4000000,-1.2092912) (3.4066667,-1.2082258) (3.4133333,-1.2071664) (3.4200000,-1.2061129) (3.4266667,-1.2050652) (3.4333333,-1.2040233) (3.4400000,-1.2029873) (3.4466667,-1.2019569) (3.4533333,-1.2009323) (3.4600000,-1.1999133) (3.4666667,-1.1989000) (3.4733333,-1.1978923) (3.4800000,-1.1968901) (3.4866667,-1.1958935) (3.4933333,-1.1949024) (3.5000000,-1.1939167) (3.5066667,-1.1929364) (3.5133333,-1.1919615) (3.5200000,-1.1909920) (3.5266667,-1.1900278) (3.5333333,-1.1890689) (3.5400000,-1.1881153) (3.5466667,-1.1871668) (3.5533333,-1.1862236) (3.5600000,-1.1852855) (3.5666667,-1.1843525) (3.5733333,-1.1834246) (3.5800000,-1.1825018) (3.5866667,-1.1815841) (3.5933333,-1.1806713) (3.6000000,-1.1797634) (3.6066667,-1.1788606) (3.6133333,-1.1779626) (3.6200000,-1.1770695) (3.6266667,-1.1761812) (3.6333333,-1.1752978) (3.6400000,-1.1744191) (3.6466667,-1.1735452) (3.6533333,-1.1726760) (3.6600000,-1.1718115) (3.6666667,-1.1709517) (3.6733333,-1.1700966) (3.6800000,-1.1692460) (3.6866667,-1.1684000) (3.6933333,-1.1675586) (3.7000000,-1.1667217) (3.7066667,-1.1658893) (3.7133333,-1.1650614) (3.7200000,-1.1642379) (3.7266667,-1.1634189) (3.7333333,-1.1626042) (3.7400000,-1.1617939) (3.7466667,-1.1609879) (3.7533333,-1.1601863) (3.7600000,-1.1593889) (3.7666667,-1.1585958) (3.7733333,-1.1578069) (3.7800000,-1.1570223) (3.7866667,-1.1562418) (3.7933333,-1.1554655) (3.8000000,-1.1546933) (3.8066667,-1.1539252) (3.8133333,-1.1531612) (3.8200000,-1.1524013) (3.8266667,-1.1516454) (3.8333333,-1.1508935) (3.8400000,-1.1501456) (3.8466667,-1.1494017) (3.8533333,-1.1486617) (3.8600000,-1.1479256) (3.8666667,-1.1471935) (3.8733333,-1.1464652) (3.8800000,-1.1457407) (3.8866667,-1.1450201) (3.8933333,-1.1443032) (3.9000000,-1.1435902) (3.9066667,-1.1428809) (3.9133333,-1.1421753) (3.9200000,-1.1414735) (3.9266667,-1.1407754) (3.9333333,-1.1400809) (3.9400000,-1.1393901) (3.9466667,-1.1387029) (3.9533333,-1.1380193) (3.9600000,-1.1373393) (3.9666667,-1.1366628) (3.9733333,-1.1359899) (3.9800000,-1.1353205) (3.9866667,-1.1346547) (3.9933333,-1.1339923) (4.0000000,-1.1333333)};
\draw[violet!75!black,densely dotted] (-1,-3.0000000)--(1,3.0000000);
\fill[violet!75!black] (-1,-3.0000000) circle(2.3pt);
\fill[violet!75!black] (1,3.0000000) circle(2.3pt);
\end{scope}
\end{scope}
\node[anchor=west,align=left,text width=4.35cm] at (8.9,4.4) {蓝实线：$a=2$\\[8pt]橙虚线：$a=\frac12$\\[14pt]左右分支关于原点对称。\\[8pt]$x=0$ 为竖直渐近线；\\图像不经过原点。};
\node[text width=12.8cm,align=center] at (7.1,.65) {$D=\mathbb R\setminus\{0\}$；$x=0$ 无定义，仍然是奇函数。};
\end{tikzpicture}\end{center}


令 $t=a^x>0$，把 $x$ 换成 $-x$ 后 $t$ 变为 $1/t$。于是
\[
\begin{aligned}
\frac{a^{-x}-1}{a^{-x}+1}&=\frac{1-t}{1+t}=-\frac{t-1}{t+1},\\[3pt]
\frac{a^{-x}+1}{a^{-x}-1}&=\frac{1+t}{1-t}=-\frac{t+1}{t-1}.
\end{aligned}
\]
奇2的定义域为 $\mathbb R$；奇3须排除 $t=1$，即 $x=0$。两式定义域都对称，且都反号，所以均为奇函数。
图中奇3的左右分支配对，直接说明“奇函数不一定经过原点”。
\end{proofbox}
\clearpage
\begin{proofbox}
\textbf{\textcolor{CProof}{证明二（续）：真数变倒数——两种对数分式}}
% Formula-based samples; see F083 checks/build_tikz_diagrams.py.
\begin{center}\begin{tikzpicture}[font=\small,>=stealth,line join=round]
\draw[gray!40,rounded corners=3pt] (0,0) rectangle (14.2,8.4);
\node at (7.1,7.75) {奇4：$y=\ln\dfrac{kx-b}{kx+b}$};
\begin{scope}[shift={(4.400000,3.900000)},x=0.850000cm,y=0.562500cm]
\fill[gray!12] (-1,-4) rectangle (1,4);
\draw[gray,dashed] (-1,-4)--(-1,4);
\draw[gray,dashed] (1,-4)--(1,4);
\draw[->,gray] (-4,0)--(4,0) node[right] {$x$};
\draw[->,gray] (0,-4)--(0,4) node[above] {$y$};
\node[anchor=north] at (-1,-4-.07/0.562500) {$-1$};
\node[anchor=north] at (1,-4-.07/0.562500) {$1$};
\draw[gray] (-2,-.035/0.562500)--(-2,.035/0.562500);
\node[anchor=north] at (-2,-.07/0.562500) {$-2$};
\draw[gray] (2,-.035/0.562500)--(2,.035/0.562500);
\node[anchor=north] at (2,-.07/0.562500) {$2$};
\begin{scope}
\clip (-4,-4) rectangle (4,4);
\draw[blue,thick,solid] plot coordinates {(-4.0000000,0.5108256) (-3.9950000,0.5114932) (-3.9900000,0.5121625) (-3.9850000,0.5128337) (-3.9800000,0.5135066) (-3.9750000,0.5141813) (-3.9700000,0.5148579) (-3.9650000,0.5155363) (-3.9600000,0.5162165) (-3.9550000,0.5168985) (-3.9500000,0.5175824) (-3.9450000,0.5182681) (-3.9400000,0.5189557) (-3.9350000,0.5196452) (-3.9300000,0.5203366) (-3.9250000,0.5210298) (-3.9200000,0.5217249) (-3.9150000,0.5224219) (-3.9100000,0.5231209) (-3.9050000,0.5238217) (-3.9000000,0.5245245) (-3.8950000,0.5252292) (-3.8900000,0.5259358) (-3.8850000,0.5266444) (-3.8800000,0.5273549) (-3.8750000,0.5280674) (-3.8700000,0.5287819) (-3.8650000,0.5294984) (-3.8600000,0.5302168) (-3.8550000,0.5309373) (-3.8500000,0.5316597) (-3.8450000,0.5323842) (-3.8400000,0.5331107) (-3.8350000,0.5338392) (-3.8300000,0.5345698) (-3.8250000,0.5353024) (-3.8200000,0.5360370) (-3.8150000,0.5367738) (-3.8100000,0.5375126) (-3.8050000,0.5382535) (-3.8000000,0.5389965) (-3.7950000,0.5397416) (-3.7900000,0.5404888) (-3.7850000,0.5412382) (-3.7800000,0.5419896) (-3.7750000,0.5427432) (-3.7700000,0.5434990) (-3.7650000,0.5442569) (-3.7600000,0.5450170) (-3.7550000,0.5457793) (-3.7500000,0.5465437) (-3.7450000,0.5473104) (-3.7400000,0.5480792) (-3.7350000,0.5488503) (-3.7300000,0.5496236) (-3.7250000,0.5503991) (-3.7200000,0.5511769) (-3.7150000,0.5519570) (-3.7100000,0.5527393) (-3.7050000,0.5535239) (-3.7000000,0.5543107) (-3.6950000,0.5550999) (-3.6900000,0.5558914) (-3.6850000,0.5566852) (-3.6800000,0.5574813) (-3.6750000,0.5582798) (-3.6700000,0.5590806) (-3.6650000,0.5598838) (-3.6600000,0.5606893) (-3.6550000,0.5614973) (-3.6500000,0.5623076) (-3.6450000,0.5631203) (-3.6400000,0.5639354) (-3.6350000,0.5647530) (-3.6300000,0.5655730) (-3.6250000,0.5663955) (-3.6200000,0.5672204) (-3.6150000,0.5680478) (-3.6100000,0.5688776) (-3.6050000,0.5697100) (-3.6000000,0.5705449) (-3.5950000,0.5713822) (-3.5900000,0.5722221) (-3.5850000,0.5730646) (-3.5800000,0.5739096) (-3.5750000,0.5747572) (-3.5700000,0.5756073) (-3.5650000,0.5764600) (-3.5600000,0.5773154) (-3.5550000,0.5781733) (-3.5500000,0.5790339) (-3.5450000,0.5798971) (-3.5400000,0.5807629) (-3.5350000,0.5816314) (-3.5300000,0.5825026) (-3.5250000,0.5833765) (-3.5200000,0.5842531) (-3.5150000,0.5851324) (-3.5100000,0.5860144) (-3.5050000,0.5868992) (-3.5000000,0.5877867) (-3.4950000,0.5886769) (-3.4900000,0.5895700) (-3.4850000,0.5904658) (-3.4800000,0.5913645) (-3.4750000,0.5922660) (-3.4700000,0.5931703) (-3.4650000,0.5940774) (-3.4600000,0.5949874) (-3.4550000,0.5959003) (-3.4500000,0.5968161) (-3.4450000,0.5977347) (-3.4400000,0.5986563) (-3.4350000,0.5995809) (-3.4300000,0.6005083) (-3.4250000,0.6014388) (-3.4200000,0.6023722) (-3.4150000,0.6033085) (-3.4100000,0.6042479) (-3.4050000,0.6051904) (-3.4000000,0.6061358) (-3.3950000,0.6070843) (-3.3900000,0.6080359) (-3.3850000,0.6089905) (-3.3800000,0.6099482) (-3.3750000,0.6109091) (-3.3700000,0.6118731) (-3.3650000,0.6128402) (-3.3600000,0.6138104) (-3.3550000,0.6147839) (-3.3500000,0.6157605) (-3.3450000,0.6167404) (-3.3400000,0.6177234) (-3.3350000,0.6187097) (-3.3300000,0.6196993) (-3.3250000,0.6206921) (-3.3200000,0.6216882) (-3.3150000,0.6226876) (-3.3100000,0.6236904) (-3.3050000,0.6246965) (-3.3000000,0.6257059) (-3.2950000,0.6267187) (-3.2900000,0.6277349) (-3.2850000,0.6287545) (-3.2800000,0.6297776) (-3.2750000,0.6308040) (-3.2700000,0.6318340) (-3.2650000,0.6328674) (-3.2600000,0.6339043) (-3.2550000,0.6349448) (-3.2500000,0.6359888) (-3.2450000,0.6370363) (-3.2400000,0.6380874) (-3.2350000,0.6391421) (-3.2300000,0.6402004) (-3.2250000,0.6412623) (-3.2200000,0.6423279) (-3.2150000,0.6433972) (-3.2100000,0.6444701) (-3.2050000,0.6455468) (-3.2000000,0.6466272) (-3.1950000,0.6477113) (-3.1900000,0.6487992) (-3.1850000,0.6498909) (-3.1800000,0.6509864) (-3.1750000,0.6520857) (-3.1700000,0.6531889) (-3.1650000,0.6542959) (-3.1600000,0.6554069) (-3.1550000,0.6565217) (-3.1500000,0.6576405) (-3.1450000,0.6587632) (-3.1400000,0.6598900) (-3.1350000,0.6610207) (-3.1300000,0.6621554) (-3.1250000,0.6632942) (-3.1200000,0.6644371) (-3.1150000,0.6655840) (-3.1100000,0.6667351) (-3.1050000,0.6678903) (-3.1000000,0.6690496) (-3.0950000,0.6702132) (-3.0900000,0.6713809) (-3.0850000,0.6725529) (-3.0800000,0.6737291) (-3.0750000,0.6749096) (-3.0700000,0.6760944) (-3.0650000,0.6772835) (-3.0600000,0.6784770) (-3.0550000,0.6796748) (-3.0500000,0.6808771) (-3.0450000,0.6820838) (-3.0400000,0.6832949) (-3.0350000,0.6845105) (-3.0300000,0.6857306) (-3.0250000,0.6869552) (-3.0200000,0.6881844) (-3.0150000,0.6894182) (-3.0100000,0.6906565) (-3.0050000,0.6918995) (-3.0000000,0.6931472) (-2.9950000,0.6943995) (-2.9900000,0.6956566) (-2.9850000,0.6969184) (-2.9800000,0.6981850) (-2.9750000,0.6994563) (-2.9700000,0.7007326) (-2.9650000,0.7020136) (-2.9600000,0.7032996) (-2.9550000,0.7045904) (-2.9500000,0.7058862) (-2.9450000,0.7071870) (-2.9400000,0.7084928) (-2.9350000,0.7098036) (-2.9300000,0.7111194) (-2.9250000,0.7124404) (-2.9200000,0.7137665) (-2.9150000,0.7150977) (-2.9100000,0.7164341) (-2.9050000,0.7177758) (-2.9000000,0.7191227) (-2.8950000,0.7204748) (-2.8900000,0.7218323) (-2.8850000,0.7231952) (-2.8800000,0.7245634) (-2.8750000,0.7259370) (-2.8700000,0.7273161) (-2.8650000,0.7287006) (-2.8600000,0.7300907) (-2.8550000,0.7314863) (-2.8500000,0.7328875) (-2.8450000,0.7342943) (-2.8400000,0.7357068) (-2.8350000,0.7371250) (-2.8300000,0.7385488) (-2.8250000,0.7399785) (-2.8200000,0.7414139) (-2.8150000,0.7428552) (-2.8100000,0.7443023) (-2.8050000,0.7457554) (-2.8000000,0.7472144) (-2.7950000,0.7486794) (-2.7900000,0.7501504) (-2.7850000,0.7516275) (-2.7800000,0.7531106) (-2.7750000,0.7546000) (-2.7700000,0.7560955) (-2.7650000,0.7575972) (-2.7600000,0.7591051) (-2.7550000,0.7606194) (-2.7500000,0.7621401) (-2.7450000,0.7636671) (-2.7400000,0.7652005) (-2.7350000,0.7667404) (-2.7300000,0.7682868) (-2.7250000,0.7698398) (-2.7200000,0.7713994) (-2.7150000,0.7729656) (-2.7100000,0.7745385) (-2.7050000,0.7761181) (-2.7000000,0.7777046) (-2.6950000,0.7792978) (-2.6900000,0.7808979) (-2.6850000,0.7825050) (-2.6800000,0.7841190) (-2.6750000,0.7857400) (-2.6700000,0.7873680) (-2.6650000,0.7890032) (-2.6600000,0.7906455) (-2.6550000,0.7922951) (-2.6500000,0.7939519) (-2.6450000,0.7956160) (-2.6400000,0.7972874) (-2.6350000,0.7989663) (-2.6300000,0.8006526) (-2.6250000,0.8023465) (-2.6200000,0.8040479) (-2.6150000,0.8057569) (-2.6100000,0.8074736) (-2.6050000,0.8091980) (-2.6000000,0.8109302) (-2.5950000,0.8126703) (-2.5900000,0.8144182) (-2.5850000,0.8161741) (-2.5800000,0.8179380) (-2.5750000,0.8197099) (-2.5700000,0.8214900) (-2.5650000,0.8232782) (-2.5600000,0.8250747) (-2.5550000,0.8268795) (-2.5500000,0.8286927) (-2.5450000,0.8305142) (-2.5400000,0.8323443) (-2.5350000,0.8341829) (-2.5300000,0.8360301) (-2.5250000,0.8378860) (-2.5200000,0.8397507) (-2.5150000,0.8416241) (-2.5100000,0.8435064) (-2.5050000,0.8453976) (-2.5000000,0.8472979) (-2.4950000,0.8492072) (-2.4900000,0.8511256) (-2.4850000,0.8530533) (-2.4800000,0.8549902) (-2.4750000,0.8569365) (-2.4700000,0.8588922) (-2.4650000,0.8608574) (-2.4600000,0.8628322) (-2.4550000,0.8648166) (-2.4500000,0.8668107) (-2.4450000,0.8688146) (-2.4400000,0.8708284) (-2.4350000,0.8728521) (-2.4300000,0.8748858) (-2.4250000,0.8769297) (-2.4200000,0.8789837) (-2.4150000,0.8810480) (-2.4100000,0.8831226) (-2.4050000,0.8852076) (-2.4000000,0.8873032) (-2.3950000,0.8894093) (-2.3900000,0.8915262) (-2.3850000,0.8936538) (-2.3800000,0.8957922) (-2.3750000,0.8979416) (-2.3700000,0.9001020) (-2.3650000,0.9022735) (-2.3600000,0.9044563) (-2.3550000,0.9066503) (-2.3500000,0.9088558) (-2.3450000,0.9110727) (-2.3400000,0.9133012) (-2.3350000,0.9155414) (-2.3300000,0.9177934) (-2.3250000,0.9200572) (-2.3200000,0.9223330) (-2.3150000,0.9246210) (-2.3100000,0.9269211) (-2.3050000,0.9292334) (-2.3000000,0.9315582) (-2.2950000,0.9338955) (-2.2900000,0.9362453) (-2.2850000,0.9386079) (-2.2800000,0.9409833) (-2.2750000,0.9433717) (-2.2700000,0.9457731) (-2.2650000,0.9481876) (-2.2600000,0.9506155) (-2.2550000,0.9530567) (-2.2500000,0.9555114) (-2.2450000,0.9579798) (-2.2400000,0.9604620) (-2.2350000,0.9629580) (-2.2300000,0.9654680) (-2.2250000,0.9679921) (-2.2200000,0.9705305) (-2.2150000,0.9730833) (-2.2100000,0.9756506) (-2.2050000,0.9782325) (-2.2000000,0.9808293) (-2.1950000,0.9834409) (-2.1900000,0.9860676) (-2.1850000,0.9887095) (-2.1800000,0.9913668) (-2.1750000,0.9940395) (-2.1700000,0.9967278) (-2.1650000,0.9994320) (-2.1600000,1.0021520) (-2.1550000,1.0048882) (-2.1500000,1.0076405) (-2.1450000,1.0104093) (-2.1400000,1.0131945) (-2.1350000,1.0159965) (-2.1300000,1.0188154) (-2.1250000,1.0216512) (-2.1200000,1.0245043) (-2.1150000,1.0273747) (-2.1100000,1.0302627) (-2.1050000,1.0331684) (-2.1000000,1.0360919) (-2.0950000,1.0390335) (-2.0900000,1.0419934) (-2.0850000,1.0449717) (-2.0800000,1.0479686) (-2.0750000,1.0509842) (-2.0700000,1.0540189) (-2.0650000,1.0570728) (-2.0600000,1.0601460) (-2.0550000,1.0632388) (-2.0500000,1.0663514) (-2.0450000,1.0694840) (-2.0400000,1.0726368) (-2.0350000,1.0758100) (-2.0300000,1.0790038) (-2.0250000,1.0822185) (-2.0200000,1.0854542) (-2.0150000,1.0887112) (-2.0100000,1.0919897) (-2.0050000,1.0952900) (-2.0000000,1.0986123) (-1.9950000,1.1019568) (-1.9900000,1.1053237) (-1.9850000,1.1087134) (-1.9800000,1.1121260) (-1.9750000,1.1155618) (-1.9700000,1.1190212) (-1.9650000,1.1225042) (-1.9600000,1.1260113) (-1.9550000,1.1295426) (-1.9500000,1.1330985) (-1.9450000,1.1366792) (-1.9400000,1.1402850) (-1.9350000,1.1439162) (-1.9300000,1.1475731) (-1.9250000,1.1512560) (-1.9200000,1.1549652) (-1.9150000,1.1587010) (-1.9100000,1.1624638) (-1.9050000,1.1662537) (-1.9000000,1.1700713) (-1.8950000,1.1739167) (-1.8900000,1.1777903) (-1.8850000,1.1816925) (-1.8800000,1.1856237) (-1.8750000,1.1895841) (-1.8700000,1.1935741) (-1.8650000,1.1975941) (-1.8600000,1.2016445) (-1.8550000,1.2057257) (-1.8500000,1.2098379) (-1.8450000,1.2139817) (-1.8400000,1.2181574) (-1.8350000,1.2223655) (-1.8300000,1.2266063) (-1.8250000,1.2308803) (-1.8200000,1.2351878) (-1.8150000,1.2395294) (-1.8100000,1.2439055) (-1.8050000,1.2483165) (-1.8000000,1.2527630) (-1.7950000,1.2572453) (-1.7900000,1.2617639) (-1.7850000,1.2663194) (-1.7800000,1.2709123) (-1.7750000,1.2755430) (-1.7700000,1.2802121) (-1.7650000,1.2849201) (-1.7600000,1.2896675) (-1.7550000,1.2944550) (-1.7500000,1.2992830) (-1.7450000,1.3041521) (-1.7400000,1.3090630) (-1.7350000,1.3140162) (-1.7300000,1.3190124) (-1.7250000,1.3240521) (-1.7200000,1.3291359) (-1.7150000,1.3342647) (-1.7100000,1.3394389) (-1.7050000,1.3446594) (-1.7000000,1.3499267) (-1.6950000,1.3552416) (-1.6900000,1.3606049) (-1.6850000,1.3660172) (-1.6800000,1.3714793) (-1.6750000,1.3769920) (-1.6700000,1.3825560) (-1.6650000,1.3881723) (-1.6600000,1.3938416) (-1.6550000,1.3995647) (-1.6500000,1.4053426) (-1.6450000,1.4111760) (-1.6400000,1.4170660) (-1.6350000,1.4230135) (-1.6300000,1.4290193) (-1.6250000,1.4350845) (-1.6200000,1.4412101) (-1.6150000,1.4473971) (-1.6100000,1.4536465) (-1.6050000,1.4599595) (-1.6000000,1.4663371) (-1.5950000,1.4727804) (-1.5900000,1.4792906) (-1.5850000,1.4858689) (-1.5800000,1.4925166) (-1.5750000,1.4992348) (-1.5700000,1.5060248) (-1.5650000,1.5128880) (-1.5600000,1.5198258) (-1.5550000,1.5268394) (-1.5500000,1.5339304) (-1.5450000,1.5411001) (-1.5400000,1.5483502) (-1.5350000,1.5556822) (-1.5300000,1.5630976) (-1.5250000,1.5705981) (-1.5200000,1.5781854) (-1.5150000,1.5858612) (-1.5100000,1.5936273) (-1.5050000,1.6014856) (-1.5000000,1.6094379) (-1.4950000,1.6174862) (-1.4900000,1.6256326) (-1.4850000,1.6338790) (-1.4800000,1.6422277) (-1.4750000,1.6506809) (-1.4700000,1.6592407) (-1.4650000,1.6679097) (-1.4600000,1.6766901) (-1.4550000,1.6855846) (-1.4500000,1.6945957) (-1.4450000,1.7037261) (-1.4400000,1.7129786) (-1.4350000,1.7223560) (-1.4300000,1.7318613) (-1.4250000,1.7414976) (-1.4200000,1.7512681) (-1.4150000,1.7611760) (-1.4100000,1.7712249) (-1.4050000,1.7814181) (-1.4000000,1.7917595) (-1.3950000,1.8022527) (-1.3900000,1.8129019) (-1.3850000,1.8237111) (-1.3800000,1.8346845) (-1.3750000,1.8458267) (-1.3700000,1.8571422) (-1.3650000,1.8686359) (-1.3600000,1.8803129) (-1.3550000,1.8921782) (-1.3500000,1.9042375) (-1.3450000,1.9164963) (-1.3400000,1.9289606) (-1.3350000,1.9416366) (-1.3300000,1.9545309) (-1.3250000,1.9676501) (-1.3200000,1.9810015) (-1.3150000,1.9945923) (-1.3100000,2.0084305) (-1.3050000,2.0225242) (-1.3000000,2.0368819) (-1.2950000,2.0515128) (-1.2900000,2.0664262) (-1.2850000,2.0816321) (-1.2800000,2.0971411) (-1.2750000,2.1129642) (-1.2700000,2.1291132) (-1.2650000,2.1456002) (-1.2600000,2.1624385) (-1.2550000,2.1796417) (-1.2500000,2.1972246) (-1.2450000,2.2152026) (-1.2400000,2.2335922) (-1.2350000,2.2524110) (-1.2300000,2.2716776) (-1.2250000,2.2914118) (-1.2200000,2.3116349) (-1.2150000,2.3323697) (-1.2100000,2.3536403) (-1.2050000,2.3754728) (-1.2000000,2.3978953) (-1.1950000,2.4209378) (-1.1900000,2.4446328) (-1.1850000,2.4690153) (-1.1800000,2.4941233) (-1.1750000,2.5199980) (-1.1700000,2.5466840) (-1.1650000,2.5742302) (-1.1600000,2.6026897) (-1.1550000,2.6321209) (-1.1500000,2.6625878) (-1.1450000,2.6941611) (-1.1400000,2.7269187) (-1.1350000,2.7609471) (-1.1300000,2.7963428) (-1.1250000,2.8332133) (-1.1200000,2.8716796) (-1.1150000,2.9118780) (-1.1100000,2.9539629) (-1.1050000,2.9981104) (-1.1000000,3.0445224) (-1.0950000,3.0934319) (-1.0900000,3.1451097) (-1.0850000,3.1998729) (-1.0800000,3.2580965) (-1.0750000,3.3202283) (-1.0700000,3.3868086) (-1.0650000,3.4584982) (-1.0600000,3.5361167) (-1.0550000,3.6206979) (-1.0500000,3.7135721) (-1.0450000,3.8164906) (-1.0400000,3.9318256)};
\draw[blue,thick,solid] plot coordinates {(1.0400000,-3.9318256) (1.0450000,-3.8164906) (1.0500000,-3.7135721) (1.0550000,-3.6206979) (1.0600000,-3.5361167) (1.0650000,-3.4584982) (1.0700000,-3.3868086) (1.0750000,-3.3202283) (1.0800000,-3.2580965) (1.0850000,-3.1998729) (1.0900000,-3.1451097) (1.0950000,-3.0934319) (1.1000000,-3.0445224) (1.1050000,-2.9981104) (1.1100000,-2.9539629) (1.1150000,-2.9118780) (1.1200000,-2.8716796) (1.1250000,-2.8332133) (1.1300000,-2.7963428) (1.1350000,-2.7609471) (1.1400000,-2.7269187) (1.1450000,-2.6941611) (1.1500000,-2.6625878) (1.1550000,-2.6321209) (1.1600000,-2.6026897) (1.1650000,-2.5742302) (1.1700000,-2.5466840) (1.1750000,-2.5199980) (1.1800000,-2.4941233) (1.1850000,-2.4690153) (1.1900000,-2.4446328) (1.1950000,-2.4209378) (1.2000000,-2.3978953) (1.2050000,-2.3754728) (1.2100000,-2.3536403) (1.2150000,-2.3323697) (1.2200000,-2.3116349) (1.2250000,-2.2914118) (1.2300000,-2.2716776) (1.2350000,-2.2524110) (1.2400000,-2.2335922) (1.2450000,-2.2152026) (1.2500000,-2.1972246) (1.2550000,-2.1796417) (1.2600000,-2.1624385) (1.2650000,-2.1456002) (1.2700000,-2.1291132) (1.2750000,-2.1129642) (1.2800000,-2.0971411) (1.2850000,-2.0816321) (1.2900000,-2.0664262) (1.2950000,-2.0515128) (1.3000000,-2.0368819) (1.3050000,-2.0225242) (1.3100000,-2.0084305) (1.3150000,-1.9945923) (1.3200000,-1.9810015) (1.3250000,-1.9676501) (1.3300000,-1.9545309) (1.3350000,-1.9416366) (1.3400000,-1.9289606) (1.3450000,-1.9164963) (1.3500000,-1.9042375) (1.3550000,-1.8921782) (1.3600000,-1.8803129) (1.3650000,-1.8686359) (1.3700000,-1.8571422) (1.3750000,-1.8458267) (1.3800000,-1.8346845) (1.3850000,-1.8237111) (1.3900000,-1.8129019) (1.3950000,-1.8022527) (1.4000000,-1.7917595) (1.4050000,-1.7814181) (1.4100000,-1.7712249) (1.4150000,-1.7611760) (1.4200000,-1.7512681) (1.4250000,-1.7414976) (1.4300000,-1.7318613) (1.4350000,-1.7223560) (1.4400000,-1.7129786) (1.4450000,-1.7037261) (1.4500000,-1.6945957) (1.4550000,-1.6855846) (1.4600000,-1.6766901) (1.4650000,-1.6679097) (1.4700000,-1.6592407) (1.4750000,-1.6506809) (1.4800000,-1.6422277) (1.4850000,-1.6338790) (1.4900000,-1.6256326) (1.4950000,-1.6174862) (1.5000000,-1.6094379) (1.5050000,-1.6014856) (1.5100000,-1.5936273) (1.5150000,-1.5858612) (1.5200000,-1.5781854) (1.5250000,-1.5705981) (1.5300000,-1.5630976) (1.5350000,-1.5556822) (1.5400000,-1.5483502) (1.5450000,-1.5411001) (1.5500000,-1.5339304) (1.5550000,-1.5268394) (1.5600000,-1.5198258) (1.5650000,-1.5128880) (1.5700000,-1.5060248) (1.5750000,-1.4992348) (1.5800000,-1.4925166) (1.5850000,-1.4858689) (1.5900000,-1.4792906) (1.5950000,-1.4727804) (1.6000000,-1.4663371) (1.6050000,-1.4599595) (1.6100000,-1.4536465) (1.6150000,-1.4473971) (1.6200000,-1.4412101) (1.6250000,-1.4350845) (1.6300000,-1.4290193) (1.6350000,-1.4230135) (1.6400000,-1.4170660) (1.6450000,-1.4111760) (1.6500000,-1.4053426) (1.6550000,-1.3995647) (1.6600000,-1.3938416) (1.6650000,-1.3881723) (1.6700000,-1.3825560) (1.6750000,-1.3769920) (1.6800000,-1.3714793) (1.6850000,-1.3660172) (1.6900000,-1.3606049) (1.6950000,-1.3552416) (1.7000000,-1.3499267) (1.7050000,-1.3446594) (1.7100000,-1.3394389) (1.7150000,-1.3342647) (1.7200000,-1.3291359) (1.7250000,-1.3240521) (1.7300000,-1.3190124) (1.7350000,-1.3140162) (1.7400000,-1.3090630) (1.7450000,-1.3041521) (1.7500000,-1.2992830) (1.7550000,-1.2944550) (1.7600000,-1.2896675) (1.7650000,-1.2849201) (1.7700000,-1.2802121) (1.7750000,-1.2755430) (1.7800000,-1.2709123) (1.7850000,-1.2663194) (1.7900000,-1.2617639) (1.7950000,-1.2572453) (1.8000000,-1.2527630) (1.8050000,-1.2483165) (1.8100000,-1.2439055) (1.8150000,-1.2395294) (1.8200000,-1.2351878) (1.8250000,-1.2308803) (1.8300000,-1.2266063) (1.8350000,-1.2223655) (1.8400000,-1.2181574) (1.8450000,-1.2139817) (1.8500000,-1.2098379) (1.8550000,-1.2057257) (1.8600000,-1.2016445) (1.8650000,-1.1975941) (1.8700000,-1.1935741) (1.8750000,-1.1895841) (1.8800000,-1.1856237) (1.8850000,-1.1816925) (1.8900000,-1.1777903) (1.8950000,-1.1739167) (1.9000000,-1.1700713) (1.9050000,-1.1662537) (1.9100000,-1.1624638) (1.9150000,-1.1587010) (1.9200000,-1.1549652) (1.9250000,-1.1512560) (1.9300000,-1.1475731) (1.9350000,-1.1439162) (1.9400000,-1.1402850) (1.9450000,-1.1366792) (1.9500000,-1.1330985) (1.9550000,-1.1295426) (1.9600000,-1.1260113) (1.9650000,-1.1225042) (1.9700000,-1.1190212) (1.9750000,-1.1155618) (1.9800000,-1.1121260) (1.9850000,-1.1087134) (1.9900000,-1.1053237) (1.9950000,-1.1019568) (2.0000000,-1.0986123) (2.0050000,-1.0952900) (2.0100000,-1.0919897) (2.0150000,-1.0887112) (2.0200000,-1.0854542) (2.0250000,-1.0822185) (2.0300000,-1.0790038) (2.0350000,-1.0758100) (2.0400000,-1.0726368) (2.0450000,-1.0694840) (2.0500000,-1.0663514) (2.0550000,-1.0632388) (2.0600000,-1.0601460) (2.0650000,-1.0570728) (2.0700000,-1.0540189) (2.0750000,-1.0509842) (2.0800000,-1.0479686) (2.0850000,-1.0449717) (2.0900000,-1.0419934) (2.0950000,-1.0390335) (2.1000000,-1.0360919) (2.1050000,-1.0331684) (2.1100000,-1.0302627) (2.1150000,-1.0273747) (2.1200000,-1.0245043) (2.1250000,-1.0216512) (2.1300000,-1.0188154) (2.1350000,-1.0159965) (2.1400000,-1.0131945) (2.1450000,-1.0104093) (2.1500000,-1.0076405) (2.1550000,-1.0048882) (2.1600000,-1.0021520) (2.1650000,-0.9994320) (2.1700000,-0.9967278) (2.1750000,-0.9940395) (2.1800000,-0.9913668) (2.1850000,-0.9887095) (2.1900000,-0.9860676) (2.1950000,-0.9834409) (2.2000000,-0.9808293) (2.2050000,-0.9782325) (2.2100000,-0.9756506) (2.2150000,-0.9730833) (2.2200000,-0.9705305) (2.2250000,-0.9679921) (2.2300000,-0.9654680) (2.2350000,-0.9629580) (2.2400000,-0.9604620) (2.2450000,-0.9579798) (2.2500000,-0.9555114) (2.2550000,-0.9530567) (2.2600000,-0.9506155) (2.2650000,-0.9481876) (2.2700000,-0.9457731) (2.2750000,-0.9433717) (2.2800000,-0.9409833) (2.2850000,-0.9386079) (2.2900000,-0.9362453) (2.2950000,-0.9338955) (2.3000000,-0.9315582) (2.3050000,-0.9292334) (2.3100000,-0.9269211) (2.3150000,-0.9246210) (2.3200000,-0.9223330) (2.3250000,-0.9200572) (2.3300000,-0.9177934) (2.3350000,-0.9155414) (2.3400000,-0.9133012) (2.3450000,-0.9110727) (2.3500000,-0.9088558) (2.3550000,-0.9066503) (2.3600000,-0.9044563) (2.3650000,-0.9022735) (2.3700000,-0.9001020) (2.3750000,-0.8979416) (2.3800000,-0.8957922) (2.3850000,-0.8936538) (2.3900000,-0.8915262) (2.3950000,-0.8894093) (2.4000000,-0.8873032) (2.4050000,-0.8852076) (2.4100000,-0.8831226) (2.4150000,-0.8810480) (2.4200000,-0.8789837) (2.4250000,-0.8769297) (2.4300000,-0.8748858) (2.4350000,-0.8728521) (2.4400000,-0.8708284) (2.4450000,-0.8688146) (2.4500000,-0.8668107) (2.4550000,-0.8648166) (2.4600000,-0.8628322) (2.4650000,-0.8608574) (2.4700000,-0.8588922) (2.4750000,-0.8569365) (2.4800000,-0.8549902) (2.4850000,-0.8530533) (2.4900000,-0.8511256) (2.4950000,-0.8492072) (2.5000000,-0.8472979) (2.5050000,-0.8453976) (2.5100000,-0.8435064) (2.5150000,-0.8416241) (2.5200000,-0.8397507) (2.5250000,-0.8378860) (2.5300000,-0.8360301) (2.5350000,-0.8341829) (2.5400000,-0.8323443) (2.5450000,-0.8305142) (2.5500000,-0.8286927) (2.5550000,-0.8268795) (2.5600000,-0.8250747) (2.5650000,-0.8232782) (2.5700000,-0.8214900) (2.5750000,-0.8197099) (2.5800000,-0.8179380) (2.5850000,-0.8161741) (2.5900000,-0.8144182) (2.5950000,-0.8126703) (2.6000000,-0.8109302) (2.6050000,-0.8091980) (2.6100000,-0.8074736) (2.6150000,-0.8057569) (2.6200000,-0.8040479) (2.6250000,-0.8023465) (2.6300000,-0.8006526) (2.6350000,-0.7989663) (2.6400000,-0.7972874) (2.6450000,-0.7956160) (2.6500000,-0.7939519) (2.6550000,-0.7922951) (2.6600000,-0.7906455) (2.6650000,-0.7890032) (2.6700000,-0.7873680) (2.6750000,-0.7857400) (2.6800000,-0.7841190) (2.6850000,-0.7825050) (2.6900000,-0.7808979) (2.6950000,-0.7792978) (2.7000000,-0.7777046) (2.7050000,-0.7761181) (2.7100000,-0.7745385) (2.7150000,-0.7729656) (2.7200000,-0.7713994) (2.7250000,-0.7698398) (2.7300000,-0.7682868) (2.7350000,-0.7667404) (2.7400000,-0.7652005) (2.7450000,-0.7636671) (2.7500000,-0.7621401) (2.7550000,-0.7606194) (2.7600000,-0.7591051) (2.7650000,-0.7575972) (2.7700000,-0.7560955) (2.7750000,-0.7546000) (2.7800000,-0.7531106) (2.7850000,-0.7516275) (2.7900000,-0.7501504) (2.7950000,-0.7486794) (2.8000000,-0.7472144) (2.8050000,-0.7457554) (2.8100000,-0.7443023) (2.8150000,-0.7428552) (2.8200000,-0.7414139) (2.8250000,-0.7399785) (2.8300000,-0.7385488) (2.8350000,-0.7371250) (2.8400000,-0.7357068) (2.8450000,-0.7342943) (2.8500000,-0.7328875) (2.8550000,-0.7314863) (2.8600000,-0.7300907) (2.8650000,-0.7287006) (2.8700000,-0.7273161) (2.8750000,-0.7259370) (2.8800000,-0.7245634) (2.8850000,-0.7231952) (2.8900000,-0.7218323) (2.8950000,-0.7204748) (2.9000000,-0.7191227) (2.9050000,-0.7177758) (2.9100000,-0.7164341) (2.9150000,-0.7150977) (2.9200000,-0.7137665) (2.9250000,-0.7124404) (2.9300000,-0.7111194) (2.9350000,-0.7098036) (2.9400000,-0.7084928) (2.9450000,-0.7071870) (2.9500000,-0.7058862) (2.9550000,-0.7045904) (2.9600000,-0.7032996) (2.9650000,-0.7020136) (2.9700000,-0.7007326) (2.9750000,-0.6994563) (2.9800000,-0.6981850) (2.9850000,-0.6969184) (2.9900000,-0.6956566) (2.9950000,-0.6943995) (3.0000000,-0.6931472) (3.0050000,-0.6918995) (3.0100000,-0.6906565) (3.0150000,-0.6894182) (3.0200000,-0.6881844) (3.0250000,-0.6869552) (3.0300000,-0.6857306) (3.0350000,-0.6845105) (3.0400000,-0.6832949) (3.0450000,-0.6820838) (3.0500000,-0.6808771) (3.0550000,-0.6796748) (3.0600000,-0.6784770) (3.0650000,-0.6772835) (3.0700000,-0.6760944) (3.0750000,-0.6749096) (3.0800000,-0.6737291) (3.0850000,-0.6725529) (3.0900000,-0.6713809) (3.0950000,-0.6702132) (3.1000000,-0.6690496) (3.1050000,-0.6678903) (3.1100000,-0.6667351) (3.1150000,-0.6655840) (3.1200000,-0.6644371) (3.1250000,-0.6632942) (3.1300000,-0.6621554) (3.1350000,-0.6610207) (3.1400000,-0.6598900) (3.1450000,-0.6587632) (3.1500000,-0.6576405) (3.1550000,-0.6565217) (3.1600000,-0.6554069) (3.1650000,-0.6542959) (3.1700000,-0.6531889) (3.1750000,-0.6520857) (3.1800000,-0.6509864) (3.1850000,-0.6498909) (3.1900000,-0.6487992) (3.1950000,-0.6477113) (3.2000000,-0.6466272) (3.2050000,-0.6455468) (3.2100000,-0.6444701) (3.2150000,-0.6433972) (3.2200000,-0.6423279) (3.2250000,-0.6412623) (3.2300000,-0.6402004) (3.2350000,-0.6391421) (3.2400000,-0.6380874) (3.2450000,-0.6370363) (3.2500000,-0.6359888) (3.2550000,-0.6349448) (3.2600000,-0.6339043) (3.2650000,-0.6328674) (3.2700000,-0.6318340) (3.2750000,-0.6308040) (3.2800000,-0.6297776) (3.2850000,-0.6287545) (3.2900000,-0.6277349) (3.2950000,-0.6267187) (3.3000000,-0.6257059) (3.3050000,-0.6246965) (3.3100000,-0.6236904) (3.3150000,-0.6226876) (3.3200000,-0.6216882) (3.3250000,-0.6206921) (3.3300000,-0.6196993) (3.3350000,-0.6187097) (3.3400000,-0.6177234) (3.3450000,-0.6167404) (3.3500000,-0.6157605) (3.3550000,-0.6147839) (3.3600000,-0.6138104) (3.3650000,-0.6128402) (3.3700000,-0.6118731) (3.3750000,-0.6109091) (3.3800000,-0.6099482) (3.3850000,-0.6089905) (3.3900000,-0.6080359) (3.3950000,-0.6070843) (3.4000000,-0.6061358) (3.4050000,-0.6051904) (3.4100000,-0.6042479) (3.4150000,-0.6033085) (3.4200000,-0.6023722) (3.4250000,-0.6014388) (3.4300000,-0.6005083) (3.4350000,-0.5995809) (3.4400000,-0.5986563) (3.4450000,-0.5977347) (3.4500000,-0.5968161) (3.4550000,-0.5959003) (3.4600000,-0.5949874) (3.4650000,-0.5940774) (3.4700000,-0.5931703) (3.4750000,-0.5922660) (3.4800000,-0.5913645) (3.4850000,-0.5904658) (3.4900000,-0.5895700) (3.4950000,-0.5886769) (3.5000000,-0.5877867) (3.5050000,-0.5868992) (3.5100000,-0.5860144) (3.5150000,-0.5851324) (3.5200000,-0.5842531) (3.5250000,-0.5833765) (3.5300000,-0.5825026) (3.5350000,-0.5816314) (3.5400000,-0.5807629) (3.5450000,-0.5798971) (3.5500000,-0.5790339) (3.5550000,-0.5781733) (3.5600000,-0.5773154) (3.5650000,-0.5764600) (3.5700000,-0.5756073) (3.5750000,-0.5747572) (3.5800000,-0.5739096) (3.5850000,-0.5730646) (3.5900000,-0.5722221) (3.5950000,-0.5713822) (3.6000000,-0.5705449) (3.6050000,-0.5697100) (3.6100000,-0.5688776) (3.6150000,-0.5680478) (3.6200000,-0.5672204) (3.6250000,-0.5663955) (3.6300000,-0.5655730) (3.6350000,-0.5647530) (3.6400000,-0.5639354) (3.6450000,-0.5631203) (3.6500000,-0.5623076) (3.6550000,-0.5614973) (3.6600000,-0.5606893) (3.6650000,-0.5598838) (3.6700000,-0.5590806) (3.6750000,-0.5582798) (3.6800000,-0.5574813) (3.6850000,-0.5566852) (3.6900000,-0.5558914) (3.6950000,-0.5550999) (3.7000000,-0.5543107) (3.7050000,-0.5535239) (3.7100000,-0.5527393) (3.7150000,-0.5519570) (3.7200000,-0.5511769) (3.7250000,-0.5503991) (3.7300000,-0.5496236) (3.7350000,-0.5488503) (3.7400000,-0.5480792) (3.7450000,-0.5473104) (3.7500000,-0.5465437) (3.7550000,-0.5457793) (3.7600000,-0.5450170) (3.7650000,-0.5442569) (3.7700000,-0.5434990) (3.7750000,-0.5427432) (3.7800000,-0.5419896) (3.7850000,-0.5412382) (3.7900000,-0.5404888) (3.7950000,-0.5397416) (3.8000000,-0.5389965) (3.8050000,-0.5382535) (3.8100000,-0.5375126) (3.8150000,-0.5367738) (3.8200000,-0.5360370) (3.8250000,-0.5353024) (3.8300000,-0.5345698) (3.8350000,-0.5338392) (3.8400000,-0.5331107) (3.8450000,-0.5323842) (3.8500000,-0.5316597) (3.8550000,-0.5309373) (3.8600000,-0.5302168) (3.8650000,-0.5294984) (3.8700000,-0.5287819) (3.8750000,-0.5280674) (3.8800000,-0.5273549) (3.8850000,-0.5266444) (3.8900000,-0.5259358) (3.8950000,-0.5252292) (3.9000000,-0.5245245) (3.9050000,-0.5238217) (3.9100000,-0.5231209) (3.9150000,-0.5224219) (3.9200000,-0.5217249) (3.9250000,-0.5210298) (3.9300000,-0.5203366) (3.9350000,-0.5196452) (3.9400000,-0.5189557) (3.9450000,-0.5182681) (3.9500000,-0.5175824) (3.9550000,-0.5168985) (3.9600000,-0.5162165) (3.9650000,-0.5155363) (3.9700000,-0.5148579) (3.9750000,-0.5141813) (3.9800000,-0.5135066) (3.9850000,-0.5128337) (3.9900000,-0.5121625) (3.9950000,-0.5114932) (4.0000000,-0.5108256)};
\draw[orange!85!black,thick,dashed] plot coordinates {(-4.0000000,-0.5108256) (-3.9950000,-0.5114932) (-3.9900000,-0.5121625) (-3.9850000,-0.5128337) (-3.9800000,-0.5135066) (-3.9750000,-0.5141813) (-3.9700000,-0.5148579) (-3.9650000,-0.5155363) (-3.9600000,-0.5162165) (-3.9550000,-0.5168985) (-3.9500000,-0.5175824) (-3.9450000,-0.5182681) (-3.9400000,-0.5189557) (-3.9350000,-0.5196452) (-3.9300000,-0.5203366) (-3.9250000,-0.5210298) (-3.9200000,-0.5217249) (-3.9150000,-0.5224219) (-3.9100000,-0.5231209) (-3.9050000,-0.5238217) (-3.9000000,-0.5245245) (-3.8950000,-0.5252292) (-3.8900000,-0.5259358) (-3.8850000,-0.5266444) (-3.8800000,-0.5273549) (-3.8750000,-0.5280674) (-3.8700000,-0.5287819) (-3.8650000,-0.5294984) (-3.8600000,-0.5302168) (-3.8550000,-0.5309373) (-3.8500000,-0.5316597) (-3.8450000,-0.5323842) (-3.8400000,-0.5331107) (-3.8350000,-0.5338392) (-3.8300000,-0.5345698) (-3.8250000,-0.5353024) (-3.8200000,-0.5360370) (-3.8150000,-0.5367738) (-3.8100000,-0.5375126) (-3.8050000,-0.5382535) (-3.8000000,-0.5389965) (-3.7950000,-0.5397416) (-3.7900000,-0.5404888) (-3.7850000,-0.5412382) (-3.7800000,-0.5419896) (-3.7750000,-0.5427432) (-3.7700000,-0.5434990) (-3.7650000,-0.5442569) (-3.7600000,-0.5450170) (-3.7550000,-0.5457793) (-3.7500000,-0.5465437) (-3.7450000,-0.5473104) (-3.7400000,-0.5480792) (-3.7350000,-0.5488503) (-3.7300000,-0.5496236) (-3.7250000,-0.5503991) (-3.7200000,-0.5511769) (-3.7150000,-0.5519570) (-3.7100000,-0.5527393) (-3.7050000,-0.5535239) (-3.7000000,-0.5543107) (-3.6950000,-0.5550999) (-3.6900000,-0.5558914) (-3.6850000,-0.5566852) (-3.6800000,-0.5574813) (-3.6750000,-0.5582798) (-3.6700000,-0.5590806) (-3.6650000,-0.5598838) (-3.6600000,-0.5606893) (-3.6550000,-0.5614973) (-3.6500000,-0.5623076) (-3.6450000,-0.5631203) (-3.6400000,-0.5639354) (-3.6350000,-0.5647530) (-3.6300000,-0.5655730) (-3.6250000,-0.5663955) (-3.6200000,-0.5672204) (-3.6150000,-0.5680478) (-3.6100000,-0.5688776) (-3.6050000,-0.5697100) (-3.6000000,-0.5705449) (-3.5950000,-0.5713822) (-3.5900000,-0.5722221) (-3.5850000,-0.5730646) (-3.5800000,-0.5739096) (-3.5750000,-0.5747572) (-3.5700000,-0.5756073) (-3.5650000,-0.5764600) (-3.5600000,-0.5773154) (-3.5550000,-0.5781733) (-3.5500000,-0.5790339) (-3.5450000,-0.5798971) (-3.5400000,-0.5807629) (-3.5350000,-0.5816314) (-3.5300000,-0.5825026) (-3.5250000,-0.5833765) (-3.5200000,-0.5842531) (-3.5150000,-0.5851324) (-3.5100000,-0.5860144) (-3.5050000,-0.5868992) (-3.5000000,-0.5877867) (-3.4950000,-0.5886769) (-3.4900000,-0.5895700) (-3.4850000,-0.5904658) (-3.4800000,-0.5913645) (-3.4750000,-0.5922660) (-3.4700000,-0.5931703) (-3.4650000,-0.5940774) (-3.4600000,-0.5949874) (-3.4550000,-0.5959003) (-3.4500000,-0.5968161) (-3.4450000,-0.5977347) (-3.4400000,-0.5986563) (-3.4350000,-0.5995809) (-3.4300000,-0.6005083) (-3.4250000,-0.6014388) (-3.4200000,-0.6023722) (-3.4150000,-0.6033085) (-3.4100000,-0.6042479) (-3.4050000,-0.6051904) (-3.4000000,-0.6061358) (-3.3950000,-0.6070843) (-3.3900000,-0.6080359) (-3.3850000,-0.6089905) (-3.3800000,-0.6099482) (-3.3750000,-0.6109091) (-3.3700000,-0.6118731) (-3.3650000,-0.6128402) (-3.3600000,-0.6138104) (-3.3550000,-0.6147839) (-3.3500000,-0.6157605) (-3.3450000,-0.6167404) (-3.3400000,-0.6177234) (-3.3350000,-0.6187097) (-3.3300000,-0.6196993) (-3.3250000,-0.6206921) (-3.3200000,-0.6216882) (-3.3150000,-0.6226876) (-3.3100000,-0.6236904) (-3.3050000,-0.6246965) (-3.3000000,-0.6257059) (-3.2950000,-0.6267187) (-3.2900000,-0.6277349) (-3.2850000,-0.6287545) (-3.2800000,-0.6297776) (-3.2750000,-0.6308040) (-3.2700000,-0.6318340) (-3.2650000,-0.6328674) (-3.2600000,-0.6339043) (-3.2550000,-0.6349448) (-3.2500000,-0.6359888) (-3.2450000,-0.6370363) (-3.2400000,-0.6380874) (-3.2350000,-0.6391421) (-3.2300000,-0.6402004) (-3.2250000,-0.6412623) (-3.2200000,-0.6423279) (-3.2150000,-0.6433972) (-3.2100000,-0.6444701) (-3.2050000,-0.6455468) (-3.2000000,-0.6466272) (-3.1950000,-0.6477113) (-3.1900000,-0.6487992) (-3.1850000,-0.6498909) (-3.1800000,-0.6509864) (-3.1750000,-0.6520857) (-3.1700000,-0.6531889) (-3.1650000,-0.6542959) (-3.1600000,-0.6554069) (-3.1550000,-0.6565217) (-3.1500000,-0.6576405) (-3.1450000,-0.6587632) (-3.1400000,-0.6598900) (-3.1350000,-0.6610207) (-3.1300000,-0.6621554) (-3.1250000,-0.6632942) (-3.1200000,-0.6644371) (-3.1150000,-0.6655840) (-3.1100000,-0.6667351) (-3.1050000,-0.6678903) (-3.1000000,-0.6690496) (-3.0950000,-0.6702132) (-3.0900000,-0.6713809) (-3.0850000,-0.6725529) (-3.0800000,-0.6737291) (-3.0750000,-0.6749096) (-3.0700000,-0.6760944) (-3.0650000,-0.6772835) (-3.0600000,-0.6784770) (-3.0550000,-0.6796748) (-3.0500000,-0.6808771) (-3.0450000,-0.6820838) (-3.0400000,-0.6832949) (-3.0350000,-0.6845105) (-3.0300000,-0.6857306) (-3.0250000,-0.6869552) (-3.0200000,-0.6881844) (-3.0150000,-0.6894182) (-3.0100000,-0.6906565) (-3.0050000,-0.6918995) (-3.0000000,-0.6931472) (-2.9950000,-0.6943995) (-2.9900000,-0.6956566) (-2.9850000,-0.6969184) (-2.9800000,-0.6981850) (-2.9750000,-0.6994563) (-2.9700000,-0.7007326) (-2.9650000,-0.7020136) (-2.9600000,-0.7032996) (-2.9550000,-0.7045904) (-2.9500000,-0.7058862) (-2.9450000,-0.7071870) (-2.9400000,-0.7084928) (-2.9350000,-0.7098036) (-2.9300000,-0.7111194) (-2.9250000,-0.7124404) (-2.9200000,-0.7137665) (-2.9150000,-0.7150977) (-2.9100000,-0.7164341) (-2.9050000,-0.7177758) (-2.9000000,-0.7191227) (-2.8950000,-0.7204748) (-2.8900000,-0.7218323) (-2.8850000,-0.7231952) (-2.8800000,-0.7245634) (-2.8750000,-0.7259370) (-2.8700000,-0.7273161) (-2.8650000,-0.7287006) (-2.8600000,-0.7300907) (-2.8550000,-0.7314863) (-2.8500000,-0.7328875) (-2.8450000,-0.7342943) (-2.8400000,-0.7357068) (-2.8350000,-0.7371250) (-2.8300000,-0.7385488) (-2.8250000,-0.7399785) (-2.8200000,-0.7414139) (-2.8150000,-0.7428552) (-2.8100000,-0.7443023) (-2.8050000,-0.7457554) (-2.8000000,-0.7472144) (-2.7950000,-0.7486794) (-2.7900000,-0.7501504) (-2.7850000,-0.7516275) (-2.7800000,-0.7531106) (-2.7750000,-0.7546000) (-2.7700000,-0.7560955) (-2.7650000,-0.7575972) (-2.7600000,-0.7591051) (-2.7550000,-0.7606194) (-2.7500000,-0.7621401) (-2.7450000,-0.7636671) (-2.7400000,-0.7652005) (-2.7350000,-0.7667404) (-2.7300000,-0.7682868) (-2.7250000,-0.7698398) (-2.7200000,-0.7713994) (-2.7150000,-0.7729656) (-2.7100000,-0.7745385) (-2.7050000,-0.7761181) (-2.7000000,-0.7777046) (-2.6950000,-0.7792978) (-2.6900000,-0.7808979) (-2.6850000,-0.7825050) (-2.6800000,-0.7841190) (-2.6750000,-0.7857400) (-2.6700000,-0.7873680) (-2.6650000,-0.7890032) (-2.6600000,-0.7906455) (-2.6550000,-0.7922951) (-2.6500000,-0.7939519) (-2.6450000,-0.7956160) (-2.6400000,-0.7972874) (-2.6350000,-0.7989663) (-2.6300000,-0.8006526) (-2.6250000,-0.8023465) (-2.6200000,-0.8040479) (-2.6150000,-0.8057569) (-2.6100000,-0.8074736) (-2.6050000,-0.8091980) (-2.6000000,-0.8109302) (-2.5950000,-0.8126703) (-2.5900000,-0.8144182) (-2.5850000,-0.8161741) (-2.5800000,-0.8179380) (-2.5750000,-0.8197099) (-2.5700000,-0.8214900) (-2.5650000,-0.8232782) (-2.5600000,-0.8250747) (-2.5550000,-0.8268795) (-2.5500000,-0.8286927) (-2.5450000,-0.8305142) (-2.5400000,-0.8323443) (-2.5350000,-0.8341829) (-2.5300000,-0.8360301) (-2.5250000,-0.8378860) (-2.5200000,-0.8397507) (-2.5150000,-0.8416241) (-2.5100000,-0.8435064) (-2.5050000,-0.8453976) (-2.5000000,-0.8472979) (-2.4950000,-0.8492072) (-2.4900000,-0.8511256) (-2.4850000,-0.8530533) (-2.4800000,-0.8549902) (-2.4750000,-0.8569365) (-2.4700000,-0.8588922) (-2.4650000,-0.8608574) (-2.4600000,-0.8628322) (-2.4550000,-0.8648166) (-2.4500000,-0.8668107) (-2.4450000,-0.8688146) (-2.4400000,-0.8708284) (-2.4350000,-0.8728521) (-2.4300000,-0.8748858) (-2.4250000,-0.8769297) (-2.4200000,-0.8789837) (-2.4150000,-0.8810480) (-2.4100000,-0.8831226) (-2.4050000,-0.8852076) (-2.4000000,-0.8873032) (-2.3950000,-0.8894093) (-2.3900000,-0.8915262) (-2.3850000,-0.8936538) (-2.3800000,-0.8957922) (-2.3750000,-0.8979416) (-2.3700000,-0.9001020) (-2.3650000,-0.9022735) (-2.3600000,-0.9044563) (-2.3550000,-0.9066503) (-2.3500000,-0.9088558) (-2.3450000,-0.9110727) (-2.3400000,-0.9133012) (-2.3350000,-0.9155414) (-2.3300000,-0.9177934) (-2.3250000,-0.9200572) (-2.3200000,-0.9223330) (-2.3150000,-0.9246210) (-2.3100000,-0.9269211) (-2.3050000,-0.9292334) (-2.3000000,-0.9315582) (-2.2950000,-0.9338955) (-2.2900000,-0.9362453) (-2.2850000,-0.9386079) (-2.2800000,-0.9409833) (-2.2750000,-0.9433717) (-2.2700000,-0.9457731) (-2.2650000,-0.9481876) (-2.2600000,-0.9506155) (-2.2550000,-0.9530567) (-2.2500000,-0.9555114) (-2.2450000,-0.9579798) (-2.2400000,-0.9604620) (-2.2350000,-0.9629580) (-2.2300000,-0.9654680) (-2.2250000,-0.9679921) (-2.2200000,-0.9705305) (-2.2150000,-0.9730833) (-2.2100000,-0.9756506) (-2.2050000,-0.9782325) (-2.2000000,-0.9808293) (-2.1950000,-0.9834409) (-2.1900000,-0.9860676) (-2.1850000,-0.9887095) (-2.1800000,-0.9913668) (-2.1750000,-0.9940395) (-2.1700000,-0.9967278) (-2.1650000,-0.9994320) (-2.1600000,-1.0021520) (-2.1550000,-1.0048882) (-2.1500000,-1.0076405) (-2.1450000,-1.0104093) (-2.1400000,-1.0131945) (-2.1350000,-1.0159965) (-2.1300000,-1.0188154) (-2.1250000,-1.0216512) (-2.1200000,-1.0245043) (-2.1150000,-1.0273747) (-2.1100000,-1.0302627) (-2.1050000,-1.0331684) (-2.1000000,-1.0360919) (-2.0950000,-1.0390335) (-2.0900000,-1.0419934) (-2.0850000,-1.0449717) (-2.0800000,-1.0479686) (-2.0750000,-1.0509842) (-2.0700000,-1.0540189) (-2.0650000,-1.0570728) (-2.0600000,-1.0601460) (-2.0550000,-1.0632388) (-2.0500000,-1.0663514) (-2.0450000,-1.0694840) (-2.0400000,-1.0726368) (-2.0350000,-1.0758100) (-2.0300000,-1.0790038) (-2.0250000,-1.0822185) (-2.0200000,-1.0854542) (-2.0150000,-1.0887112) (-2.0100000,-1.0919897) (-2.0050000,-1.0952900) (-2.0000000,-1.0986123) (-1.9950000,-1.1019568) (-1.9900000,-1.1053237) (-1.9850000,-1.1087134) (-1.9800000,-1.1121260) (-1.9750000,-1.1155618) (-1.9700000,-1.1190212) (-1.9650000,-1.1225042) (-1.9600000,-1.1260113) (-1.9550000,-1.1295426) (-1.9500000,-1.1330985) (-1.9450000,-1.1366792) (-1.9400000,-1.1402850) (-1.9350000,-1.1439162) (-1.9300000,-1.1475731) (-1.9250000,-1.1512560) (-1.9200000,-1.1549652) (-1.9150000,-1.1587010) (-1.9100000,-1.1624638) (-1.9050000,-1.1662537) (-1.9000000,-1.1700713) (-1.8950000,-1.1739167) (-1.8900000,-1.1777903) (-1.8850000,-1.1816925) (-1.8800000,-1.1856237) (-1.8750000,-1.1895841) (-1.8700000,-1.1935741) (-1.8650000,-1.1975941) (-1.8600000,-1.2016445) (-1.8550000,-1.2057257) (-1.8500000,-1.2098379) (-1.8450000,-1.2139817) (-1.8400000,-1.2181574) (-1.8350000,-1.2223655) (-1.8300000,-1.2266063) (-1.8250000,-1.2308803) (-1.8200000,-1.2351878) (-1.8150000,-1.2395294) (-1.8100000,-1.2439055) (-1.8050000,-1.2483165) (-1.8000000,-1.2527630) (-1.7950000,-1.2572453) (-1.7900000,-1.2617639) (-1.7850000,-1.2663194) (-1.7800000,-1.2709123) (-1.7750000,-1.2755430) (-1.7700000,-1.2802121) (-1.7650000,-1.2849201) (-1.7600000,-1.2896675) (-1.7550000,-1.2944550) (-1.7500000,-1.2992830) (-1.7450000,-1.3041521) (-1.7400000,-1.3090630) (-1.7350000,-1.3140162) (-1.7300000,-1.3190124) (-1.7250000,-1.3240521) (-1.7200000,-1.3291359) (-1.7150000,-1.3342647) (-1.7100000,-1.3394389) (-1.7050000,-1.3446594) (-1.7000000,-1.3499267) (-1.6950000,-1.3552416) (-1.6900000,-1.3606049) (-1.6850000,-1.3660172) (-1.6800000,-1.3714793) (-1.6750000,-1.3769920) (-1.6700000,-1.3825560) (-1.6650000,-1.3881723) (-1.6600000,-1.3938416) (-1.6550000,-1.3995647) (-1.6500000,-1.4053426) (-1.6450000,-1.4111760) (-1.6400000,-1.4170660) (-1.6350000,-1.4230135) (-1.6300000,-1.4290193) (-1.6250000,-1.4350845) (-1.6200000,-1.4412101) (-1.6150000,-1.4473971) (-1.6100000,-1.4536465) (-1.6050000,-1.4599595) (-1.6000000,-1.4663371) (-1.5950000,-1.4727804) (-1.5900000,-1.4792906) (-1.5850000,-1.4858689) (-1.5800000,-1.4925166) (-1.5750000,-1.4992348) (-1.5700000,-1.5060248) (-1.5650000,-1.5128880) (-1.5600000,-1.5198258) (-1.5550000,-1.5268394) (-1.5500000,-1.5339304) (-1.5450000,-1.5411001) (-1.5400000,-1.5483502) (-1.5350000,-1.5556822) (-1.5300000,-1.5630976) (-1.5250000,-1.5705981) (-1.5200000,-1.5781854) (-1.5150000,-1.5858612) (-1.5100000,-1.5936273) (-1.5050000,-1.6014856) (-1.5000000,-1.6094379) (-1.4950000,-1.6174862) (-1.4900000,-1.6256326) (-1.4850000,-1.6338790) (-1.4800000,-1.6422277) (-1.4750000,-1.6506809) (-1.4700000,-1.6592407) (-1.4650000,-1.6679097) (-1.4600000,-1.6766901) (-1.4550000,-1.6855846) (-1.4500000,-1.6945957) (-1.4450000,-1.7037261) (-1.4400000,-1.7129786) (-1.4350000,-1.7223560) (-1.4300000,-1.7318613) (-1.4250000,-1.7414976) (-1.4200000,-1.7512681) (-1.4150000,-1.7611760) (-1.4100000,-1.7712249) (-1.4050000,-1.7814181) (-1.4000000,-1.7917595) (-1.3950000,-1.8022527) (-1.3900000,-1.8129019) (-1.3850000,-1.8237111) (-1.3800000,-1.8346845) (-1.3750000,-1.8458267) (-1.3700000,-1.8571422) (-1.3650000,-1.8686359) (-1.3600000,-1.8803129) (-1.3550000,-1.8921782) (-1.3500000,-1.9042375) (-1.3450000,-1.9164963) (-1.3400000,-1.9289606) (-1.3350000,-1.9416366) (-1.3300000,-1.9545309) (-1.3250000,-1.9676501) (-1.3200000,-1.9810015) (-1.3150000,-1.9945923) (-1.3100000,-2.0084305) (-1.3050000,-2.0225242) (-1.3000000,-2.0368819) (-1.2950000,-2.0515128) (-1.2900000,-2.0664262) (-1.2850000,-2.0816321) (-1.2800000,-2.0971411) (-1.2750000,-2.1129642) (-1.2700000,-2.1291132) (-1.2650000,-2.1456002) (-1.2600000,-2.1624385) (-1.2550000,-2.1796417) (-1.2500000,-2.1972246) (-1.2450000,-2.2152026) (-1.2400000,-2.2335922) (-1.2350000,-2.2524110) (-1.2300000,-2.2716776) (-1.2250000,-2.2914118) (-1.2200000,-2.3116349) (-1.2150000,-2.3323697) (-1.2100000,-2.3536403) (-1.2050000,-2.3754728) (-1.2000000,-2.3978953) (-1.1950000,-2.4209378) (-1.1900000,-2.4446328) (-1.1850000,-2.4690153) (-1.1800000,-2.4941233) (-1.1750000,-2.5199980) (-1.1700000,-2.5466840) (-1.1650000,-2.5742302) (-1.1600000,-2.6026897) (-1.1550000,-2.6321209) (-1.1500000,-2.6625878) (-1.1450000,-2.6941611) (-1.1400000,-2.7269187) (-1.1350000,-2.7609471) (-1.1300000,-2.7963428) (-1.1250000,-2.8332133) (-1.1200000,-2.8716796) (-1.1150000,-2.9118780) (-1.1100000,-2.9539629) (-1.1050000,-2.9981104) (-1.1000000,-3.0445224) (-1.0950000,-3.0934319) (-1.0900000,-3.1451097) (-1.0850000,-3.1998729) (-1.0800000,-3.2580965) (-1.0750000,-3.3202283) (-1.0700000,-3.3868086) (-1.0650000,-3.4584982) (-1.0600000,-3.5361167) (-1.0550000,-3.6206979) (-1.0500000,-3.7135721) (-1.0450000,-3.8164906) (-1.0400000,-3.9318256)};
\draw[orange!85!black,thick,dashed] plot coordinates {(1.0400000,3.9318256) (1.0450000,3.8164906) (1.0500000,3.7135721) (1.0550000,3.6206979) (1.0600000,3.5361167) (1.0650000,3.4584982) (1.0700000,3.3868086) (1.0750000,3.3202283) (1.0800000,3.2580965) (1.0850000,3.1998729) (1.0900000,3.1451097) (1.0950000,3.0934319) (1.1000000,3.0445224) (1.1050000,2.9981104) (1.1100000,2.9539629) (1.1150000,2.9118780) (1.1200000,2.8716796) (1.1250000,2.8332133) (1.1300000,2.7963428) (1.1350000,2.7609471) (1.1400000,2.7269187) (1.1450000,2.6941611) (1.1500000,2.6625878) (1.1550000,2.6321209) (1.1600000,2.6026897) (1.1650000,2.5742302) (1.1700000,2.5466840) (1.1750000,2.5199980) (1.1800000,2.4941233) (1.1850000,2.4690153) (1.1900000,2.4446328) (1.1950000,2.4209378) (1.2000000,2.3978953) (1.2050000,2.3754728) (1.2100000,2.3536403) (1.2150000,2.3323697) (1.2200000,2.3116349) (1.2250000,2.2914118) (1.2300000,2.2716776) (1.2350000,2.2524110) (1.2400000,2.2335922) (1.2450000,2.2152026) (1.2500000,2.1972246) (1.2550000,2.1796417) (1.2600000,2.1624385) (1.2650000,2.1456002) (1.2700000,2.1291132) (1.2750000,2.1129642) (1.2800000,2.0971411) (1.2850000,2.0816321) (1.2900000,2.0664262) (1.2950000,2.0515128) (1.3000000,2.0368819) (1.3050000,2.0225242) (1.3100000,2.0084305) (1.3150000,1.9945923) (1.3200000,1.9810015) (1.3250000,1.9676501) (1.3300000,1.9545309) (1.3350000,1.9416366) (1.3400000,1.9289606) (1.3450000,1.9164963) (1.3500000,1.9042375) (1.3550000,1.8921782) (1.3600000,1.8803129) (1.3650000,1.8686359) (1.3700000,1.8571422) (1.3750000,1.8458267) (1.3800000,1.8346845) (1.3850000,1.8237111) (1.3900000,1.8129019) (1.3950000,1.8022527) (1.4000000,1.7917595) (1.4050000,1.7814181) (1.4100000,1.7712249) (1.4150000,1.7611760) (1.4200000,1.7512681) (1.4250000,1.7414976) (1.4300000,1.7318613) (1.4350000,1.7223560) (1.4400000,1.7129786) (1.4450000,1.7037261) (1.4500000,1.6945957) (1.4550000,1.6855846) (1.4600000,1.6766901) (1.4650000,1.6679097) (1.4700000,1.6592407) (1.4750000,1.6506809) (1.4800000,1.6422277) (1.4850000,1.6338790) (1.4900000,1.6256326) (1.4950000,1.6174862) (1.5000000,1.6094379) (1.5050000,1.6014856) (1.5100000,1.5936273) (1.5150000,1.5858612) (1.5200000,1.5781854) (1.5250000,1.5705981) (1.5300000,1.5630976) (1.5350000,1.5556822) (1.5400000,1.5483502) (1.5450000,1.5411001) (1.5500000,1.5339304) (1.5550000,1.5268394) (1.5600000,1.5198258) (1.5650000,1.5128880) (1.5700000,1.5060248) (1.5750000,1.4992348) (1.5800000,1.4925166) (1.5850000,1.4858689) (1.5900000,1.4792906) (1.5950000,1.4727804) (1.6000000,1.4663371) (1.6050000,1.4599595) (1.6100000,1.4536465) (1.6150000,1.4473971) (1.6200000,1.4412101) (1.6250000,1.4350845) (1.6300000,1.4290193) (1.6350000,1.4230135) (1.6400000,1.4170660) (1.6450000,1.4111760) (1.6500000,1.4053426) (1.6550000,1.3995647) (1.6600000,1.3938416) (1.6650000,1.3881723) (1.6700000,1.3825560) (1.6750000,1.3769920) (1.6800000,1.3714793) (1.6850000,1.3660172) (1.6900000,1.3606049) (1.6950000,1.3552416) (1.7000000,1.3499267) (1.7050000,1.3446594) (1.7100000,1.3394389) (1.7150000,1.3342647) (1.7200000,1.3291359) (1.7250000,1.3240521) (1.7300000,1.3190124) (1.7350000,1.3140162) (1.7400000,1.3090630) (1.7450000,1.3041521) (1.7500000,1.2992830) (1.7550000,1.2944550) (1.7600000,1.2896675) (1.7650000,1.2849201) (1.7700000,1.2802121) (1.7750000,1.2755430) (1.7800000,1.2709123) (1.7850000,1.2663194) (1.7900000,1.2617639) (1.7950000,1.2572453) (1.8000000,1.2527630) (1.8050000,1.2483165) (1.8100000,1.2439055) (1.8150000,1.2395294) (1.8200000,1.2351878) (1.8250000,1.2308803) (1.8300000,1.2266063) (1.8350000,1.2223655) (1.8400000,1.2181574) (1.8450000,1.2139817) (1.8500000,1.2098379) (1.8550000,1.2057257) (1.8600000,1.2016445) (1.8650000,1.1975941) (1.8700000,1.1935741) (1.8750000,1.1895841) (1.8800000,1.1856237) (1.8850000,1.1816925) (1.8900000,1.1777903) (1.8950000,1.1739167) (1.9000000,1.1700713) (1.9050000,1.1662537) (1.9100000,1.1624638) (1.9150000,1.1587010) (1.9200000,1.1549652) (1.9250000,1.1512560) (1.9300000,1.1475731) (1.9350000,1.1439162) (1.9400000,1.1402850) (1.9450000,1.1366792) (1.9500000,1.1330985) (1.9550000,1.1295426) (1.9600000,1.1260113) (1.9650000,1.1225042) (1.9700000,1.1190212) (1.9750000,1.1155618) (1.9800000,1.1121260) (1.9850000,1.1087134) (1.9900000,1.1053237) (1.9950000,1.1019568) (2.0000000,1.0986123) (2.0050000,1.0952900) (2.0100000,1.0919897) (2.0150000,1.0887112) (2.0200000,1.0854542) (2.0250000,1.0822185) (2.0300000,1.0790038) (2.0350000,1.0758100) (2.0400000,1.0726368) (2.0450000,1.0694840) (2.0500000,1.0663514) (2.0550000,1.0632388) (2.0600000,1.0601460) (2.0650000,1.0570728) (2.0700000,1.0540189) (2.0750000,1.0509842) (2.0800000,1.0479686) (2.0850000,1.0449717) (2.0900000,1.0419934) (2.0950000,1.0390335) (2.1000000,1.0360919) (2.1050000,1.0331684) (2.1100000,1.0302627) (2.1150000,1.0273747) (2.1200000,1.0245043) (2.1250000,1.0216512) (2.1300000,1.0188154) (2.1350000,1.0159965) (2.1400000,1.0131945) (2.1450000,1.0104093) (2.1500000,1.0076405) (2.1550000,1.0048882) (2.1600000,1.0021520) (2.1650000,0.9994320) (2.1700000,0.9967278) (2.1750000,0.9940395) (2.1800000,0.9913668) (2.1850000,0.9887095) (2.1900000,0.9860676) (2.1950000,0.9834409) (2.2000000,0.9808293) (2.2050000,0.9782325) (2.2100000,0.9756506) (2.2150000,0.9730833) (2.2200000,0.9705305) (2.2250000,0.9679921) (2.2300000,0.9654680) (2.2350000,0.9629580) (2.2400000,0.9604620) (2.2450000,0.9579798) (2.2500000,0.9555114) (2.2550000,0.9530567) (2.2600000,0.9506155) (2.2650000,0.9481876) (2.2700000,0.9457731) (2.2750000,0.9433717) (2.2800000,0.9409833) (2.2850000,0.9386079) (2.2900000,0.9362453) (2.2950000,0.9338955) (2.3000000,0.9315582) (2.3050000,0.9292334) (2.3100000,0.9269211) (2.3150000,0.9246210) (2.3200000,0.9223330) (2.3250000,0.9200572) (2.3300000,0.9177934) (2.3350000,0.9155414) (2.3400000,0.9133012) (2.3450000,0.9110727) (2.3500000,0.9088558) (2.3550000,0.9066503) (2.3600000,0.9044563) (2.3650000,0.9022735) (2.3700000,0.9001020) (2.3750000,0.8979416) (2.3800000,0.8957922) (2.3850000,0.8936538) (2.3900000,0.8915262) (2.3950000,0.8894093) (2.4000000,0.8873032) (2.4050000,0.8852076) (2.4100000,0.8831226) (2.4150000,0.8810480) (2.4200000,0.8789837) (2.4250000,0.8769297) (2.4300000,0.8748858) (2.4350000,0.8728521) (2.4400000,0.8708284) (2.4450000,0.8688146) (2.4500000,0.8668107) (2.4550000,0.8648166) (2.4600000,0.8628322) (2.4650000,0.8608574) (2.4700000,0.8588922) (2.4750000,0.8569365) (2.4800000,0.8549902) (2.4850000,0.8530533) (2.4900000,0.8511256) (2.4950000,0.8492072) (2.5000000,0.8472979) (2.5050000,0.8453976) (2.5100000,0.8435064) (2.5150000,0.8416241) (2.5200000,0.8397507) (2.5250000,0.8378860) (2.5300000,0.8360301) (2.5350000,0.8341829) (2.5400000,0.8323443) (2.5450000,0.8305142) (2.5500000,0.8286927) (2.5550000,0.8268795) (2.5600000,0.8250747) (2.5650000,0.8232782) (2.5700000,0.8214900) (2.5750000,0.8197099) (2.5800000,0.8179380) (2.5850000,0.8161741) (2.5900000,0.8144182) (2.5950000,0.8126703) (2.6000000,0.8109302) (2.6050000,0.8091980) (2.6100000,0.8074736) (2.6150000,0.8057569) (2.6200000,0.8040479) (2.6250000,0.8023465) (2.6300000,0.8006526) (2.6350000,0.7989663) (2.6400000,0.7972874) (2.6450000,0.7956160) (2.6500000,0.7939519) (2.6550000,0.7922951) (2.6600000,0.7906455) (2.6650000,0.7890032) (2.6700000,0.7873680) (2.6750000,0.7857400) (2.6800000,0.7841190) (2.6850000,0.7825050) (2.6900000,0.7808979) (2.6950000,0.7792978) (2.7000000,0.7777046) (2.7050000,0.7761181) (2.7100000,0.7745385) (2.7150000,0.7729656) (2.7200000,0.7713994) (2.7250000,0.7698398) (2.7300000,0.7682868) (2.7350000,0.7667404) (2.7400000,0.7652005) (2.7450000,0.7636671) (2.7500000,0.7621401) (2.7550000,0.7606194) (2.7600000,0.7591051) (2.7650000,0.7575972) (2.7700000,0.7560955) (2.7750000,0.7546000) (2.7800000,0.7531106) (2.7850000,0.7516275) (2.7900000,0.7501504) (2.7950000,0.7486794) (2.8000000,0.7472144) (2.8050000,0.7457554) (2.8100000,0.7443023) (2.8150000,0.7428552) (2.8200000,0.7414139) (2.8250000,0.7399785) (2.8300000,0.7385488) (2.8350000,0.7371250) (2.8400000,0.7357068) (2.8450000,0.7342943) (2.8500000,0.7328875) (2.8550000,0.7314863) (2.8600000,0.7300907) (2.8650000,0.7287006) (2.8700000,0.7273161) (2.8750000,0.7259370) (2.8800000,0.7245634) (2.8850000,0.7231952) (2.8900000,0.7218323) (2.8950000,0.7204748) (2.9000000,0.7191227) (2.9050000,0.7177758) (2.9100000,0.7164341) (2.9150000,0.7150977) (2.9200000,0.7137665) (2.9250000,0.7124404) (2.9300000,0.7111194) (2.9350000,0.7098036) (2.9400000,0.7084928) (2.9450000,0.7071870) (2.9500000,0.7058862) (2.9550000,0.7045904) (2.9600000,0.7032996) (2.9650000,0.7020136) (2.9700000,0.7007326) (2.9750000,0.6994563) (2.9800000,0.6981850) (2.9850000,0.6969184) (2.9900000,0.6956566) (2.9950000,0.6943995) (3.0000000,0.6931472) (3.0050000,0.6918995) (3.0100000,0.6906565) (3.0150000,0.6894182) (3.0200000,0.6881844) (3.0250000,0.6869552) (3.0300000,0.6857306) (3.0350000,0.6845105) (3.0400000,0.6832949) (3.0450000,0.6820838) (3.0500000,0.6808771) (3.0550000,0.6796748) (3.0600000,0.6784770) (3.0650000,0.6772835) (3.0700000,0.6760944) (3.0750000,0.6749096) (3.0800000,0.6737291) (3.0850000,0.6725529) (3.0900000,0.6713809) (3.0950000,0.6702132) (3.1000000,0.6690496) (3.1050000,0.6678903) (3.1100000,0.6667351) (3.1150000,0.6655840) (3.1200000,0.6644371) (3.1250000,0.6632942) (3.1300000,0.6621554) (3.1350000,0.6610207) (3.1400000,0.6598900) (3.1450000,0.6587632) (3.1500000,0.6576405) (3.1550000,0.6565217) (3.1600000,0.6554069) (3.1650000,0.6542959) (3.1700000,0.6531889) (3.1750000,0.6520857) (3.1800000,0.6509864) (3.1850000,0.6498909) (3.1900000,0.6487992) (3.1950000,0.6477113) (3.2000000,0.6466272) (3.2050000,0.6455468) (3.2100000,0.6444701) (3.2150000,0.6433972) (3.2200000,0.6423279) (3.2250000,0.6412623) (3.2300000,0.6402004) (3.2350000,0.6391421) (3.2400000,0.6380874) (3.2450000,0.6370363) (3.2500000,0.6359888) (3.2550000,0.6349448) (3.2600000,0.6339043) (3.2650000,0.6328674) (3.2700000,0.6318340) (3.2750000,0.6308040) (3.2800000,0.6297776) (3.2850000,0.6287545) (3.2900000,0.6277349) (3.2950000,0.6267187) (3.3000000,0.6257059) (3.3050000,0.6246965) (3.3100000,0.6236904) (3.3150000,0.6226876) (3.3200000,0.6216882) (3.3250000,0.6206921) (3.3300000,0.6196993) (3.3350000,0.6187097) (3.3400000,0.6177234) (3.3450000,0.6167404) (3.3500000,0.6157605) (3.3550000,0.6147839) (3.3600000,0.6138104) (3.3650000,0.6128402) (3.3700000,0.6118731) (3.3750000,0.6109091) (3.3800000,0.6099482) (3.3850000,0.6089905) (3.3900000,0.6080359) (3.3950000,0.6070843) (3.4000000,0.6061358) (3.4050000,0.6051904) (3.4100000,0.6042479) (3.4150000,0.6033085) (3.4200000,0.6023722) (3.4250000,0.6014388) (3.4300000,0.6005083) (3.4350000,0.5995809) (3.4400000,0.5986563) (3.4450000,0.5977347) (3.4500000,0.5968161) (3.4550000,0.5959003) (3.4600000,0.5949874) (3.4650000,0.5940774) (3.4700000,0.5931703) (3.4750000,0.5922660) (3.4800000,0.5913645) (3.4850000,0.5904658) (3.4900000,0.5895700) (3.4950000,0.5886769) (3.5000000,0.5877867) (3.5050000,0.5868992) (3.5100000,0.5860144) (3.5150000,0.5851324) (3.5200000,0.5842531) (3.5250000,0.5833765) (3.5300000,0.5825026) (3.5350000,0.5816314) (3.5400000,0.5807629) (3.5450000,0.5798971) (3.5500000,0.5790339) (3.5550000,0.5781733) (3.5600000,0.5773154) (3.5650000,0.5764600) (3.5700000,0.5756073) (3.5750000,0.5747572) (3.5800000,0.5739096) (3.5850000,0.5730646) (3.5900000,0.5722221) (3.5950000,0.5713822) (3.6000000,0.5705449) (3.6050000,0.5697100) (3.6100000,0.5688776) (3.6150000,0.5680478) (3.6200000,0.5672204) (3.6250000,0.5663955) (3.6300000,0.5655730) (3.6350000,0.5647530) (3.6400000,0.5639354) (3.6450000,0.5631203) (3.6500000,0.5623076) (3.6550000,0.5614973) (3.6600000,0.5606893) (3.6650000,0.5598838) (3.6700000,0.5590806) (3.6750000,0.5582798) (3.6800000,0.5574813) (3.6850000,0.5566852) (3.6900000,0.5558914) (3.6950000,0.5550999) (3.7000000,0.5543107) (3.7050000,0.5535239) (3.7100000,0.5527393) (3.7150000,0.5519570) (3.7200000,0.5511769) (3.7250000,0.5503991) (3.7300000,0.5496236) (3.7350000,0.5488503) (3.7400000,0.5480792) (3.7450000,0.5473104) (3.7500000,0.5465437) (3.7550000,0.5457793) (3.7600000,0.5450170) (3.7650000,0.5442569) (3.7700000,0.5434990) (3.7750000,0.5427432) (3.7800000,0.5419896) (3.7850000,0.5412382) (3.7900000,0.5404888) (3.7950000,0.5397416) (3.8000000,0.5389965) (3.8050000,0.5382535) (3.8100000,0.5375126) (3.8150000,0.5367738) (3.8200000,0.5360370) (3.8250000,0.5353024) (3.8300000,0.5345698) (3.8350000,0.5338392) (3.8400000,0.5331107) (3.8450000,0.5323842) (3.8500000,0.5316597) (3.8550000,0.5309373) (3.8600000,0.5302168) (3.8650000,0.5294984) (3.8700000,0.5287819) (3.8750000,0.5280674) (3.8800000,0.5273549) (3.8850000,0.5266444) (3.8900000,0.5259358) (3.8950000,0.5252292) (3.9000000,0.5245245) (3.9050000,0.5238217) (3.9100000,0.5231209) (3.9150000,0.5224219) (3.9200000,0.5217249) (3.9250000,0.5210298) (3.9300000,0.5203366) (3.9350000,0.5196452) (3.9400000,0.5189557) (3.9450000,0.5182681) (3.9500000,0.5175824) (3.9550000,0.5168985) (3.9600000,0.5162165) (3.9650000,0.5155363) (3.9700000,0.5148579) (3.9750000,0.5141813) (3.9800000,0.5135066) (3.9850000,0.5128337) (3.9900000,0.5121625) (3.9950000,0.5114932) (4.0000000,0.5108256)};
\draw[violet!75!black,densely dotted] (-2,1.0986123)--(2,-1.0986123);
\fill[violet!75!black] (-2,1.0986123) circle(2.3pt);
\fill[violet!75!black] (2,-1.0986123) circle(2.3pt);
\end{scope}
\end{scope}
\node[anchor=west,align=left,text width=4.35cm] at (8.9,4.4) {固定 $k=1$。\\[8pt]蓝实线：$b=1$\\[8pt]橙虚线：$b=-1$\\[14pt]阴影内没有函数值。\\两条边界均不能取。};
\node[text width=12.8cm,align=center] at (7.1,.65) {$k,b\ne0$；$D=\{x:|x|>|b/k|\}$。图中禁区为 $[-1,1]$。};
\end{tikzpicture}\end{center}


设 $k,b\ne0$。真数为正当且仅当 $(kx-b)(kx+b)>0$，即
$k^2x^2>b^2$，故定义域为 $|x|>|b/k|$，关于原点对称。
令 $R(x)=\dfrac{kx-b}{kx+b}>0$，则
\[
R(-x)=\frac{kx+b}{kx-b}=\frac1{R(x)},\qquad
\ln R(-x)=-\ln R(x).
\]
所以奇4为奇函数。奇5等于 $-\ln R(x)$，也为奇函数。
图中阴影区间对应真数不为正的位置；仅排除分母为零远远不够。
\end{proofbox}
\clearpage
\begin{proofbox}
\textbf{\textcolor{CProof}{证明二（续）：先证明根式真数互为正倒数}}
\begin{center}\begin{tikzpicture}[font=\small,>=stealth]
\draw[gray!40,rounded corners=3pt] (0,0) rectangle (14.2,8.4);
\node at (7.1,7.75) {两个根式真数都为正；互为倒数，取对数后才互为相反数};
\begin{scope}[shift={(4.4,1.8)},x=1.2cm,y=.75cm]
\draw[->,gray] (-2.8,0)--(2.8,0) node[right] {$x$};
\draw[->,gray] (0,0)--(0,6) node[above] {$y$};
\draw[gray,dotted] (-2.8,1)--(2.8,1);
\draw[blue,thick,solid] plot coordinates {(-2.7000000,0.1792360) (-2.6775000,0.1806473) (-2.6550000,0.1820804) (-2.6325000,0.1835356) (-2.6100000,0.1850134) (-2.5875000,0.1865145) (-2.5650000,0.1880392) (-2.5425000,0.1895883) (-2.5200000,0.1911621) (-2.4975000,0.1927614) (-2.4750000,0.1943866) (-2.4525000,0.1960385) (-2.4300000,0.1977176) (-2.4075000,0.1994247) (-2.3850000,0.2011603) (-2.3625000,0.2029252) (-2.3400000,0.2047200) (-2.3175000,0.2065456) (-2.2950000,0.2084027) (-2.2725000,0.2102920) (-2.2500000,0.2122145) (-2.2275000,0.2141708) (-2.2050000,0.2161619) (-2.1825000,0.2181887) (-2.1600000,0.2202521) (-2.1375000,0.2223530) (-2.1150000,0.2244925) (-2.0925000,0.2266715) (-2.0700000,0.2288910) (-2.0475000,0.2311523) (-2.0250000,0.2334563) (-2.0025000,0.2358043) (-1.9800000,0.2381975) (-1.9575000,0.2406370) (-1.9350000,0.2431242) (-1.9125000,0.2456604) (-1.8900000,0.2482469) (-1.8675000,0.2508853) (-1.8450000,0.2535769) (-1.8225000,0.2563233) (-1.8000000,0.2591260) (-1.7775000,0.2619868) (-1.7550000,0.2649072) (-1.7325000,0.2678890) (-1.7100000,0.2709341) (-1.6875000,0.2740444) (-1.6650000,0.2772217) (-1.6425000,0.2804681) (-1.6200000,0.2837857) (-1.5975000,0.2871767) (-1.5750000,0.2906433) (-1.5525000,0.2941879) (-1.5300000,0.2978129) (-1.5075000,0.3015208) (-1.4850000,0.3053142) (-1.4625000,0.3091959) (-1.4400000,0.3131686) (-1.4175000,0.3172352) (-1.3950000,0.3213988) (-1.3725000,0.3256626) (-1.3500000,0.3300298) (-1.3275000,0.3345037) (-1.3050000,0.3390879) (-1.2825000,0.3437860) (-1.2600000,0.3486019) (-1.2375000,0.3535394) (-1.2150000,0.3586026) (-1.1925000,0.3637957) (-1.1700000,0.3691231) (-1.1475000,0.3745894) (-1.1250000,0.3801993) (-1.1025000,0.3859577) (-1.0800000,0.3918696) (-1.0575000,0.3979402) (-1.0350000,0.4041751) (-1.0125000,0.4105798) (-0.9900000,0.4171603) (-0.9675000,0.4239224) (-0.9450000,0.4308725) (-0.9225000,0.4380169) (-0.9000000,0.4453624) (-0.8775000,0.4529158) (-0.8550000,0.4606842) (-0.8325000,0.4686749) (-0.8100000,0.4768955) (-0.7875000,0.4853536) (-0.7650000,0.4940572) (-0.7425000,0.5030145) (-0.7200000,0.5122337) (-0.6975000,0.5217236) (-0.6750000,0.5314929) (-0.6525000,0.5415504) (-0.6300000,0.5519052) (-0.6075000,0.5625668) (-0.5850000,0.5735443) (-0.5625000,0.5848475) (-0.5400000,0.5964858) (-0.5175000,0.6084690) (-0.4950000,0.6208069) (-0.4725000,0.6335092) (-0.4500000,0.6465856) (-0.4275000,0.6600460) (-0.4050000,0.6738999) (-0.3825000,0.6881569) (-0.3600000,0.7028264) (-0.3375000,0.7179176) (-0.3150000,0.7334393) (-0.2925000,0.7494003) (-0.2700000,0.7658089) (-0.2475000,0.7826729) (-0.2250000,0.8000000) (-0.2025000,0.8177971) (-0.1800000,0.8360709) (-0.1575000,0.8548271) (-0.1350000,0.8740714) (-0.1125000,0.8938082) (-0.0900000,0.9140418) (-0.0675000,0.9347755) (-0.0450000,0.9560120) (-0.0225000,0.9777531) (0.0000000,1.0000000) (0.0225000,1.0227531) (0.0450000,1.0460120) (0.0675000,1.0697755) (0.0900000,1.0940418) (0.1125000,1.1188082) (0.1350000,1.1440714) (0.1575000,1.1698271) (0.1800000,1.1960709) (0.2025000,1.2227971) (0.2250000,1.2500000) (0.2475000,1.2776729) (0.2700000,1.3058089) (0.2925000,1.3344003) (0.3150000,1.3634393) (0.3375000,1.3929176) (0.3600000,1.4228264) (0.3825000,1.4531569) (0.4050000,1.4838999) (0.4275000,1.5150460) (0.4500000,1.5465856) (0.4725000,1.5785092) (0.4950000,1.6108069) (0.5175000,1.6434690) (0.5400000,1.6764858) (0.5625000,1.7098475) (0.5850000,1.7435443) (0.6075000,1.7775668) (0.6300000,1.8119052) (0.6525000,1.8465504) (0.6750000,1.8814929) (0.6975000,1.9167236) (0.7200000,1.9522337) (0.7425000,1.9880145) (0.7650000,2.0240572) (0.7875000,2.0603536) (0.8100000,2.0968955) (0.8325000,2.1336749) (0.8550000,2.1706842) (0.8775000,2.2079158) (0.9000000,2.2453624) (0.9225000,2.2830169) (0.9450000,2.3208725) (0.9675000,2.3589224) (0.9900000,2.3971603) (1.0125000,2.4355798) (1.0350000,2.4741751) (1.0575000,2.5129402) (1.0800000,2.5518696) (1.1025000,2.5909577) (1.1250000,2.6301993) (1.1475000,2.6695894) (1.1700000,2.7091231) (1.1925000,2.7487957) (1.2150000,2.7886026) (1.2375000,2.8285394) (1.2600000,2.8686019) (1.2825000,2.9087860) (1.3050000,2.9490879) (1.3275000,2.9895037) (1.3500000,3.0300298) (1.3725000,3.0706626) (1.3950000,3.1113988) (1.4175000,3.1522352) (1.4400000,3.1931686) (1.4625000,3.2341959) (1.4850000,3.2753142) (1.5075000,3.3165208) (1.5300000,3.3578129) (1.5525000,3.3991879) (1.5750000,3.4406433) (1.5975000,3.4821767) (1.6200000,3.5237857) (1.6425000,3.5654681) (1.6650000,3.6072217) (1.6875000,3.6490444) (1.7100000,3.6909341) (1.7325000,3.7328890) (1.7550000,3.7749072) (1.7775000,3.8169868) (1.8000000,3.8591260) (1.8225000,3.9013233) (1.8450000,3.9435769) (1.8675000,3.9858853) (1.8900000,4.0282469) (1.9125000,4.0706604) (1.9350000,4.1131242) (1.9575000,4.1556370) (1.9800000,4.1981975) (2.0025000,4.2408043) (2.0250000,4.2834563) (2.0475000,4.3261523) (2.0700000,4.3688910) (2.0925000,4.4116715) (2.1150000,4.4544925) (2.1375000,4.4973530) (2.1600000,4.5402521) (2.1825000,4.5831887) (2.2050000,4.6261619) (2.2275000,4.6691708) (2.2500000,4.7122145) (2.2725000,4.7552920) (2.2950000,4.7984027) (2.3175000,4.8415456) (2.3400000,4.8847200) (2.3625000,4.9279252) (2.3850000,4.9711603) (2.4075000,5.0144247) (2.4300000,5.0577176) (2.4525000,5.1010385) (2.4750000,5.1443866) (2.4975000,5.1877614) (2.5200000,5.2311621) (2.5425000,5.2745883) (2.5650000,5.3180392) (2.5875000,5.3615145) (2.6100000,5.4050134) (2.6325000,5.4485356) (2.6550000,5.4920804) (2.6775000,5.5356473) (2.7000000,5.5792360)};
\fill[blue] (1,2.4142136) circle(2.3pt);
\draw[orange!85!black,thick,dashed] plot coordinates {(-2.7000000,5.5792360) (-2.6775000,5.5356473) (-2.6550000,5.4920804) (-2.6325000,5.4485356) (-2.6100000,5.4050134) (-2.5875000,5.3615145) (-2.5650000,5.3180392) (-2.5425000,5.2745883) (-2.5200000,5.2311621) (-2.4975000,5.1877614) (-2.4750000,5.1443866) (-2.4525000,5.1010385) (-2.4300000,5.0577176) (-2.4075000,5.0144247) (-2.3850000,4.9711603) (-2.3625000,4.9279252) (-2.3400000,4.8847200) (-2.3175000,4.8415456) (-2.2950000,4.7984027) (-2.2725000,4.7552920) (-2.2500000,4.7122145) (-2.2275000,4.6691708) (-2.2050000,4.6261619) (-2.1825000,4.5831887) (-2.1600000,4.5402521) (-2.1375000,4.4973530) (-2.1150000,4.4544925) (-2.0925000,4.4116715) (-2.0700000,4.3688910) (-2.0475000,4.3261523) (-2.0250000,4.2834563) (-2.0025000,4.2408043) (-1.9800000,4.1981975) (-1.9575000,4.1556370) (-1.9350000,4.1131242) (-1.9125000,4.0706604) (-1.8900000,4.0282469) (-1.8675000,3.9858853) (-1.8450000,3.9435769) (-1.8225000,3.9013233) (-1.8000000,3.8591260) (-1.7775000,3.8169868) (-1.7550000,3.7749072) (-1.7325000,3.7328890) (-1.7100000,3.6909341) (-1.6875000,3.6490444) (-1.6650000,3.6072217) (-1.6425000,3.5654681) (-1.6200000,3.5237857) (-1.5975000,3.4821767) (-1.5750000,3.4406433) (-1.5525000,3.3991879) (-1.5300000,3.3578129) (-1.5075000,3.3165208) (-1.4850000,3.2753142) (-1.4625000,3.2341959) (-1.4400000,3.1931686) (-1.4175000,3.1522352) (-1.3950000,3.1113988) (-1.3725000,3.0706626) (-1.3500000,3.0300298) (-1.3275000,2.9895037) (-1.3050000,2.9490879) (-1.2825000,2.9087860) (-1.2600000,2.8686019) (-1.2375000,2.8285394) (-1.2150000,2.7886026) (-1.1925000,2.7487957) (-1.1700000,2.7091231) (-1.1475000,2.6695894) (-1.1250000,2.6301993) (-1.1025000,2.5909577) (-1.0800000,2.5518696) (-1.0575000,2.5129402) (-1.0350000,2.4741751) (-1.0125000,2.4355798) (-0.9900000,2.3971603) (-0.9675000,2.3589224) (-0.9450000,2.3208725) (-0.9225000,2.2830169) (-0.9000000,2.2453624) (-0.8775000,2.2079158) (-0.8550000,2.1706842) (-0.8325000,2.1336749) (-0.8100000,2.0968955) (-0.7875000,2.0603536) (-0.7650000,2.0240572) (-0.7425000,1.9880145) (-0.7200000,1.9522337) (-0.6975000,1.9167236) (-0.6750000,1.8814929) (-0.6525000,1.8465504) (-0.6300000,1.8119052) (-0.6075000,1.7775668) (-0.5850000,1.7435443) (-0.5625000,1.7098475) (-0.5400000,1.6764858) (-0.5175000,1.6434690) (-0.4950000,1.6108069) (-0.4725000,1.5785092) (-0.4500000,1.5465856) (-0.4275000,1.5150460) (-0.4050000,1.4838999) (-0.3825000,1.4531569) (-0.3600000,1.4228264) (-0.3375000,1.3929176) (-0.3150000,1.3634393) (-0.2925000,1.3344003) (-0.2700000,1.3058089) (-0.2475000,1.2776729) (-0.2250000,1.2500000) (-0.2025000,1.2227971) (-0.1800000,1.1960709) (-0.1575000,1.1698271) (-0.1350000,1.1440714) (-0.1125000,1.1188082) (-0.0900000,1.0940418) (-0.0675000,1.0697755) (-0.0450000,1.0460120) (-0.0225000,1.0227531) (0.0000000,1.0000000) (0.0225000,0.9777531) (0.0450000,0.9560120) (0.0675000,0.9347755) (0.0900000,0.9140418) (0.1125000,0.8938082) (0.1350000,0.8740714) (0.1575000,0.8548271) (0.1800000,0.8360709) (0.2025000,0.8177971) (0.2250000,0.8000000) (0.2475000,0.7826729) (0.2700000,0.7658089) (0.2925000,0.7494003) (0.3150000,0.7334393) (0.3375000,0.7179176) (0.3600000,0.7028264) (0.3825000,0.6881569) (0.4050000,0.6738999) (0.4275000,0.6600460) (0.4500000,0.6465856) (0.4725000,0.6335092) (0.4950000,0.6208069) (0.5175000,0.6084690) (0.5400000,0.5964858) (0.5625000,0.5848475) (0.5850000,0.5735443) (0.6075000,0.5625668) (0.6300000,0.5519052) (0.6525000,0.5415504) (0.6750000,0.5314929) (0.6975000,0.5217236) (0.7200000,0.5122337) (0.7425000,0.5030145) (0.7650000,0.4940572) (0.7875000,0.4853536) (0.8100000,0.4768955) (0.8325000,0.4686749) (0.8550000,0.4606842) (0.8775000,0.4529158) (0.9000000,0.4453624) (0.9225000,0.4380169) (0.9450000,0.4308725) (0.9675000,0.4239224) (0.9900000,0.4171603) (1.0125000,0.4105798) (1.0350000,0.4041751) (1.0575000,0.3979402) (1.0800000,0.3918696) (1.1025000,0.3859577) (1.1250000,0.3801993) (1.1475000,0.3745894) (1.1700000,0.3691231) (1.1925000,0.3637957) (1.2150000,0.3586026) (1.2375000,0.3535394) (1.2600000,0.3486019) (1.2825000,0.3437860) (1.3050000,0.3390879) (1.3275000,0.3345037) (1.3500000,0.3300298) (1.3725000,0.3256626) (1.3950000,0.3213988) (1.4175000,0.3172352) (1.4400000,0.3131686) (1.4625000,0.3091959) (1.4850000,0.3053142) (1.5075000,0.3015208) (1.5300000,0.2978129) (1.5525000,0.2941879) (1.5750000,0.2906433) (1.5975000,0.2871767) (1.6200000,0.2837857) (1.6425000,0.2804681) (1.6650000,0.2772217) (1.6875000,0.2740444) (1.7100000,0.2709341) (1.7325000,0.2678890) (1.7550000,0.2649072) (1.7775000,0.2619868) (1.8000000,0.2591260) (1.8225000,0.2563233) (1.8450000,0.2535769) (1.8675000,0.2508853) (1.8900000,0.2482469) (1.9125000,0.2456604) (1.9350000,0.2431242) (1.9575000,0.2406370) (1.9800000,0.2381975) (2.0025000,0.2358043) (2.0250000,0.2334563) (2.0475000,0.2311523) (2.0700000,0.2288910) (2.0925000,0.2266715) (2.1150000,0.2244925) (2.1375000,0.2223530) (2.1600000,0.2202521) (2.1825000,0.2181887) (2.2050000,0.2161619) (2.2275000,0.2141708) (2.2500000,0.2122145) (2.2725000,0.2102920) (2.2950000,0.2084027) (2.3175000,0.2065456) (2.3400000,0.2047200) (2.3625000,0.2029252) (2.3850000,0.2011603) (2.4075000,0.1994247) (2.4300000,0.1977176) (2.4525000,0.1960385) (2.4750000,0.1943866) (2.4975000,0.1927614) (2.5200000,0.1911621) (2.5425000,0.1895883) (2.5650000,0.1880392) (2.5875000,0.1865145) (2.6100000,0.1850134) (2.6325000,0.1835356) (2.6550000,0.1820804) (2.6775000,0.1806473) (2.7000000,0.1792360)};
\fill[orange!85!black] (1,0.4142136) circle(2.3pt);
\draw[violet,densely dotted] (1,0)--(1,2.41421356);
\node[anchor=north] at (1,-.1) {$1$};
\end{scope}
\node[anchor=west,align=left,text width=4.5cm] at (8.9,4.3) {蓝实线：\\[4pt]$U(x)=\sqrt{x^2+1}+x$\\[10pt]橙虚线：\\[4pt]$V(x)=\sqrt{x^2+1}-x$\\[14pt]$U(1)=\sqrt2+1$\\[10pt]$V(1)=\sqrt2-1$};
\node at (7.1,.65) {$U(x)>0$、$V(x)>0$，且 $U(-x)=V(x)=1/U(x)$。};
\end{tikzpicture}\end{center}


令 $U(x)=\sqrt{x^2+1}+x$、$V(x)=\sqrt{x^2+1}-x$。因为 $\sqrt{x^2+1}>|x|$，两者均为正，且
\[
U(x)V(x)=(x^2+1)-x^2=1,\quad U(-x)=V(x).
\]
由对数运算法则，
\[
\begin{aligned}
\log_a U(-x)&=\log_a V(x)=\log_a\frac1{U(x)}=-\log_a U(x),\\
\log_a V(-x)&=\log_a U(x)=-\log_a V(x).
\end{aligned}
\]
因此奇6、奇7在 $\mathbb R$ 上均为奇函数。图中两个真数始终在横轴上方；它们取对数后才互为相反数。
\end{proofbox}
\clearpage
\begin{proofbox}
\textbf{\textcolor{CProof}{证明三：先拆因子，再用奇乘奇为偶}}
% Formula-based samples; see F083 checks/build_tikz_diagrams.py.
\begin{center}\begin{tikzpicture}[font=\small,>=stealth,line join=round]
\draw[gray!40,rounded corners=3pt] (0,0) rectangle (14.2,8.4);
\node at (7.1,7.75) {偶5：$y=x\left(\dfrac12+\dfrac1{2^x-1}\right)$};
\begin{scope}[shift={(4.400000,1.650000)},x=0.850000cm,y=1.551724cm]
\draw[->,gray] (-4,0)--(4,0) node[right] {$x$};
\draw[->,gray] (0,0)--(0,2.9) node[above] {$y$};
\draw[gray] (-2,-.035/1.551724)--(-2,.035/1.551724);
\node[anchor=north] at (-2,-.07/1.551724) {$-2$};
\draw[gray] (2,-.035/1.551724)--(2,.035/1.551724);
\node[anchor=north] at (2,-.07/1.551724) {$2$};
\begin{scope}
\clip (-4,0) rectangle (4,2.9);
\draw[blue,thick,solid] plot coordinates {(-4.0000000,2.2666667) (-3.9933333,2.2642046) (-3.9866667,2.2617449) (-3.9800000,2.2592879) (-3.9733333,2.2568333) (-3.9666667,2.2543812) (-3.9600000,2.2519317) (-3.9533333,2.2494847) (-3.9466667,2.2470403) (-3.9400000,2.2445984) (-3.9333333,2.2421591) (-3.9266667,2.2397223) (-3.9200000,2.2372881) (-3.9133333,2.2348564) (-3.9066667,2.2324274) (-3.9000000,2.2300009) (-3.8933333,2.2275770) (-3.8866667,2.2251557) (-3.8800000,2.2227370) (-3.8733333,2.2203209) (-3.8666667,2.2179074) (-3.8600000,2.2154965) (-3.8533333,2.2130882) (-3.8466667,2.2106826) (-3.8400000,2.2082796) (-3.8333333,2.2058793) (-3.8266667,2.2034815) (-3.8200000,2.2010865) (-3.8133333,2.1986941) (-3.8066667,2.1963043) (-3.8000000,2.1939173) (-3.7933333,2.1915328) (-3.7866667,2.1891511) (-3.7800000,2.1867721) (-3.7733333,2.1843957) (-3.7666667,2.1820221) (-3.7600000,2.1796511) (-3.7533333,2.1772829) (-3.7466667,2.1749174) (-3.7400000,2.1725546) (-3.7333333,2.1701945) (-3.7266667,2.1678371) (-3.7200000,2.1654825) (-3.7133333,2.1631306) (-3.7066667,2.1607815) (-3.7000000,2.1584352) (-3.6933333,2.1560916) (-3.6866667,2.1537507) (-3.6800000,2.1514126) (-3.6733333,2.1490774) (-3.6666667,2.1467449) (-3.6600000,2.1444151) (-3.6533333,2.1420882) (-3.6466667,2.1397641) (-3.6400000,2.1374428) (-3.6333333,2.1351243) (-3.6266667,2.1328086) (-3.6200000,2.1304958) (-3.6133333,2.1281857) (-3.6066667,2.1258785) (-3.6000000,2.1235742) (-3.5933333,2.1212727) (-3.5866667,2.1189741) (-3.5800000,2.1166783) (-3.5733333,2.1143854) (-3.5666667,2.1120953) (-3.5600000,2.1098082) (-3.5533333,2.1075239) (-3.5466667,2.1052425) (-3.5400000,2.1029640) (-3.5333333,2.1006884) (-3.5266667,2.0984157) (-3.5200000,2.0961459) (-3.5133333,2.0938791) (-3.5066667,2.0916152) (-3.5000000,2.0893542) (-3.4933333,2.0870961) (-3.4866667,2.0848410) (-3.4800000,2.0825888) (-3.4733333,2.0803396) (-3.4666667,2.0780934) (-3.4600000,2.0758501) (-3.4533333,2.0736098) (-3.4466667,2.0713724) (-3.4400000,2.0691381) (-3.4333333,2.0669067) (-3.4266667,2.0646784) (-3.4200000,2.0624530) (-3.4133333,2.0602307) (-3.4066667,2.0580114) (-3.4000000,2.0557950) (-3.3933333,2.0535818) (-3.3866667,2.0513715) (-3.3800000,2.0491643) (-3.3733333,2.0469601) (-3.3666667,2.0447590) (-3.3600000,2.0425610) (-3.3533333,2.0403660) (-3.3466667,2.0381740) (-3.3400000,2.0359852) (-3.3333333,2.0337994) (-3.3266667,2.0316167) (-3.3200000,2.0294371) (-3.3133333,2.0272606) (-3.3066667,2.0250872) (-3.3000000,2.0229169) (-3.2933333,2.0207498) (-3.2866667,2.0185857) (-3.2800000,2.0164248) (-3.2733333,2.0142671) (-3.2666667,2.0121124) (-3.2600000,2.0099609) (-3.2533333,2.0078126) (-3.2466667,2.0056674) (-3.2400000,2.0035254) (-3.2333333,2.0013865) (-3.2266667,1.9992509) (-3.2200000,1.9971184) (-3.2133333,1.9949891) (-3.2066667,1.9928630) (-3.2000000,1.9907401) (-3.1933333,1.9886204) (-3.1866667,1.9865039) (-3.1800000,1.9843907) (-3.1733333,1.9822806) (-3.1666667,1.9801738) (-3.1600000,1.9780703) (-3.1533333,1.9759699) (-3.1466667,1.9738729) (-3.1400000,1.9717790) (-3.1333333,1.9696885) (-3.1266667,1.9676012) (-3.1200000,1.9655172) (-3.1133333,1.9634364) (-3.1066667,1.9613590) (-3.1000000,1.9592848) (-3.0933333,1.9572139) (-3.0866667,1.9551464) (-3.0800000,1.9530821) (-3.0733333,1.9510212) (-3.0666667,1.9489636) (-3.0600000,1.9469093) (-3.0533333,1.9448583) (-3.0466667,1.9428107) (-3.0400000,1.9407665) (-3.0333333,1.9387255) (-3.0266667,1.9366880) (-3.0200000,1.9346538) (-3.0133333,1.9326229) (-3.0066667,1.9305955) (-3.0000000,1.9285714) (-2.9933333,1.9265507) (-2.9866667,1.9245335) (-2.9800000,1.9225196) (-2.9733333,1.9205091) (-2.9666667,1.9185020) (-2.9600000,1.9164984) (-2.9533333,1.9144981) (-2.9466667,1.9125013) (-2.9400000,1.9105080) (-2.9333333,1.9085181) (-2.9266667,1.9065316) (-2.9200000,1.9045486) (-2.9133333,1.9025690) (-2.9066667,1.9005929) (-2.9000000,1.8986203) (-2.8933333,1.8966512) (-2.8866667,1.8946855) (-2.8800000,1.8927233) (-2.8733333,1.8907647) (-2.8666667,1.8888095) (-2.8600000,1.8868578) (-2.8533333,1.8849097) (-2.8466667,1.8829651) (-2.8400000,1.8810240) (-2.8333333,1.8790864) (-2.8266667,1.8771524) (-2.8200000,1.8752219) (-2.8133333,1.8732949) (-2.8066667,1.8713715) (-2.8000000,1.8694517) (-2.7933333,1.8675355) (-2.7866667,1.8656228) (-2.7800000,1.8637137) (-2.7733333,1.8618082) (-2.7666667,1.8599063) (-2.7600000,1.8580080) (-2.7533333,1.8561132) (-2.7466667,1.8542221) (-2.7400000,1.8523346) (-2.7333333,1.8504508) (-2.7266667,1.8485705) (-2.7200000,1.8466939) (-2.7133333,1.8448210) (-2.7066667,1.8429516) (-2.7000000,1.8410860) (-2.6933333,1.8392239) (-2.6866667,1.8373656) (-2.6800000,1.8355109) (-2.6733333,1.8336599) (-2.6666667,1.8318126) (-2.6600000,1.8299689) (-2.6533333,1.8281290) (-2.6466667,1.8262927) (-2.6400000,1.8244602) (-2.6333333,1.8226314) (-2.6266667,1.8208063) (-2.6200000,1.8189849) (-2.6133333,1.8171672) (-2.6066667,1.8153533) (-2.6000000,1.8135431) (-2.5933333,1.8117367) (-2.5866667,1.8099340) (-2.5800000,1.8081350) (-2.5733333,1.8063399) (-2.5666667,1.8045485) (-2.5600000,1.8027608) (-2.5533333,1.8009770) (-2.5466667,1.7991969) (-2.5400000,1.7974207) (-2.5333333,1.7956482) (-2.5266667,1.7938796) (-2.5200000,1.7921147) (-2.5133333,1.7903537) (-2.5066667,1.7885965) (-2.5000000,1.7868431) (-2.4933333,1.7850935) (-2.4866667,1.7833478) (-2.4800000,1.7816060) (-2.4733333,1.7798680) (-2.4666667,1.7781338) (-2.4600000,1.7764035) (-2.4533333,1.7746771) (-2.4466667,1.7729545) (-2.4400000,1.7712359) (-2.4333333,1.7695211) (-2.4266667,1.7678102) (-2.4200000,1.7661032) (-2.4133333,1.7644001) (-2.4066667,1.7627009) (-2.4000000,1.7610057) (-2.3933333,1.7593143) (-2.3866667,1.7576269) (-2.3800000,1.7559434) (-2.3733333,1.7542639) (-2.3666667,1.7525883) (-2.3600000,1.7509166) (-2.3533333,1.7492489) (-2.3466667,1.7475852) (-2.3400000,1.7459254) (-2.3333333,1.7442696) (-2.3266667,1.7426177) (-2.3200000,1.7409699) (-2.3133333,1.7393260) (-2.3066667,1.7376862) (-2.3000000,1.7360503) (-2.2933333,1.7344185) (-2.2866667,1.7327906) (-2.2800000,1.7311668) (-2.2733333,1.7295470) (-2.2666667,1.7279312) (-2.2600000,1.7263195) (-2.2533333,1.7247118) (-2.2466667,1.7231081) (-2.2400000,1.7215085) (-2.2333333,1.7199129) (-2.2266667,1.7183215) (-2.2200000,1.7167340) (-2.2133333,1.7151507) (-2.2066667,1.7135714) (-2.2000000,1.7119962) (-2.1933333,1.7104251) (-2.1866667,1.7088581) (-2.1800000,1.7072952) (-2.1733333,1.7057364) (-2.1666667,1.7041817) (-2.1600000,1.7026312) (-2.1533333,1.7010847) (-2.1466667,1.6995424) (-2.1400000,1.6980042) (-2.1333333,1.6964702) (-2.1266667,1.6949403) (-2.1200000,1.6934145) (-2.1133333,1.6918929) (-2.1066667,1.6903755) (-2.1000000,1.6888622) (-2.0933333,1.6873531) (-2.0866667,1.6858482) (-2.0800000,1.6843474) (-2.0733333,1.6828509) (-2.0666667,1.6813585) (-2.0600000,1.6798703) (-2.0533333,1.6783864) (-2.0466667,1.6769066) (-2.0400000,1.6754311) (-2.0333333,1.6739598) (-2.0266667,1.6724927) (-2.0200000,1.6710298) (-2.0133333,1.6695712) (-2.0066667,1.6681168) (-2.0000000,1.6666667) (-1.9933333,1.6652208) (-1.9866667,1.6637791) (-1.9800000,1.6623418) (-1.9733333,1.6609087) (-1.9666667,1.6594798) (-1.9600000,1.6580553) (-1.9533333,1.6566350) (-1.9466667,1.6552190) (-1.9400000,1.6538073) (-1.9333333,1.6523999) (-1.9266667,1.6509968) (-1.9200000,1.6495981) (-1.9133333,1.6482036) (-1.9066667,1.6468134) (-1.9000000,1.6454276) (-1.8933333,1.6440461) (-1.8866667,1.6426689) (-1.8800000,1.6412961) (-1.8733333,1.6399276) (-1.8666667,1.6385635) (-1.8600000,1.6372037) (-1.8533333,1.6358483) (-1.8466667,1.6344973) (-1.8400000,1.6331506) (-1.8333333,1.6318083) (-1.8266667,1.6304703) (-1.8200000,1.6291368) (-1.8133333,1.6278076) (-1.8066667,1.6264829) (-1.8000000,1.6251625) (-1.7933333,1.6238465) (-1.7866667,1.6225350) (-1.7800000,1.6212279) (-1.7733333,1.6199251) (-1.7666667,1.6186268) (-1.7600000,1.6173330) (-1.7533333,1.6160435) (-1.7466667,1.6147586) (-1.7400000,1.6134780) (-1.7333333,1.6122019) (-1.7266667,1.6109303) (-1.7200000,1.6096631) (-1.7133333,1.6084003) (-1.7066667,1.6071421) (-1.7000000,1.6058883) (-1.6933333,1.6046390) (-1.6866667,1.6033941) (-1.6800000,1.6021538) (-1.6733333,1.6009179) (-1.6666667,1.5996866) (-1.6600000,1.5984597) (-1.6533333,1.5972374) (-1.6466667,1.5960195) (-1.6400000,1.5948062) (-1.6333333,1.5935974) (-1.6266667,1.5923931) (-1.6200000,1.5911934) (-1.6133333,1.5899981) (-1.6066667,1.5888074) (-1.6000000,1.5876213) (-1.5933333,1.5864397) (-1.5866667,1.5852627) (-1.5800000,1.5840902) (-1.5733333,1.5829222) (-1.5666667,1.5817589) (-1.5600000,1.5806001) (-1.5533333,1.5794458) (-1.5466667,1.5782962) (-1.5400000,1.5771511) (-1.5333333,1.5760107) (-1.5266667,1.5748748) (-1.5200000,1.5737435) (-1.5133333,1.5726168) (-1.5066667,1.5714947) (-1.5000000,1.5703772) (-1.4933333,1.5692644) (-1.4866667,1.5681561) (-1.4800000,1.5670525) (-1.4733333,1.5659535) (-1.4666667,1.5648592) (-1.4600000,1.5637694) (-1.4533333,1.5626843) (-1.4466667,1.5616039) (-1.4400000,1.5605281) (-1.4333333,1.5594569) (-1.4266667,1.5583905) (-1.4200000,1.5573286) (-1.4133333,1.5562715) (-1.4066667,1.5552190) (-1.4000000,1.5541711) (-1.3933333,1.5531280) (-1.3866667,1.5520895) (-1.3800000,1.5510557) (-1.3733333,1.5500266) (-1.3666667,1.5490022) (-1.3600000,1.5479825) (-1.3533333,1.5469675) (-1.3466667,1.5459572) (-1.3400000,1.5449517) (-1.3333333,1.5439508) (-1.3266667,1.5429546) (-1.3200000,1.5419632) (-1.3133333,1.5409765) (-1.3066667,1.5399945) (-1.3000000,1.5390173) (-1.2933333,1.5380447) (-1.2866667,1.5370770) (-1.2800000,1.5361140) (-1.2733333,1.5351557) (-1.2666667,1.5342022) (-1.2600000,1.5332534) (-1.2533333,1.5323094) (-1.2466667,1.5313701) (-1.2400000,1.5304357) (-1.2333333,1.5295060) (-1.2266667,1.5285810) (-1.2200000,1.5276609) (-1.2133333,1.5267455) (-1.2066667,1.5258349) (-1.2000000,1.5249291) (-1.1933333,1.5240281) (-1.1866667,1.5231319) (-1.1800000,1.5222405) (-1.1733333,1.5213539) (-1.1666667,1.5204721) (-1.1600000,1.5195951) (-1.1533333,1.5187230) (-1.1466667,1.5178556) (-1.1400000,1.5169931) (-1.1333333,1.5161354) (-1.1266667,1.5152825) (-1.1200000,1.5144345) (-1.1133333,1.5135913) (-1.1066667,1.5127529) (-1.1000000,1.5119194) (-1.0933333,1.5110908) (-1.0866667,1.5102669) (-1.0800000,1.5094480) (-1.0733333,1.5086338) (-1.0666667,1.5078246) (-1.0600000,1.5070202) (-1.0533333,1.5062207) (-1.0466667,1.5054260) (-1.0400000,1.5046362) (-1.0333333,1.5038513) (-1.0266667,1.5030713) (-1.0200000,1.5022961) (-1.0133333,1.5015259) (-1.0066667,1.5007605) (-1.0000000,1.5000000) (-0.9933333,1.4992444) (-0.9866667,1.4984937) (-0.9800000,1.4977479) (-0.9733333,1.4970070) (-0.9666667,1.4962711) (-0.9600000,1.4955400) (-0.9533333,1.4948138) (-0.9466667,1.4940926) (-0.9400000,1.4933763) (-0.9333333,1.4926649) (-0.9266667,1.4919584) (-0.9200000,1.4912568) (-0.9133333,1.4905602) (-0.9066667,1.4898685) (-0.9000000,1.4891818) (-0.8933333,1.4885000) (-0.8866667,1.4878231) (-0.8800000,1.4871512) (-0.8733333,1.4864842) (-0.8666667,1.4858222) (-0.8600000,1.4851651) (-0.8533333,1.4845130) (-0.8466667,1.4838659) (-0.8400000,1.4832237) (-0.8333333,1.4825864) (-0.8266667,1.4819542) (-0.8200000,1.4813269) (-0.8133333,1.4807045) (-0.8066667,1.4800872) (-0.8000000,1.4794748) (-0.7933333,1.4788674) (-0.7866667,1.4782650) (-0.7800000,1.4776676) (-0.7733333,1.4770751) (-0.7666667,1.4764877) (-0.7600000,1.4759052) (-0.7533333,1.4753278) (-0.7466667,1.4747553) (-0.7400000,1.4741878) (-0.7333333,1.4736254) (-0.7266667,1.4730679) (-0.7200000,1.4725154) (-0.7133333,1.4719680) (-0.7066667,1.4714255) (-0.7000000,1.4708881) (-0.6933333,1.4703557) (-0.6866667,1.4698283) (-0.6800000,1.4693059) (-0.6733333,1.4687886) (-0.6666667,1.4682763) (-0.6600000,1.4677690) (-0.6533333,1.4672667) (-0.6466667,1.4667694) (-0.6400000,1.4662772) (-0.6333333,1.4657900) (-0.6266667,1.4653079) (-0.6200000,1.4648308) (-0.6133333,1.4643587) (-0.6066667,1.4638917) (-0.6000000,1.4634298) (-0.5933333,1.4629728) (-0.5866667,1.4625210) (-0.5800000,1.4620741) (-0.5733333,1.4616323) (-0.5666667,1.4611956) (-0.5600000,1.4607640) (-0.5533333,1.4603374) (-0.5466667,1.4599158) (-0.5400000,1.4594993) (-0.5333333,1.4590879) (-0.5266667,1.4586815) (-0.5200000,1.4582802) (-0.5133333,1.4578840) (-0.5066667,1.4574929) (-0.5000000,1.4571068) (-0.4933333,1.4567258) (-0.4866667,1.4563498) (-0.4800000,1.4559790) (-0.4733333,1.4556132) (-0.4666667,1.4552525) (-0.4600000,1.4548969) (-0.4533333,1.4545463) (-0.4466667,1.4542009) (-0.4400000,1.4538605) (-0.4333333,1.4535252) (-0.4266667,1.4531950) (-0.4200000,1.4528699) (-0.4133333,1.4525499) (-0.4066667,1.4522350) (-0.4000000,1.4519252) (-0.3933333,1.4516205) (-0.3866667,1.4513208) (-0.3800000,1.4510263) (-0.3733333,1.4507368) (-0.3666667,1.4504525) (-0.3600000,1.4501733) (-0.3533333,1.4498991) (-0.3466667,1.4496301) (-0.3400000,1.4493662) (-0.3333333,1.4491074) (-0.3266667,1.4488537) (-0.3200000,1.4486051) (-0.3133333,1.4483616) (-0.3066667,1.4481232) (-0.3000000,1.4478899) (-0.2933333,1.4476617) (-0.2866667,1.4474387) (-0.2800000,1.4472208) (-0.2733333,1.4470079) (-0.2666667,1.4468002) (-0.2600000,1.4465977) (-0.2533333,1.4464002) (-0.2466667,1.4462078) (-0.2400000,1.4460206) (-0.2333333,1.4458385) (-0.2266667,1.4456615) (-0.2200000,1.4454897) (-0.2133333,1.4453229) (-0.2066667,1.4451613) (-0.2000000,1.4450048) (-0.1933333,1.4448534) (-0.1866667,1.4447072) (-0.1800000,1.4445661) (-0.1733333,1.4444301) (-0.1666667,1.4442992) (-0.1600000,1.4441735) (-0.1533333,1.4440528) (-0.1466667,1.4439374) (-0.1400000,1.4438270) (-0.1333333,1.4437218) (-0.1266667,1.4436217) (-0.1200000,1.4435267) (-0.1133333,1.4434369) (-0.1066667,1.4433522) (-0.1000000,1.4432726) (-0.0933333,1.4431982) (-0.0866667,1.4431289) (-0.0800000,1.4430647) (-0.0733333,1.4430057) (-0.0666667,1.4429518) (-0.0600000,1.4429030) (-0.0533333,1.4428593) (-0.0466667,1.4428208) (-0.0400000,1.4427875) (-0.0333333,1.4427592) (-0.0266667,1.4427361) (-0.0200000,1.4427181) (-0.0133333,1.4427053) (-0.0066667,1.4426976)};
\draw[blue,thick,solid] plot coordinates {(0.0066667,1.4426976) (0.0133333,1.4427053) (0.0200000,1.4427181) (0.0266667,1.4427361) (0.0333333,1.4427592) (0.0400000,1.4427875) (0.0466667,1.4428208) (0.0533333,1.4428593) (0.0600000,1.4429030) (0.0666667,1.4429518) (0.0733333,1.4430057) (0.0800000,1.4430647) (0.0866667,1.4431289) (0.0933333,1.4431982) (0.1000000,1.4432726) (0.1066667,1.4433522) (0.1133333,1.4434369) (0.1200000,1.4435267) (0.1266667,1.4436217) (0.1333333,1.4437218) (0.1400000,1.4438270) (0.1466667,1.4439374) (0.1533333,1.4440528) (0.1600000,1.4441735) (0.1666667,1.4442992) (0.1733333,1.4444301) (0.1800000,1.4445661) (0.1866667,1.4447072) (0.1933333,1.4448534) (0.2000000,1.4450048) (0.2066667,1.4451613) (0.2133333,1.4453229) (0.2200000,1.4454897) (0.2266667,1.4456615) (0.2333333,1.4458385) (0.2400000,1.4460206) (0.2466667,1.4462078) (0.2533333,1.4464002) (0.2600000,1.4465977) (0.2666667,1.4468002) (0.2733333,1.4470079) (0.2800000,1.4472208) (0.2866667,1.4474387) (0.2933333,1.4476617) (0.3000000,1.4478899) (0.3066667,1.4481232) (0.3133333,1.4483616) (0.3200000,1.4486051) (0.3266667,1.4488537) (0.3333333,1.4491074) (0.3400000,1.4493662) (0.3466667,1.4496301) (0.3533333,1.4498991) (0.3600000,1.4501733) (0.3666667,1.4504525) (0.3733333,1.4507368) (0.3800000,1.4510263) (0.3866667,1.4513208) (0.3933333,1.4516205) (0.4000000,1.4519252) (0.4066667,1.4522350) (0.4133333,1.4525499) (0.4200000,1.4528699) (0.4266667,1.4531950) (0.4333333,1.4535252) (0.4400000,1.4538605) (0.4466667,1.4542009) (0.4533333,1.4545463) (0.4600000,1.4548969) (0.4666667,1.4552525) (0.4733333,1.4556132) (0.4800000,1.4559790) (0.4866667,1.4563498) (0.4933333,1.4567258) (0.5000000,1.4571068) (0.5066667,1.4574929) (0.5133333,1.4578840) (0.5200000,1.4582802) (0.5266667,1.4586815) (0.5333333,1.4590879) (0.5400000,1.4594993) (0.5466667,1.4599158) (0.5533333,1.4603374) (0.5600000,1.4607640) (0.5666667,1.4611956) (0.5733333,1.4616323) (0.5800000,1.4620741) (0.5866667,1.4625210) (0.5933333,1.4629728) (0.6000000,1.4634298) (0.6066667,1.4638917) (0.6133333,1.4643587) (0.6200000,1.4648308) (0.6266667,1.4653079) (0.6333333,1.4657900) (0.6400000,1.4662772) (0.6466667,1.4667694) (0.6533333,1.4672667) (0.6600000,1.4677690) (0.6666667,1.4682763) (0.6733333,1.4687886) (0.6800000,1.4693059) (0.6866667,1.4698283) (0.6933333,1.4703557) (0.7000000,1.4708881) (0.7066667,1.4714255) (0.7133333,1.4719680) (0.7200000,1.4725154) (0.7266667,1.4730679) (0.7333333,1.4736254) (0.7400000,1.4741878) (0.7466667,1.4747553) (0.7533333,1.4753278) (0.7600000,1.4759052) (0.7666667,1.4764877) (0.7733333,1.4770751) (0.7800000,1.4776676) (0.7866667,1.4782650) (0.7933333,1.4788674) (0.8000000,1.4794748) (0.8066667,1.4800872) (0.8133333,1.4807045) (0.8200000,1.4813269) (0.8266667,1.4819542) (0.8333333,1.4825864) (0.8400000,1.4832237) (0.8466667,1.4838659) (0.8533333,1.4845130) (0.8600000,1.4851651) (0.8666667,1.4858222) (0.8733333,1.4864842) (0.8800000,1.4871512) (0.8866667,1.4878231) (0.8933333,1.4885000) (0.9000000,1.4891818) (0.9066667,1.4898685) (0.9133333,1.4905602) (0.9200000,1.4912568) (0.9266667,1.4919584) (0.9333333,1.4926649) (0.9400000,1.4933763) (0.9466667,1.4940926) (0.9533333,1.4948138) (0.9600000,1.4955400) (0.9666667,1.4962711) (0.9733333,1.4970070) (0.9800000,1.4977479) (0.9866667,1.4984937) (0.9933333,1.4992444) (1.0000000,1.5000000) (1.0066667,1.5007605) (1.0133333,1.5015259) (1.0200000,1.5022961) (1.0266667,1.5030713) (1.0333333,1.5038513) (1.0400000,1.5046362) (1.0466667,1.5054260) (1.0533333,1.5062207) (1.0600000,1.5070202) (1.0666667,1.5078246) (1.0733333,1.5086338) (1.0800000,1.5094480) (1.0866667,1.5102669) (1.0933333,1.5110908) (1.1000000,1.5119194) (1.1066667,1.5127529) (1.1133333,1.5135913) (1.1200000,1.5144345) (1.1266667,1.5152825) (1.1333333,1.5161354) (1.1400000,1.5169931) (1.1466667,1.5178556) (1.1533333,1.5187230) (1.1600000,1.5195951) (1.1666667,1.5204721) (1.1733333,1.5213539) (1.1800000,1.5222405) (1.1866667,1.5231319) (1.1933333,1.5240281) (1.2000000,1.5249291) (1.2066667,1.5258349) (1.2133333,1.5267455) (1.2200000,1.5276609) (1.2266667,1.5285810) (1.2333333,1.5295060) (1.2400000,1.5304357) (1.2466667,1.5313701) (1.2533333,1.5323094) (1.2600000,1.5332534) (1.2666667,1.5342022) (1.2733333,1.5351557) (1.2800000,1.5361140) (1.2866667,1.5370770) (1.2933333,1.5380447) (1.3000000,1.5390173) (1.3066667,1.5399945) (1.3133333,1.5409765) (1.3200000,1.5419632) (1.3266667,1.5429546) (1.3333333,1.5439508) (1.3400000,1.5449517) (1.3466667,1.5459572) (1.3533333,1.5469675) (1.3600000,1.5479825) (1.3666667,1.5490022) (1.3733333,1.5500266) (1.3800000,1.5510557) (1.3866667,1.5520895) (1.3933333,1.5531280) (1.4000000,1.5541711) (1.4066667,1.5552190) (1.4133333,1.5562715) (1.4200000,1.5573286) (1.4266667,1.5583905) (1.4333333,1.5594569) (1.4400000,1.5605281) (1.4466667,1.5616039) (1.4533333,1.5626843) (1.4600000,1.5637694) (1.4666667,1.5648592) (1.4733333,1.5659535) (1.4800000,1.5670525) (1.4866667,1.5681561) (1.4933333,1.5692644) (1.5000000,1.5703772) (1.5066667,1.5714947) (1.5133333,1.5726168) (1.5200000,1.5737435) (1.5266667,1.5748748) (1.5333333,1.5760107) (1.5400000,1.5771511) (1.5466667,1.5782962) (1.5533333,1.5794458) (1.5600000,1.5806001) (1.5666667,1.5817589) (1.5733333,1.5829222) (1.5800000,1.5840902) (1.5866667,1.5852627) (1.5933333,1.5864397) (1.6000000,1.5876213) (1.6066667,1.5888074) (1.6133333,1.5899981) (1.6200000,1.5911934) (1.6266667,1.5923931) (1.6333333,1.5935974) (1.6400000,1.5948062) (1.6466667,1.5960195) (1.6533333,1.5972374) (1.6600000,1.5984597) (1.6666667,1.5996866) (1.6733333,1.6009179) (1.6800000,1.6021538) (1.6866667,1.6033941) (1.6933333,1.6046390) (1.7000000,1.6058883) (1.7066667,1.6071421) (1.7133333,1.6084003) (1.7200000,1.6096631) (1.7266667,1.6109303) (1.7333333,1.6122019) (1.7400000,1.6134780) (1.7466667,1.6147586) (1.7533333,1.6160435) (1.7600000,1.6173330) (1.7666667,1.6186268) (1.7733333,1.6199251) (1.7800000,1.6212279) (1.7866667,1.6225350) (1.7933333,1.6238465) (1.8000000,1.6251625) (1.8066667,1.6264829) (1.8133333,1.6278076) (1.8200000,1.6291368) (1.8266667,1.6304703) (1.8333333,1.6318083) (1.8400000,1.6331506) (1.8466667,1.6344973) (1.8533333,1.6358483) (1.8600000,1.6372037) (1.8666667,1.6385635) (1.8733333,1.6399276) (1.8800000,1.6412961) (1.8866667,1.6426689) (1.8933333,1.6440461) (1.9000000,1.6454276) (1.9066667,1.6468134) (1.9133333,1.6482036) (1.9200000,1.6495981) (1.9266667,1.6509968) (1.9333333,1.6523999) (1.9400000,1.6538073) (1.9466667,1.6552190) (1.9533333,1.6566350) (1.9600000,1.6580553) (1.9666667,1.6594798) (1.9733333,1.6609087) (1.9800000,1.6623418) (1.9866667,1.6637791) (1.9933333,1.6652208) (2.0000000,1.6666667) (2.0066667,1.6681168) (2.0133333,1.6695712) (2.0200000,1.6710298) (2.0266667,1.6724927) (2.0333333,1.6739598) (2.0400000,1.6754311) (2.0466667,1.6769066) (2.0533333,1.6783864) (2.0600000,1.6798703) (2.0666667,1.6813585) (2.0733333,1.6828509) (2.0800000,1.6843474) (2.0866667,1.6858482) (2.0933333,1.6873531) (2.1000000,1.6888622) (2.1066667,1.6903755) (2.1133333,1.6918929) (2.1200000,1.6934145) (2.1266667,1.6949403) (2.1333333,1.6964702) (2.1400000,1.6980042) (2.1466667,1.6995424) (2.1533333,1.7010847) (2.1600000,1.7026312) (2.1666667,1.7041817) (2.1733333,1.7057364) (2.1800000,1.7072952) (2.1866667,1.7088581) (2.1933333,1.7104251) (2.2000000,1.7119962) (2.2066667,1.7135714) (2.2133333,1.7151507) (2.2200000,1.7167340) (2.2266667,1.7183215) (2.2333333,1.7199129) (2.2400000,1.7215085) (2.2466667,1.7231081) (2.2533333,1.7247118) (2.2600000,1.7263195) (2.2666667,1.7279312) (2.2733333,1.7295470) (2.2800000,1.7311668) (2.2866667,1.7327906) (2.2933333,1.7344185) (2.3000000,1.7360503) (2.3066667,1.7376862) (2.3133333,1.7393260) (2.3200000,1.7409699) (2.3266667,1.7426177) (2.3333333,1.7442696) (2.3400000,1.7459254) (2.3466667,1.7475852) (2.3533333,1.7492489) (2.3600000,1.7509166) (2.3666667,1.7525883) (2.3733333,1.7542639) (2.3800000,1.7559434) (2.3866667,1.7576269) (2.3933333,1.7593143) (2.4000000,1.7610057) (2.4066667,1.7627009) (2.4133333,1.7644001) (2.4200000,1.7661032) (2.4266667,1.7678102) (2.4333333,1.7695211) (2.4400000,1.7712359) (2.4466667,1.7729545) (2.4533333,1.7746771) (2.4600000,1.7764035) (2.4666667,1.7781338) (2.4733333,1.7798680) (2.4800000,1.7816060) (2.4866667,1.7833478) (2.4933333,1.7850935) (2.5000000,1.7868431) (2.5066667,1.7885965) (2.5133333,1.7903537) (2.5200000,1.7921147) (2.5266667,1.7938796) (2.5333333,1.7956482) (2.5400000,1.7974207) (2.5466667,1.7991969) (2.5533333,1.8009770) (2.5600000,1.8027608) (2.5666667,1.8045485) (2.5733333,1.8063399) (2.5800000,1.8081350) (2.5866667,1.8099340) (2.5933333,1.8117367) (2.6000000,1.8135431) (2.6066667,1.8153533) (2.6133333,1.8171672) (2.6200000,1.8189849) (2.6266667,1.8208063) (2.6333333,1.8226314) (2.6400000,1.8244602) (2.6466667,1.8262927) (2.6533333,1.8281290) (2.6600000,1.8299689) (2.6666667,1.8318126) (2.6733333,1.8336599) (2.6800000,1.8355109) (2.6866667,1.8373656) (2.6933333,1.8392239) (2.7000000,1.8410860) (2.7066667,1.8429516) (2.7133333,1.8448210) (2.7200000,1.8466939) (2.7266667,1.8485705) (2.7333333,1.8504508) (2.7400000,1.8523346) (2.7466667,1.8542221) (2.7533333,1.8561132) (2.7600000,1.8580080) (2.7666667,1.8599063) (2.7733333,1.8618082) (2.7800000,1.8637137) (2.7866667,1.8656228) (2.7933333,1.8675355) (2.8000000,1.8694517) (2.8066667,1.8713715) (2.8133333,1.8732949) (2.8200000,1.8752219) (2.8266667,1.8771524) (2.8333333,1.8790864) (2.8400000,1.8810240) (2.8466667,1.8829651) (2.8533333,1.8849097) (2.8600000,1.8868578) (2.8666667,1.8888095) (2.8733333,1.8907647) (2.8800000,1.8927233) (2.8866667,1.8946855) (2.8933333,1.8966512) (2.9000000,1.8986203) (2.9066667,1.9005929) (2.9133333,1.9025690) (2.9200000,1.9045486) (2.9266667,1.9065316) (2.9333333,1.9085181) (2.9400000,1.9105080) (2.9466667,1.9125013) (2.9533333,1.9144981) (2.9600000,1.9164984) (2.9666667,1.9185020) (2.9733333,1.9205091) (2.9800000,1.9225196) (2.9866667,1.9245335) (2.9933333,1.9265507) (3.0000000,1.9285714) (3.0066667,1.9305955) (3.0133333,1.9326229) (3.0200000,1.9346538) (3.0266667,1.9366880) (3.0333333,1.9387255) (3.0400000,1.9407665) (3.0466667,1.9428107) (3.0533333,1.9448583) (3.0600000,1.9469093) (3.0666667,1.9489636) (3.0733333,1.9510212) (3.0800000,1.9530821) (3.0866667,1.9551464) (3.0933333,1.9572139) (3.1000000,1.9592848) (3.1066667,1.9613590) (3.1133333,1.9634364) (3.1200000,1.9655172) (3.1266667,1.9676012) (3.1333333,1.9696885) (3.1400000,1.9717790) (3.1466667,1.9738729) (3.1533333,1.9759699) (3.1600000,1.9780703) (3.1666667,1.9801738) (3.1733333,1.9822806) (3.1800000,1.9843907) (3.1866667,1.9865039) (3.1933333,1.9886204) (3.2000000,1.9907401) (3.2066667,1.9928630) (3.2133333,1.9949891) (3.2200000,1.9971184) (3.2266667,1.9992509) (3.2333333,2.0013865) (3.2400000,2.0035254) (3.2466667,2.0056674) (3.2533333,2.0078126) (3.2600000,2.0099609) (3.2666667,2.0121124) (3.2733333,2.0142671) (3.2800000,2.0164248) (3.2866667,2.0185857) (3.2933333,2.0207498) (3.3000000,2.0229169) (3.3066667,2.0250872) (3.3133333,2.0272606) (3.3200000,2.0294371) (3.3266667,2.0316167) (3.3333333,2.0337994) (3.3400000,2.0359852) (3.3466667,2.0381740) (3.3533333,2.0403660) (3.3600000,2.0425610) (3.3666667,2.0447590) (3.3733333,2.0469601) (3.3800000,2.0491643) (3.3866667,2.0513715) (3.3933333,2.0535818) (3.4000000,2.0557950) (3.4066667,2.0580114) (3.4133333,2.0602307) (3.4200000,2.0624530) (3.4266667,2.0646784) (3.4333333,2.0669067) (3.4400000,2.0691381) (3.4466667,2.0713724) (3.4533333,2.0736098) (3.4600000,2.0758501) (3.4666667,2.0780934) (3.4733333,2.0803396) (3.4800000,2.0825888) (3.4866667,2.0848410) (3.4933333,2.0870961) (3.5000000,2.0893542) (3.5066667,2.0916152) (3.5133333,2.0938791) (3.5200000,2.0961459) (3.5266667,2.0984157) (3.5333333,2.1006884) (3.5400000,2.1029640) (3.5466667,2.1052425) (3.5533333,2.1075239) (3.5600000,2.1098082) (3.5666667,2.1120953) (3.5733333,2.1143854) (3.5800000,2.1166783) (3.5866667,2.1189741) (3.5933333,2.1212727) (3.6000000,2.1235742) (3.6066667,2.1258785) (3.6133333,2.1281857) (3.6200000,2.1304958) (3.6266667,2.1328086) (3.6333333,2.1351243) (3.6400000,2.1374428) (3.6466667,2.1397641) (3.6533333,2.1420882) (3.6600000,2.1444151) (3.6666667,2.1467449) (3.6733333,2.1490774) (3.6800000,2.1514126) (3.6866667,2.1537507) (3.6933333,2.1560916) (3.7000000,2.1584352) (3.7066667,2.1607815) (3.7133333,2.1631306) (3.7200000,2.1654825) (3.7266667,2.1678371) (3.7333333,2.1701945) (3.7400000,2.1725546) (3.7466667,2.1749174) (3.7533333,2.1772829) (3.7600000,2.1796511) (3.7666667,2.1820221) (3.7733333,2.1843957) (3.7800000,2.1867721) (3.7866667,2.1891511) (3.7933333,2.1915328) (3.8000000,2.1939173) (3.8066667,2.1963043) (3.8133333,2.1986941) (3.8200000,2.2010865) (3.8266667,2.2034815) (3.8333333,2.2058793) (3.8400000,2.2082796) (3.8466667,2.2106826) (3.8533333,2.2130882) (3.8600000,2.2154965) (3.8666667,2.2179074) (3.8733333,2.2203209) (3.8800000,2.2227370) (3.8866667,2.2251557) (3.8933333,2.2275770) (3.9000000,2.2300009) (3.9066667,2.2324274) (3.9133333,2.2348564) (3.9200000,2.2372881) (3.9266667,2.2397223) (3.9333333,2.2421591) (3.9400000,2.2445984) (3.9466667,2.2470403) (3.9533333,2.2494847) (3.9600000,2.2519317) (3.9666667,2.2543812) (3.9733333,2.2568333) (3.9800000,2.2592879) (3.9866667,2.2617449) (3.9933333,2.2642046) (4.0000000,2.2666667)};
\draw[violet!75!black,densely dotted] (-2,1.6666667)--(2,1.6666667);
\fill[violet!75!black] (-2,1.6666667) circle(2.3pt);
\fill[violet!75!black] (2,1.6666667) circle(2.3pt);
\draw[blue,fill=white,thick] (0,1.44269504) circle(2.7pt);
\end{scope}
\end{scope}
\node[anchor=west,align=left,text width=4.35cm] at (8.9,4.4) {没有可变底数参数。\\[14pt]$f(-2)=f(2)=\frac53$。\\[14pt]左右同高，是偶函数；\\中间空点不能补入。};
\node[text width=12.8cm,align=center] at (7.1,.65) {$D=\mathbb R\setminus\{0\}$；括号内是奇函数，乘上奇函数 $x$ 后为偶函数。};
\end{tikzpicture}\end{center}


原式要求 $x\ne0$。令
\[
H(x)=\frac12+\frac1{2^x-1}=\frac{2^x+1}{2(2^x-1)}.
\]
由奇3的结论，$H(-x)=-H(x)$。因为 $x$ 也是奇函数，所以
\[
f(-x)=(-x)H(-x)=(-x)(-H(x))=xH(x)=f(x).
\]
故偶5是定义在 $\mathbb R\setminus\{0\}$ 上的偶函数，图中空点不能补入。
这一步使用的是 F017 的奇偶函数乘法规则；前三类证明已覆盖全部12式。
\end{proofbox}
